% How fast may the pushing policy move the fork? Tracking error, apparent mass and sim-to-real transfer.
% Written for the supervisors (robotics/Franka background, little RL background). Content-only section, included by
% main.tex (% How fast may the pushing policy move the fork? Tracking error, apparent mass and sim-to-real transfer.
% Written for the supervisors (robotics/Franka background, little RL background). Content-only section, included by
% main.tex (% How fast may the pushing policy move the fork? Tracking error, apparent mass and sim-to-real transfer.
% Written for the supervisors (robotics/Franka background, little RL background). Content-only section, included by
% main.tex (% How fast may the pushing policy move the fork? Tracking error, apparent mass and sim-to-real transfer.
% Written for the supervisors (robotics/Franka background, little RL background). Content-only section, included by
% main.tex (\input{open_questions/action_limits_problem}) or compiled alone via action_limits_report.tex.
% Uses amsmath, booktabs and array (added to main.tex for this section). Numbers: scripts/sweep_action_poses.py
% (default and --full_speed), scripts/check_push_env.py and docs/notes/2026-10-07.md, 2026-10-08.md in the project repo.

% ======================================= BEGIN CONTENT =======================================
\section{Open problem: how fast may the policy move the fork?}
\label{sec:al}

\paragraph{Status (7--8 October 2026).} Problem identified in simulation; values for the first training runs set
on 7 October (option B1) and revised on 8 October after longer test moves (Section~\ref{sec:al-plan}); the
long-term solution is open. We ask for feedback on the analysis and on the
questions at the end (Section~\ref{sec:al-questions}).

\paragraph{In one paragraph.} The policy commands the fork through a Cartesian impedance controller, the same in
simulation and on the FR3. What blocks us is a trade-off between \emph{agility} and a \emph{small tracking error}.
A small error transfers better~\cite{al-aljalbout2024}, because then the fork's motion is set by our controller
(identical on both sides) rather than by the arm's dynamics (not identical, not yet identified). But an impedance
controller is a spring of stiffness $K$ between the commanded and the actual tool position, and it can only
accelerate the arm by being stretched: while the commanded motion accelerates at $a$, the tool lags behind by about
$\Lambda a/K$, where $\Lambda$ is the arm's apparent mass at the tool tip (8--22\,kg depending on pose and direction,
half of it from the motors' reflected inertia). With a small lag (2\,mm) and a stiffness we consider acceptable near a
plant (400--1000\,N/m), the fork may accelerate only at about 0.04--0.09\,m/s$^2$; we use 0.08\,m/s$^2$ at
1000\,N/m. This section explains the
trade-off, what we measured, and every design choice that moves it.

% ---------------------------------------------------------------------------------------------------------------------
\subsection{Background}
\label{sec:al-background}

\paragraph{What the policy does.} The policy is a neural network trained by reinforcement learning (RL) in
simulation. At 31.25\,Hz (every 32\,ms, one ``policy step'') it reads the state (fork pose, stem shape, target) and
outputs a small displacement and rotation of the fork's \emph{target} pose $x_d$. The next target is the previous
target plus this step, so commanded steps are not lost even if the tool lags behind. Within the 32\,ms the target
moves at constant velocity $\dot x_d$ from the old to the new value. A Cartesian impedance controller at 1\,kHz
(2\,ms in simulation) turns the target into joint torques:
\begin{equation}
  \tau = J^\top\!\left[\,K\,(x_d - x) + D\,(\dot x_d - \dot x)\,\right] + \tau_\text{null} + \tau_\text{Coriolis}.
  \label{eq:al-law}
\end{equation}
$x$ is the measured tool pose, $J$ the tool-tip Jacobian, $K$ the stiffness and $D$ the damping (written here for
the position; the orientation has its own stiffness $K_o$). This is the law of Franka's example controllers, with
two changes: the target velocity $\dot x_d$ is fed forward (Franka's examples use $\dot x_d = 0$), and the damping
$D = 2\sqrt{K}\,\Lambda^{1/2}$ is chosen critical for the arm's actual inertia $\Lambda$ (Section~\ref{sec:al-problem})
instead of Franka's $2\sqrt K$, which is critical only for a 1\,kg tool.

\paragraph{Null-space spring $\tau_\text{null}$.} The FR3 has 7 joints, but the tool pose has only 6 degrees of
freedom. For every tool pose there is therefore a one-parameter family of arm postures (the elbow can swing about the
line from shoulder to wrist without moving the tool): the \emph{null space}. The first term of
Eq.~\eqref{eq:al-law} does not act in this direction, so the elbow would drift freely. A weak joint-space spring
(0.5\,N\,m/rad, as in Franka's examples) pulls the posture towards a rest posture; it is projected so that it does
not push the tool. It plays no role in the problem below (we checked: switching it off does not change the lag).

\paragraph{What RL learns, and why the action space matters for the transfer.} During training the policy tries
millions of commands in simulation and keeps what earns reward. It learns \emph{how the simulated fork reacts to its
commands}, including any lag, overshoot or coupling, and relies on it. If the real robot reacts differently, the
learned behaviour is wrong by that difference: this is the sim-to-real (``reality'') gap. The policy cannot correct
what it never saw in simulation. The usual tools against the gap are:
\begin{center}
{\renewcommand{\arraystretch}{1.15}
\begin{tabular}{>{\raggedright\arraybackslash}p{0.24\textwidth}>{\raggedright\arraybackslash}p{0.42\textwidth}>{\raggedright\arraybackslash}p{0.24\textwidth}}
  \toprule
  Tool & What it does & Status \\
  \midrule
  A. Action-space design & choose commands whose result depends little on what differs between sim and real
    (same controller, small tracking error) & controller matched; this section \\
  \addlinespace
  B. Identification & measure the real robot and correct the simulation & planned, before the first real run \\
  \addlinespace
  C. Domain randomization & vary unknown parameters in training, so the policy works over a range & planned \\
  \addlinespace
  D. Model the robot's safety layer & limits and reflexes of the FR3 become terminations/penalties in sim, so the
    policy does not learn motions the robot refuses & planned \\
  \addlinespace
  E. Adaptation & infer differences online or fine-tune on the robot & not planned \\
  \bottomrule
\end{tabular}}
\end{center}
A large study of action spaces on a Franka arm~\cite{al-aljalbout2024} (13 action spaces, about 250 policies,
reaching and pushing, transferred to the real robot) concludes that \emph{``an action space which yields or
naturally limits the tracking error transfers better''}; torque actions, whose result is set by the robot's
dynamics alone, could not be transferred safely. The same study warns that a limiter which acts on the robot but is
not visible to the policy in simulation (they tried a rate limiter) made learning collapse. A caveat: their controllers were joint
impedance controllers (Cartesian actions were converted to joint targets by inverse kinematics), with gains that
the paper does not report, and the tasks were reaching and pushing a box. The conclusion carries over in principle,
since in every impedance-type controller a larger tracking error means more dependence on the arm's dynamics; the
numbers do not, and the paper does not test a soft Cartesian impedance controller like ours.

\paragraph{Where the gap enters our system.} The chain from a policy output to the stem, step by step:
\begin{center}
{\renewcommand{\arraystretch}{1.15}
\begin{tabular}{>{\raggedright\arraybackslash}p{0.32\textwidth}>{\raggedright\arraybackslash}p{0.20\textwidth}>{\raggedright\arraybackslash}p{0.38\textwidth}}
  \toprule
  Step & Same in sim and real? & Source of differences \\
  \midrule
  policy step $\to$ target pose & yes (same code) & none \\
  target $\to$ torques, law~\eqref{eq:al-law} & yes (ported to C++, checked) & port errors; $J$, $M$ from different
    models \\
  \textbf{torques $\to$ fork motion} & \textbf{no} & arm inertia, joint friction, motor inertia, delays, 1 vs.\ 2\,ms
    control period \\
  fork $\leftrightarrow$ stem contact & no & friction, stem stiffness and damping (randomized) \\
  stem $\to$ observation & no & camera noise, latency \\
  FR3 safety layer & \textbf{missing in sim} & reflexes stop the robot \\
  \bottomrule
\end{tabular}}
\end{center}
The step ``torques $\to$ fork motion'' is where the arm's dynamics act. How much of it reaches the fork's motion
depends on the tracking error: if the tool follows the target within a small error, its motion is essentially the
target's motion, which our code computes identically in sim and real; the larger the error, the more the motion is
shaped by the arm's dynamics, which differ.

% ---------------------------------------------------------------------------------------------------------------------
\subsection{The problem}
\label{sec:al-problem}

\paragraph{What blocks us: agility versus a small tracking error.} We want two things from the fork's motion:
\begin{enumerate}
  \item \textbf{Small lag behind the commanded target}, so that the fork's motion is set by our controller, which
        is the same in simulation and on the FR3, and not by the arm's dynamics, which differ and are not yet
        identified. Action spaces that keep the tracking error small transfer better~\cite{al-aljalbout2024}.
  \item \textbf{Agility}: reach the stem well within the 10\,s episode, brake late before contact, and react quickly
        to how the stem moves.
\end{enumerate}
With an impedance controller the two conflict. The lag grows with how quickly the commanded motion accelerates, so
a small lag allows only gentle accelerations. In numbers: with $K = 1000$\,N/m and at most 2\,mm of lag, the fork may
accelerate at only about 0.09\,m/s$^2$ where the arm is heaviest; at the 0.08\,m/s$^2$ we use, it needs 1.25\,s to
reach 10\,cm/s and 6\,cm to stop from that speed (2.5\,s and 25\,cm for 20\,cm/s). Either
we accept this sluggishness, or we accept more lag (and with it a larger dependence on the real arm), or we change
the controller. This is the decision we need help with. The rest of this section explains where the conflict comes
from, how severe it is for the task, and how severe it will be on the real robot.

\paragraph{Why the tool lags: the impedance law as a mass on a spring.} Consider one direction, free space,
gravity compensated. The controller makes the tool tip behave like a mass $\Lambda$ (the arm's inertia felt at the
tool tip, explained in the next paragraph) attached by a spring $K$ and a damper $D$ to the moving target:
\begin{equation}
  \Lambda\,\ddot x = K\,(x_d - x) + D\,(\dot x_d - \dot x).
\end{equation}
Write the lag as $e = x_d - x$ (target minus tool). Then $\ddot x = \ddot x_d - \ddot e$, and substituting gives
\begin{equation}
  \Lambda\,\ddot e + D\,\dot e + K\,e = \Lambda\,\ddot x_d.
  \label{eq:al-error}
\end{equation}
The lag itself behaves like a damped mass-spring system, driven by the force $\Lambda \ddot x_d$, i.e.\ by the
\emph{acceleration of the target}. Three cases follow:
\begin{enumerate}
  \item \textbf{Target at rest or moving at constant speed} ($\ddot x_d = 0$): the right-hand side is zero, so any lag
        dies out and $e \to 0$. Physically: moving at constant speed in free space needs no force (no friction, no
        gravity), so the spring does not need to be stretched. This holds only because the target velocity is fed
        forward. Without it the damper pulls back with $D\dot x$ all the time, and the spring must balance it:
        $e = D\dot x_d / K$. With $D = 2\sqrt{K\Lambda} = 180$\,N\,s/m ($K = 400$\,N/m, $\Lambda = 20$\,kg) and
        20\,cm/s this gives 90\,mm; we measured 82\,mm (Table~\ref{tab:al-time}).
  \item \textbf{Target accelerating at constant $a$}: the right-hand side is the constant force $\Lambda a$. Once the
        transient has died out ($\ddot e = \dot e = 0$), $K e = \Lambda a$, so
        \begin{equation}
          e = \frac{\Lambda\, a}{K}.
          \label{eq:al-lag}
        \end{equation}
        Physically: Newton's law. To accelerate the mass $\Lambda$ at $a$, the force $\Lambda a$ is needed, and its only
        source is the stretched spring. Example: $\Lambda = 20$\,kg, $a = 0.1$\,m/s$^2$, $K = 1000$\,N/m: $e = 2$\,mm.
        The same holds when braking (the tool then runs ahead of the target, which shows up as overshoot).
  \item \textbf{How fast the lag builds up and dies out}: Eq.~\eqref{eq:al-error} is critically damped, with the time
        constant $\sqrt{\Lambda/K}$: 0.14\,s for 20\,kg and 1000\,N/m, 0.22\,s for 400\,N/m, i.e.\ 4--7 policy steps.
        The lag reaches the value of Eq.~\eqref{eq:al-lag} if the acceleration lasts longer than about three time
        constants, and it disappears again within a similar time once the target moves at constant speed. This
        matters for testing: short test moves end before the lag has settled. Our first sweep (20 policy steps per
        segment, reversals that only braked) therefore showed 1.9\,mm at 0.1\,m/s$^2$; moves that brake \emph{and}
        reverse, or that run up to the speed cap, accelerate long enough and reached 2.15\,mm
        (Table~\ref{tab:al-final}).
\end{enumerate}

\paragraph{The mass in this equation: the apparent mass $\Lambda$ (``8--20\,kg at the tool tip'').} The trade-off
is severe because the mass that the spring has to accelerate is large. If one pushes the tool tip of
the arm (joints free, gravity compensated) with a force $F$, the tip starts to accelerate with $\ddot x = \Lambda^{-1}
F$. $\Lambda$ is the arm's inertia as seen at the tool tip, the \emph{apparent mass} (or operational-space inertia):
\begin{equation}
  \Lambda = \left(J M^{-1} J^\top\right)^{-1},
  \label{eq:al-lambda}
\end{equation}
with $M$ the $7\times7$ joint-space mass matrix and $J$ the $6\times7$ tool-tip Jacobian. $\Lambda$ is a $6\times6$
matrix (force and moment versus linear and angular acceleration) and depends on the pose. Its translational
diagonal entries say how heavy the arm feels when the tip is pushed along $x$, $y$ or $z$. At the earlier start pose
(2026-10-07; tool tip at $(0.30, 0, 0.55)$\,m in front of the robot, fork horizontal), computed from the simulated FR3:
\begin{center}
{\renewcommand{\arraystretch}{1.15}
\begin{tabular}{lrrr}
  \toprule
  Apparent mass, start pose & along $x$ (forward) & along $y$ (sideways) & along $z$ (up) \\
  \midrule
  links only & 11.9\,kg & 4.2\,kg & 5.1\,kg \\
  links + motors' reflected inertia (our simulation) & 20.5\,kg & 7.6\,kg & 10.6\,kg \\
  \bottomrule
\end{tabular}}
\end{center}
Pushing the tip forward with 20\,N accelerates it at about 1\,m/s$^2$, as if it were a 20\,kg mass; the same push
sideways gives 2.6\,m/s$^2$. Forward motion needs the heavy shoulder and elbow joints, sideways motion mostly the
base rotation and the wrist. The second row includes the motors' reflected inertia $n^2 I_\text{motor}$ ($n$: gear
ratio), which we add to $M$ from Franka's description files: 0.61, 0.61, 0.46, 0.46, 0.21, 0.21, 0.21\,kg\,m$^2$ for
joints 1--7 (at the wrist about 100 times the link's own inertia). Without these terms the simulated arm oscillated
and pressed the stem 5 times harder. The lag measurements below give the same value independently: from
Eq.~\eqref{eq:al-lag}, $\Lambda = K e / a \approx 16$--20\,kg along $x$ at the start pose, and up to 21.5\,kg
at the far, high start points of the sweep ($x$ 0.60, $z$ 0.45\,m).

\paragraph{How severe the trade-off is for the task.} Turning Eq.~\eqref{eq:al-lag} around: for a largest acceptable lag $e_\text{max}$, the
target may accelerate at most at
\begin{equation}
  a_\text{max} \approx \frac{K\, e_\text{max}}{\Lambda}.
  \label{eq:al-amax}
\end{equation}
With $\Lambda = 21.5$\,kg (the heaviest pose and direction) and $e_\text{max} = 2$\,mm: 0.04\,m/s$^2$ at
$K = 400$\,N/m and 0.09\,m/s$^2$ at 1000\,N/m. We use 0.08\,m/s$^2$, which leaves about 15\,\% margin (lag
$\le 1.72$\,mm; measured 1.68\,mm, Table~\ref{tab:al-final}). What this means for the task (the
target speeds up at $a_\text{max}$ until the speed cap, then brakes at $a_\text{max}$):
\begin{center}
{\renewcommand{\arraystretch}{1.15}
\begin{tabular}{rrrrr}
  \toprule
  largest acceleration & time from rest & braking distance & 31\,cm, rest to rest & 31\,cm, rest to rest \\
  $a_\text{max}$ & to 20\,cm/s & from 20\,cm/s & (20\,cm/s cap) & (10\,cm/s cap) \\
  \midrule
  0.05\,m/s$^2$ & 4.0\,s & 40\,cm & 5.0\,s & 5.1\,s \\
  \textbf{0.08\,m/s}$^\mathbf{2}$ & 2.5\,s & 25\,cm & 3.9\,s & \textbf{4.4\,s} \\
  0.10\,m/s$^2$ & 2.0\,s & 20\,cm & 3.5\,s & 4.1\,s \\
  0.25\,m/s$^2$ & 0.8\,s & 8\,cm  & 2.3\,s & 3.5\,s \\
  0.50\,m/s$^2$ & 0.4\,s & 4\,cm  & 1.9\,s & 3.3\,s \\
  \bottomrule
\end{tabular}}
\end{center}
31\,cm is the distance from the start pose to the farthest stem tip. The episode lasts 10\,s, and the stem must
still be pushed after the approach. The cost of a small $a_\text{max}$ is therefore not the top speed but agility:
the policy must start braking long before it touches the stem and cannot react quickly to how the stem moves.

\paragraph{Why not simply make the episode longer?} A longer episode gives a slow policy time to finish, but it is
not free. RL learns from simulated experience, so the training time grows roughly in proportion to the simulated
time per episode: a 20\,s episode costs about twice the training time of a 10\,s episode for the same number of
attempts. In addition, the reward for reaching the target arrives later after the actions that caused it, which
makes it harder for the policy to learn which actions were good. We therefore keep 10\,s as long as the task fits,
and lengthen it only if the trained policy demonstrably runs out of time.

\paragraph{Rotations: the same effect, plus a coupling.} Turning the fork works in the same way: the orientation
spring $K_o$ must twist to give the rotational inertia an angular acceleration $\alpha$, so the angle lags by about
$\Lambda_o \alpha / K_o$ ($\Lambda_o$: the rotational part of $\Lambda$). In addition, $\Lambda$ is a full $6\times6$
matrix: its off-diagonal blocks couple translation and rotation. The reason is geometric: the tool tip is 10--20\,cm
away from the wrist's joint axes, and the joints that turn the fork also move the tip, while those that move the
tip also turn the fork. The control law, however, treats position and orientation as two independent springs.
Two consequences, both measured:
\begin{itemize}
  \item \textbf{Tilt while translating.} Accelerating the tip at $a$ needs, besides the force, a moment
        $\Lambda_{op}\, a$ ($\Lambda_{op}$: off-diagonal block). It comes from the orientation spring, which must twist:
        the fork tilts by about $\Lambda_{op} a / K_o$. Table~\ref{tab:al-diag} shows the tilt falling with $1/K_o$
        (2.0$^\circ$ at $K_o = 30$, 0.26$^\circ$ at 300\,N\,m/rad).
  \item \textbf{Tip drift while turning.} Turning the fork at angular acceleration $\alpha$ needs, besides the moment, a
        force at the tip, which the position spring supplies by stretching: the tip moves away from its target by
        about $\Lambda_{po}\,\alpha / K$, although the position command is constant (Tables~\ref{tab:al-rot}
        and~\ref{tab:al-b1}).
\end{itemize}
Both effects scale with the acceleration, like the lag. A controller with inertia feedforward (operational-space
control, option B3 below) would supply these coupling forces and moments from the model and remove both.

\paragraph{Contact.} When the tool is blocked (e.g.\ pressed against the stem), the target keeps moving and the
lag becomes force: $K e$. Small lags therefore also mean small unintended forces. A hard bound is added by
clamping the target: at every control step, $\lVert x_d - x \rVert \le 4$\,mm (spring force $\le 4$\,N at
$K = 1000$\,N/m) and the angle between target and tool orientation $\le 3^\circ$ (moment $\le 5.2$\,N\,m at
$K_o = 100$\,N\,m/rad). In free motion the clamp is inactive, because the lag stays below 2\,mm. While it holds the
target, the target moves with the tool, so the target velocity fed forward is the tool's own velocity: the damper
then adds no force to the bounded spring force, and the controller does not switch back and forth between holding and
releasing. Not bounded by the clamp: the momentum of the arm when the moving fork first hits something (as with any
impedance controller); the policy is meant to learn to approach gently (contact-force and smoothness penalties).

\paragraph{How severe the trade-off will be on the real robot: which $\Lambda$ does the FR3 have?} All numbers above
use the simulated $\Lambda$. The real arm's $\Lambda$ is not known yet, and two observations from Franka's software
suggest that it may differ:
\begin{itemize}
  \item libfranka (current version, September 2026) computes its model (\texttt{franka::Model::mass}) with the
        Pinocchio library from the URDF that the robot provides. The FR3 URDF shipped with libfranka's tests has
        the same link inertias as Franka's description files but \emph{no} motor-inertia entries, and libfranka never
        sets Pinocchio's rotor-inertia (``armature'') term. So libfranka's mass matrix most likely \emph{excludes}
        the motors' reflected inertia. (Read in the source code, not yet checked on our robot.)
  \item The FR3 executes commanded torques with an internal joint torque controller that uses the link-side
        torque sensors. On the DLR lightweight arms, from which Franka's design derives, such torque feedback
        reduces the apparent motor inertia. If the FR3 does the same, the real arm feels lighter than our
        simulation (between the two rows of the table above), and libfranka's model would be the right one.
\end{itemize}
How much this matters depends on the option chosen (Section~\ref{sec:al-choices}):
\begin{center}
{\renewcommand{\arraystretch}{1.15}
\begin{tabular}{>{\raggedright\arraybackslash}p{0.24\textwidth}>{\raggedright\arraybackslash}p{0.66\textwidth}}
  \toprule
  Option & Role of the real $\Lambda$ \\
  \midrule
  B1 (small lag, now) & \textbf{Little}, and this is exactly why small lags transfer well: the unknown $\Lambda$
    acts on the motion only through the lag. The acceleration limit was sized for $\le 2$\,mm with the simulated
    $\le 21.5$\,kg; a real arm of 12\,kg would lag 1\,mm, a heavier one (30\,kg) 2.4\,mm. The policy is then wrong by about a
    millimetre. \textbf{One exception: the damping} $D = 2\sqrt K \Lambda^{1/2}$ is computed from a model of $\Lambda$.
    Before the first real run we must decide which mass matrix the controller on the robot uses (libfranka's, or
    libfranka's plus the motor terms); with the wrong one the tool is under- or over-damped. \\
  \addlinespace
  B2 (larger lag) & \textbf{More}: with 5\,mm of lag set by the arm, a 40\,\% error in $\Lambda$ shifts the motion by
    2\,mm. \\
  \addlinespace
  B3 (inertia feedforward, later) & \textbf{Central}: the controller supplies the accelerating force from the model
    of $\Lambda$, so the remaining lag is set by the model error. Requires the identification first. \\
  \bottomrule
\end{tabular}}
\end{center}

% ---------------------------------------------------------------------------------------------------------------------
\subsection{Measurements}
\label{sec:al-measurements}

\paragraph{Setup.} Simulation (Isaac Lab, PhysX) of the FR3 with the fork, gravity compensated, no joint friction,
control period 2\,ms, policy step 32\,ms, no stem (free space). Every test commands the fork with the same commands
a policy could send, through the full controller.

\paragraph{Commands and limits.} A command is a target step per policy step. Three limits act on it:
\begin{itemize}
  \item \emph{speed cap}: the largest step, e.g.\ 3.2\,mm per 32\,ms = 10\,cm/s (rotation: 1.44$^\circ$ = 45$^\circ$/s);
  \item \emph{acceleration limit}: the largest change of the step from one policy step to the next. A change of
        0.1\,mm per step means $a = 0.1\,\text{mm}/(32\,\text{ms})^2 \approx 0.1$\,m/s$^2$, 0.08\,mm per step
        0.08\,m/s$^2$ (rotation: 0.02$^\circ$ per step $\approx$ 20$^\circ$/s$^2$). Commanding ``full speed'' from rest therefore makes the target speed up at
        this acceleration until the speed cap; commanding ``stop'' makes it brake at the same rate;
  \item the first design (Table~\ref{tab:al-first}) had a speed cap only, so the target jumped from rest to full
        speed in one step.
\end{itemize}

\paragraph{Test moves.} \emph{Straight move}: full command along one axis of the robot base frame ($\pm x$, $\pm y$,
$\pm z$) for a number of policy steps, then zero command (``hold'') for 45 steps (1.4\,s), during which the target
brakes and stops. \emph{Reversal}: full command along $+x$, then immediately along $-x$ (same for $y$, $z$), then hold;
this contains the largest change of velocity. Since 8 October the second segment is twice as long, so that the
target brakes \emph{and} reverses (before, it only braked to rest). \emph{Full-speed move} (since 8 October): from
rest up to the speed cap, 10 steps at the cap, then braking to rest (about 16\,cm; rotations about 120$^\circ$),
one per axis, in a direction with room for it (towards the middle of the workspace; up along $z$).
\emph{Rotation moves}: the same about the base axes, with zero position command. \emph{Start points}: either the
start pose only, or 9 points spanning the workspace in which the stems stand (tool tip at $x$ = 0.35 and 0.60\,m,
$y = \pm0.20$\,m, $z$ = 0.20 and 0.45\,m, plus the start pose). Moves that come within 0.25\,rad of a joint limit
are reported but not counted (decision 8 October): there the arm, not the controller, limits the motion. This
concerns the two start points low and close to the robot ($x$ 0.35, $z$ 0.20\,m), where joint 4 is within
0.21\,rad of its limit.

\paragraph{What is measured (all tables).}
\begin{itemize}
  \item \emph{Lag} (following error): the largest distance between the target position and the actual tool-tip
        position during the move and the hold. The quantity of Eq.~\eqref{eq:al-lag}.
  \item \emph{Overshoot}: how far the tool tip goes past the target's final position after it stops.
  \item \emph{Tilt}: the largest change of the fork's orientation during a straight move, although no rotation is
        commanded.
  \item \emph{Angle lag}: during rotation moves, the largest angle between the target and the actual orientation.
  \item \emph{Tip drift}: during rotation moves, the largest displacement of the tool tip, although no translation is
        commanded.
\end{itemize}
Each table reports the worst value over all moves (and start points). The criterion used so far: lag $\le 2$\,mm,
overshoot $\le 5$\,mm, tilt $\le 1^\circ$; angle lag $\le 2^\circ$, tip drift $\le 2$\,mm.

\begin{table}[h]
  \centering
  \caption{\textbf{First design}: speed cap only, no acceleration limit, $K = 400$\,N/m, $K_o = 30$\,N\,m/rad. Straight
  moves of 30 policy steps at full command from 9 start points, so each move starts with a jump from rest to full speed.
  Lag measured over the last quarter of each move (as defined at the time), overshoot after it. Worst case over all
  moves and start points. This design kept the lag below 2\,mm by limiting the speed to 5.5\,cm/s.}
  \label{tab:al-first}
  {\renewcommand{\arraystretch}{1.15}
  \begin{tabular}{rrrr}
    \toprule
    speed cap (step per 32\,ms) & speed & lag & overshoot \\
    \midrule
    2.0\,mm  & 6.3\,cm/s & 2.25\,mm & 3.9\,mm \\
    1.75\,mm & 5.5\,cm/s & 1.94\,mm & 3.4\,mm \\
    1.5\,mm  & 4.7\,cm/s & 1.63\,mm & 2.9\,mm \\
    \bottomrule
  \end{tabular}}

  \smallskip
  {\small Reading: with a cap of 5.5\,cm/s and 7.8$^\circ$/s the approach to a far stem took more than half of the
  episode, and a 90$^\circ$ turn did not fit in one. We therefore raised the caps to 20\,cm/s and 45$^\circ$/s and
  added the acceleration limit; Tables~\ref{tab:al-time}--\ref{tab:al-b1} show what that does.}
\end{table}

\begin{table}[h]
  \centering
  \caption{\textbf{How the lag evolves over time} (illustrates Eq.~\eqref{eq:al-error}). Start pose, $K = 400$\,N/m,
  $K_o = 30$\,N\,m/rad. One straight move along $+x$ at full command for 40 policy steps. With the acceleration limit
  of 0.49\,m/s$^2$ the target speeds up during the first 13 steps (0.42\,s) and then moves at constant speed.
  Lag of the tool tip behind the target after 10, 20, 30 and 40 policy steps.}
  \label{tab:al-time}
  {\renewcommand{\arraystretch}{1.15}
  \begin{tabular}{lrrrr}
    \toprule
    & \multicolumn{4}{c}{lag after $n$ policy steps} \\
    \cmidrule(lr){2-5}
    & 10 (0.32\,s) & 20 (0.64\,s) & 30 (0.96\,s) & 40 (1.28\,s) \\
    & accelerating & constant speed & constant speed & constant speed \\
    \midrule
    20\,cm/s, velocity feedforward on  & 9.4\,mm & 10.4\,mm & 4.7\,mm & 1.9\,mm \\
    10\,cm/s, velocity feedforward on  & 7.7\,mm & 4.7\,mm  & 2.0\,mm & 0.7\,mm \\
    20\,cm/s, velocity feedforward off & 23\,mm  & 59\,mm   & 75\,mm  & 82\,mm \\
    \bottomrule
  \end{tabular}}

  \smallskip
  {\small With feedforward the lag builds up while the target accelerates and dies out at constant speed (case 1 and 2
  of Section~\ref{sec:al-problem}); without it, the lag grows towards $D\dot x_d/K$. Switching off the null-space
  spring gave the same numbers.}
\end{table}

\begin{table}[h]
  \centering
  \caption{\textbf{Effect of the stiffnesses $K$ and $K_o$.} Start pose, speed cap 20\,cm/s. A straight move along $+x$
  at full command for 25 policy steps (0.8\,s, enough to reach about the speed cap), hold for 45 steps, then the same
  along $-x$ and hold. Two acceleration limits (columns). Largest lag of the tool tip and largest tilt of the fork
  over the whole test.}
  \label{tab:al-diag}
  {\renewcommand{\arraystretch}{1.15}
  \begin{tabular}{rrrrrr}
    \toprule
    & & \multicolumn{2}{c}{acceleration 0.24\,m/s$^2$} & \multicolumn{2}{c}{acceleration 0.49\,m/s$^2$} \\
    \cmidrule(lr){3-4} \cmidrule(lr){5-6}
    $K$ [N/m] & $K_o$ [N\,m/rad] & lag & tilt & lag & tilt \\
    \midrule
    400  & 30  & 9.7\,mm  & 1.80$^\circ$ & 13.6\,mm & 2.40$^\circ$ \\
    400  & 100 & 10.5\,mm & 0.66$^\circ$ & 14.4\,mm & 0.84$^\circ$ \\
    400  & 300 & 10.6\,mm & 0.23$^\circ$ & 14.8\,mm & 0.29$^\circ$ \\
    1000 & 30  & 4.0\,mm  & 2.03$^\circ$ & 6.4\,mm  & 2.69$^\circ$ \\
    1000 & 100 & 4.6\,mm  & 0.73$^\circ$ & 7.1\,mm  & 1.00$^\circ$ \\
    1000 & 300 & 4.8\,mm  & 0.26$^\circ$ & 7.5\,mm  & 0.37$^\circ$ \\
    \bottomrule
  \end{tabular}}

  \smallskip
  {\small Reading: the lag depends on $K$ (2.5 times stiffer, about 2.4 times less lag) and on the acceleration, hardly
  on $K_o$; the tilt depends on $K_o$ (10 times stiffer, about 8 times less tilt). Every entry fails the 2\,mm criterion.}
\end{table}

\begin{table}[h]
  \centering
  \caption{\textbf{Workspace sweep at $K = 1000$\,N/m, $K_o = 30$\,N\,m/rad}, caps 20\,cm/s and 45$^\circ$/s. At each
  start point: straight moves along $\pm x, \pm y, \pm z$ and reversals in $x, y, z$ (30 policy steps per segment, then
  hold), then the same as rotations about the base axes. Each row is one pair of acceleration limits (translation,
  rotation). Worst case over all moves and over the five start points closest to the robot: from the start points at
  $x = 0.60$\,m these moves of up to 20\,cm ran to the edge of the arm's reach, which is a flaw of the test, not of the
  controller.}
  \label{tab:al-rot}
  {\renewcommand{\arraystretch}{1.15}
  \begin{tabular}{rrrrrr}
    \toprule
    \multicolumn{3}{c}{translation moves} & \multicolumn{3}{c}{rotation moves} \\
    \cmidrule(lr){1-3} \cmidrule(lr){4-6}
    acceleration limit & lag & tilt & angular acceleration limit & angle lag & tip drift \\
    \midrule
    0.24\,m/s$^2$ & 5.0\,mm  & 3.0$^\circ$ & 20$^\circ$/s$^2$  & 9.9$^\circ$$^*$ & 14.0\,mm$^*$ \\
    0.49\,m/s$^2$ & 8.5\,mm  & 5.0$^\circ$ & 49$^\circ$/s$^2$  & 4.7$^\circ$ & 4.0\,mm \\
    0.73\,m/s$^2$ & 11.0\,mm & 5.7$^\circ$ & 98$^\circ$/s$^2$  & 6.8$^\circ$ & 7.0\,mm \\
    0.98\,m/s$^2$ & 12.6\,mm & 6.1$^\circ$ & 195$^\circ$/s$^2$ & 9.4$^\circ$ & 14.3\,mm \\
    \bottomrule
  \end{tabular}}

  \smallskip
  {\small Reading: all rows fail the criterion; the lag grows with the acceleration as expected, the tilt is large
  because $K_o$ is soft. $^*$ one start point ($x$ 0.35, $y$ 0.20, $z$ 0.45\,m), reversal about $x$; at the other four
  points at most 4.1$^\circ$ / 3.7\,mm. Not yet understood (likely a wrist joint near its limit).}
\end{table}

\begin{table}[h]
  \centering
  \caption{\textbf{Workspace sweep for the interim choice of 7 October: $K = 1000$\,N/m, $K_o = 100$\,N\,m/rad}, caps
  20\,cm/s and 45$^\circ$/s, smaller acceleration limits. Same moves as Table~\ref{tab:al-rot} but 20 policy steps per
  segment (moves of at most about 6\,cm, so all 9 start points are valid). Worst case over all moves and all 9 start
  points. \emph{Superseded by Table~\ref{tab:al-final}}: these moves reached at most 2\,mm per step (6.25\,cm/s), not
  the caps, and the reversals only braked, so the lag had not settled.}
  \label{tab:al-b1}
  {\renewcommand{\arraystretch}{1.15}
  \begin{tabular}{rrrrrrr}
    \toprule
    \multicolumn{4}{c}{translation moves} & \multicolumn{3}{c}{rotation moves} \\
    \cmidrule(lr){1-4} \cmidrule(lr){5-7}
    acceleration limit & lag & overshoot & tilt & angular accel.\ limit & angle lag & tip drift \\
    \midrule
    0.05\,m/s$^2$ & 1.3\,mm & 0.9\,mm & 0.17$^\circ$ & 10$^\circ$/s$^2$ & 0.30$^\circ$ & 0.55\,mm \\
    \textbf{0.10\,m/s}$^\mathbf{2}$ & \textbf{1.9\,mm} (2.4$^*$) & \textbf{1.8\,mm} & \textbf{0.33}$^\circ$ &
      \textbf{20}$^\circ$\textbf{/s}$^\mathbf{2}$ & \textbf{0.51}$^\circ$ & \textbf{1.35\,mm} \\
    0.15\,m/s$^2$ & 2.9\,mm (3.5$^*$) & 2.6\,mm & 0.47$^\circ$ & 39$^\circ$/s$^2$ & 1.31$^\circ$ & 4.41\,mm \\
    \bottomrule
  \end{tabular}}

  \smallskip
  {\small Reading: the criterion (lag $\le 2$\,mm, tilt $\le 1^\circ$, angle lag $\le 2^\circ$, tip drift $\le 2$\,mm)
  is met at 0.1\,m/s$^2$ and 20$^\circ$/s$^2$ (bold, chosen) and fails at 0.15\,m/s$^2$ / 39$^\circ$/s$^2$. The lag
  of 1.9\,mm matches Eq.~\eqref{eq:al-lag} with $\Lambda \approx 20$\,kg. $^*$ at the two start points low and close
  to the robot ($x$ 0.35, $z$ 0.20\,m), where joint 4 is within 0.21\,rad of its limit; the task will keep the tool
  away from there.}
\end{table}

\begin{table}[h]
  \centering
  \caption{\textbf{Revised sweep (8 October): real reversals and full-speed moves.} $K = 1000$\,N/m,
  $K_o = 100$\,N\,m/rad, rotation limits unchanged (45$^\circ$/s, 20$^\circ$/s$^2$). \emph{Short moves}: as in
  Table~\ref{tab:al-b1}, but reversals that brake and reverse. \emph{Full speed}: one move per axis up to the cap and
  back to rest. Worst case over the counted moves of the 9 start points (moves within 0.25\,rad of a joint limit are
  not counted).}
  \label{tab:al-final}
  {\renewcommand{\arraystretch}{1.15}
  \begin{tabular}{rrrrrrrr}
    \toprule
    & & \multicolumn{3}{c}{short moves, real reversals} & \multicolumn{3}{c}{full speed} \\
    \cmidrule(lr){3-5} \cmidrule(lr){6-8}
    accel.\ limit & speed cap & lag & overshoot & tilt & lag & overshoot & moves counted \\
    \midrule
    0.10\,m/s$^2$ & 20\,cm/s & 2.15\,mm & 1.65\,mm & 0.35$^\circ$ & 3.35\,mm$^*$ & 3.01\,mm & 7 of 27 \\
    0.10\,m/s$^2$ & 10\,cm/s & 2.15\,mm & 1.65\,mm & 0.35$^\circ$ & 2.00\,mm & 1.93\,mm & 21 of 27 \\
    \textbf{0.08\,m/s}$^\mathbf{2}$ & \textbf{10\,cm/s} & \textbf{1.68\,mm} & \textbf{1.32\,mm} & \textbf{0.27}$^\circ$ &
      \textbf{1.65\,mm} & \textbf{1.58\,mm} & 21 of 27 \\
    \bottomrule
  \end{tabular}}

  \smallskip
  {\small Reading: at 0.1\,m/s$^2$ the real reversals at the far, high points ($x$ 0.60, $z$ 0.45\,m) lag 2.15\,mm,
  i.e.\ $\Lambda \approx 21.5$\,kg in Eq.~\eqref{eq:al-lag}, just above the criterion. At 20\,cm/s a full-speed move is
  about 42\,cm long and leaves the comfortable workspace: only 7 of 27 moves stay clear of the joint limits, and the
  counted ones lag up to 3.35\,mm ($^*$ start pose, $+y$, while braking far out; two further moves ran beyond the
  arm's reach). At 10\,cm/s all full-speed moves of the 7 counted start points fit. \textbf{0.08\,m/s$^2$ and
  10\,cm/s} (bold, chosen) meet the criterion everywhere with about 15\,\% margin. Rotations at the same limits: angle
  lag $\le 0.62^\circ$, tip drift $\le 1.23$\,mm (full-speed turns of about 120$^\circ$ about $x$ and $y$ reach the
  wrist's joint limits at most points, so only 1--2 of 3 turns per point are counted). Not counted: the two points
  near joint 4's limit lag up to 2.3\,mm.}
\end{table}

\clearpage
% ---------------------------------------------------------------------------------------------------------------------
\subsection{Design choices and their trade-offs}
\label{sec:al-choices}
These are the building blocks of the design. Every one moves the balance between \emph{transfer robustness}
(small lag, so that the motion is set by the controller), \emph{agility} (how fast the policy can move and react),
\emph{gentleness} (forces on the plant) and \emph{effort} (implementation, identification, training time). The
current value is given in the first column.

\paragraph{Controller (what turns the target into torques).}
\begin{center}
{\small\renewcommand{\arraystretch}{1.2}
\begin{tabular}{>{\raggedright\arraybackslash}p{0.16\textwidth}>{\raggedright\arraybackslash}p{0.23\textwidth}>{\raggedright\arraybackslash}p{0.245\textwidth}>{\raggedright\arraybackslash}p{0.245\textwidth}}
  \toprule
  Building block (current) & What it is & Gains & Costs \\
  \midrule
  Translational stiffness $K$ \newline (\textbf{1000\,N/m}, was 400)
    & spring between target and tool position. Franka's examples use 150--400\,N/m; Franka's internal Cartesian
      impedance accepts up to 3000\,N/m
    & lag $\propto 1/K$, so $a_\text{max} \propto K$ (Eq.~\ref{eq:al-amax}); less position error from joint friction on
      the real arm (error $\approx$ friction$/K$)
    & more force per mm of lag on contact; the damping grows with $\sqrt K$, so encoder noise and loop delay matter
      more; stability margin on the real robot to be checked \\
  \addlinespace\cmidrule(lr){1-4}\addlinespace
  Rotational stiffness $K_o$ \newline (\textbf{100\,N\,m/rad}, was 30)
    & spring between target and tool orientation (10--300\,N\,m/rad)
    & tilt $\propto 1/K_o$ while translating (Table~\ref{tab:al-diag}); smaller angle lag when turning
    & twists a stem caught in the fork more stiffly; same noise/stability caveat as $K$ \\
  \addlinespace\cmidrule(lr){1-4}\addlinespace
  Damping $D$ \newline (\textbf{$2\sqrt K\Lambda^{1/2}$})
    & damper between target and tool velocity
    & critical in every pose and direction; Franka's $2\sqrt K$ (critical only for 1\,kg) overshot by up to 7\,mm
    & uses the model of $\Lambda$ (only for the damping) \\
  \addlinespace\cmidrule(lr){1-4}\addlinespace
  Velocity feedforward $\dot x_d$ \newline (\textbf{on})
    & the damper acts on the velocity \emph{difference}, not on the tool velocity; while the clamp holds the target,
      $\dot x_d$ is the tool's velocity (the held target moves with the tool)
    & no lag at constant speed (without it: 82\,mm at 20\,cm/s, Table~\ref{tab:al-time}); no damping force on top of
      the clamp's bound
    & none found \\
  \addlinespace\cmidrule(lr){1-4}\addlinespace
  Inertia feedforward \newline $J^\top\Lambda\ddot x_d$ \newline (\textbf{off}; option B3)
    & the controller supplies the force needed to accelerate the arm from the model (operational-space
      control~\cite{al-khatib1987}) instead of waiting for the spring to stretch
    & lag drops to $\Delta\Lambda\, a / K$, set only by the \emph{model error} $\Delta\Lambda$ (20\,\% error: 5 times less
      lag); also removes the tilt and tip drift from the coupling
    & the motion now depends on the model of $\Lambda$ (libfranka on the robot); needs the motor-inertia question
      settled and the model checked (identification) \\
  \addlinespace\cmidrule(lr){1-4}\addlinespace
  Control period \newline (\textbf{2\,ms} sim, 1\,ms robot)
    & how often the law is evaluated
    & 1\,ms in sim matches the robot exactly
    & simulation twice as slow, so training twice as long \\
  \bottomrule
\end{tabular}}
\end{center}

\clearpage
\paragraph{Commands (what the policy outputs, and the limits on it).}
\begin{center}
{\small\renewcommand{\arraystretch}{1.2}
\begin{tabular}{>{\raggedright\arraybackslash}p{0.16\textwidth}>{\raggedright\arraybackslash}p{0.23\textwidth}>{\raggedright\arraybackslash}p{0.245\textwidth}>{\raggedright\arraybackslash}p{0.245\textwidth}}
  \toprule
  Building block (current) & What it is & Gains & Costs \\
  \midrule
  Speed cap \newline (\textbf{10\,cm/s, 45$^\circ$/s}; 20\,cm/s until 8 Oct.)
    & largest target step per policy step
    & the policy can approach and turn within one episode; 10\,cm/s is reached and left again within about 13\,cm,
      i.e.\ within the stem area
    & 20\,cm/s needed 42\,cm to reach and leave, more than the stem area, and lagged up to 3.35\,mm on such moves
      (Table~\ref{tab:al-final}) \\
  \addlinespace\cmidrule(lr){1-4}\addlinespace
  Acceleration limit \newline (\textbf{0.08\,m/s$^2$, 20$^\circ$/s$^2$}; 0.1\,m/s$^2$ until 8 Oct.)
    & largest change of the target step from one policy step to the next
    & \emph{the} knob for the lag (Eq.~\ref{eq:al-lag}); part of the policy-side code, so identical in sim and on the
      robot; the policy sees the applied step, so it can learn the limit
    & less agility: long time to speed and long braking distance (table in Section~\ref{sec:al-problem}) \\
  \addlinespace\cmidrule(lr){1-4}\addlinespace
  Lag criterion $e_\text{max}$ \newline (\textbf{2\,mm, 2$^\circ$})
    & how much lag we accept; sets the acceleration limit via Eq.~\eqref{eq:al-amax}
    & smaller: the motion depends less on the arm's dynamics (better transfer~\cite{al-aljalbout2024})
    & smaller $a_\text{max}$; a larger value (e.g.\ 5\,mm, option B2) gives 2.5 times more acceleration \\
  \addlinespace\cmidrule(lr){1-4}\addlinespace
  Low-pass filter on the commands \newline (\textbf{not used})
    & smooths the commands instead of limiting their change
    & smooth motion; common in sim-to-real
    & adds hidden state; Ref.~\cite{al-aljalbout2024} saw learning collapse with a limiter the policy could not see;
      no hard bound on the acceleration \\
  \addlinespace\cmidrule(lr){1-4}\addlinespace
  Target clamp \newline $\lVert x_d - x\rVert \le e_\text{clamp}$ \newline (\textbf{4\,mm, 3$^\circ$}, every control
    step)
    & the target is never further than $e_\text{clamp}$ from the measured tool; checked at every control step, not
      once per policy step (which would also cut steps in free motion)
    & hard bound on the spring force when the tool is blocked: 4\,mm $\times$ 1000\,N/m = 4\,N (about 3 times the
      measured push forces of 0.8--1.4\,N), 3$^\circ \times$ 100\,N\,m/rad = 5.2\,N\,m (43\,\% of the FR3's 12\,N\,m
      wrist torque limit; 10$^\circ$ would exceed it), even if the policy keeps pushing
    & must stay above the free-space lag (2\,mm), otherwise it acts while moving; while it acts, the target does not
      follow the commands (the policy sees this in the applied step); runs in the 1\,kHz loop on the robot \\
  \addlinespace\cmidrule(lr){1-4}\addlinespace
  Policy rate \newline (\textbf{31.25\,Hz})
    & how often the policy sends a new step
    & higher: finer commands, quicker reaction
    & more policy evaluations per second of simulation, so longer training; more sensitive to latency \\
  \addlinespace\cmidrule(lr){1-4}\addlinespace
  Action space \newline (\textbf{Cartesian steps on the previous target})
    & alternatives: joint velocities, Cartesian velocities, absolute poses, joint torques
    & joint velocities transferred best in Ref.~\cite{al-aljalbout2024} (stiff joint-level tracking)
    & joint velocities: no Cartesian compliance at contact, which matters for a fragile stem; torques: motion set by
      the dynamics alone, could not be transferred safely~\cite{al-aljalbout2024} \\
  \bottomrule
\end{tabular}}
\end{center}

\clearpage
\paragraph{Task and training.}
\begin{center}
{\small\renewcommand{\arraystretch}{1.2}
\begin{tabular}{>{\raggedright\arraybackslash}p{0.16\textwidth}>{\raggedright\arraybackslash}p{0.23\textwidth}>{\raggedright\arraybackslash}p{0.245\textwidth}>{\raggedright\arraybackslash}p{0.245\textwidth}}
  \toprule
  Building block (current) & What it is & Gains & Costs \\
  \midrule
  Episode length \newline (\textbf{10\,s})
    & how long the policy has per attempt
    & longer: a slow approach still leaves time to push the stem
    & training time grows about in proportion; the reward comes later after the actions that earned it, which makes
      learning harder \\
  \addlinespace\cmidrule(lr){1-4}\addlinespace
  Workspace bound \newline (\textbf{planned})
    & keep the tool away from poses near joint limits and the edge of the reach
    & there the lag grows and the real FR3 stops with a reflex
    & excludes some stem placements \\
  \addlinespace\cmidrule(lr){1-4}\addlinespace
  Reward terms \newline (\textbf{planned}: contact force, action rate, joint-limit margin)
    & penalties the policy learns to avoid, instead of hard limits
    & the policy learns gentle motion itself, where it matters
    & soft: no guarantee on the robot \\
  \addlinespace\cmidrule(lr){1-4}\addlinespace
  Domain randomization \newline (\textbf{planned}, after a first baseline run)
    & vary $K$, $K_o$, $\Lambda$ (motor inertia, link masses), latency and friction in training
    & policy robust to the remaining differences between sim and real~\cite{al-peng2018}
    & ranges must be realistic; too wide makes the policy cautious and slow to train \\
  \addlinespace\cmidrule(lr){1-4}\addlinespace
  Identification \newline (\textbf{planned}, before the first real run)
    & compare libfranka's model with the simulation; record the real arm's response to the same commands
    & tells how large the dynamics gap is; prerequisite for inertia feedforward and for loosening the criterion
    & robot time \\
  \addlinespace\cmidrule(lr){1-4}\addlinespace
  FR3 safety layer in sim \newline (\textbf{planned})
    & joint velocity, torque and torque-rate limits and the collision reflex as terminations/penalties
    & the policy does not learn motions the robot refuses
    & thresholds must match the robot's configuration \\
  \bottomrule
\end{tabular}}
\end{center}

\paragraph{Main options, in short.}
\begin{center}
{\renewcommand{\arraystretch}{1.2}
\begin{tabular}{>{\raggedright\arraybackslash}p{0.24\textwidth}>{\raggedright\arraybackslash}p{0.18\textwidth}>{\raggedright\arraybackslash}p{0.22\textwidth}>{\raggedright\arraybackslash}p{0.26\textwidth}}
  \toprule
  Option & $a_\text{max}$ at 2\,mm lag & Depends on the arm model & Main cost \\
  \midrule
  B1: $K$ 1000, $K_o$ 100, no inertia feedforward & 0.08\,m/s$^2$ (measured; 0.1 lags 2.15\,mm) & only through the
    damping & sluggish \\
  \addlinespace
  B2: B1 with a 5\,mm criterion & $\approx 0.25$\,m/s$^2$ & yes: 5\,mm of lag are set by the arm & against the transfer
    evidence \\
  \addlinespace
  B3: B1 + inertia feedforward & $\approx 0.5$\,m/s$^2$ or more (not run) & yes: model of $\Lambda$ & identification
    first \\
  \addlinespace
  Stiffer $K$ (3000\,N/m), no inertia feedforward & $\approx 0.3$\,m/s$^2$ & only through the damping & real-robot
    stability, contact forces \\
  \bottomrule
\end{tabular}}
\end{center}

% ---------------------------------------------------------------------------------------------------------------------
\subsection{Plan}
\label{sec:al-plan}
Decision (7 October 2026, revised 8 October): \textbf{B1 for the first training and transfer}, B3 (inertia
feedforward) once a policy trains successfully and the robot's $\Lambda$ has been checked. Values from
Table~\ref{tab:al-final}: $K = 1000$\,N/m, $K_o = 100$\,N\,m/rad, speed caps 10\,cm/s (3.2\,mm per step) and
45$^\circ$/s, acceleration limits 0.08\,m/s$^2$ (0.08\,mm per step) and 20$^\circ$/s$^2$ (0.02$^\circ$ per step),
target clamp 4\,mm and 3$^\circ$ at every control step. On 8 October the speed cap was halved and the acceleration
limit lowered from 0.1\,m/s$^2$, after moves that run up to the cap and reversals that really reverse showed lags of
up to 3.35\,mm and 2.15\,mm (Table~\ref{tab:al-final}).

\paragraph{Is this fast enough for the task?} With these limits (speed up, then brake, from rest to rest): the
farthest stem tip, 31\,cm from the start pose, is reached in 4.4\,s (1.25\,s to reach 10\,cm/s, 1.85\,s at
10\,cm/s, 1.25\,s braking); a 90$^\circ$ turn takes 4.3\,s (peak 42$^\circ$/s), and turning can happen while
approaching. That leaves about 5.5\,s of the 10\,s episode for pushing. Braking from 10\,cm/s takes 6\,cm and 1.25\,s.
We consider this sufficient for a first training. If the trained policy runs out of time, we lengthen the episode
first (cost: training time) and loosen the limits only after the identification.

% ---------------------------------------------------------------------------------------------------------------------
\subsection{Questions}
\label{sec:al-questions}
\begin{enumerate}
  \item Is the analysis right: is the lag of a Cartesian impedance controller on the FR3 set by its apparent mass
        (8--22\,kg at the tool in simulation, half of it from the motors' reflected inertia)?
  \item Does the FR3's internal joint torque controller reduce the apparent motor inertia (as on the DLR
        lightweight arms), and is that why libfranka's model leaves the motor inertia out? This decides how the
        simulation must model the motors, and which mass matrix the damping of our controller should use on the
        robot.
  \item Stiffness: is $K = 1000$\,N/m (and $K_o = 100$\,N\,m/rad) in a libfranka torque loop acceptable on our FR3
        next to a plant? Has the lab run similar gains, and with which velocity filtering?
  \item Inertia feedforward (option B3): how accurate is libfranka's model in your experience? Has the lab used
        operational-space control with acceleration feedforward on the FR3?
  \item Is an acceleration of about 0.08\,m/s$^2$ (1.25\,s to 10\,cm/s) a reasonable limit near a plant, or would
        you accept larger tracking errors for more agility?
  \item Which reflexes and thresholds (collision behaviour, joint velocity limits) does the lab configure, so that we
        model the same ones in simulation?
\end{enumerate}

{\raggedright
\begin{thebibliography}{9}
\bibitem{al-aljalbout2024} E.~Aljalbout, F.~Frank, M.~Karl, P.~van der Smagt. On the role of the action space in
  robot manipulation learning and sim-to-real transfer. \emph{IEEE Robotics and Automation Letters}, 2024.
  arXiv:2312.03673.
\bibitem{al-khatib1987} O.~Khatib. A unified approach for motion and force control of robot manipulators: the
  operational space formulation. \emph{IEEE Journal on Robotics and Automation}, 3(1), 1987.
\bibitem{al-peng2018} X.~B.~Peng, M.~Andrychowicz, W.~Zaremba, P.~Abbeel. Sim-to-real transfer of robotic control
  with dynamics randomization. \emph{IEEE ICRA}, 2018.
\end{thebibliography}
}
% ======================================== END CONTENT ========================================
) or compiled alone via action_limits_report.tex.
% Uses amsmath, booktabs and array (added to main.tex for this section). Numbers: scripts/sweep_action_poses.py
% (default and --full_speed), scripts/check_push_env.py and docs/notes/2026-10-07.md, 2026-10-08.md in the project repo.

% ======================================= BEGIN CONTENT =======================================
\section{Open problem: how fast may the policy move the fork?}
\label{sec:al}

\paragraph{Status (7--8 October 2026).} Problem identified in simulation; values for the first training runs set
on 7 October (option B1) and revised on 8 October after longer test moves (Section~\ref{sec:al-plan}); the
long-term solution is open. We ask for feedback on the analysis and on the
questions at the end (Section~\ref{sec:al-questions}).

\paragraph{In one paragraph.} The policy commands the fork through a Cartesian impedance controller, the same in
simulation and on the FR3. What blocks us is a trade-off between \emph{agility} and a \emph{small tracking error}.
A small error transfers better~\cite{al-aljalbout2024}, because then the fork's motion is set by our controller
(identical on both sides) rather than by the arm's dynamics (not identical, not yet identified). But an impedance
controller is a spring of stiffness $K$ between the commanded and the actual tool position, and it can only
accelerate the arm by being stretched: while the commanded motion accelerates at $a$, the tool lags behind by about
$\Lambda a/K$, where $\Lambda$ is the arm's apparent mass at the tool tip (8--22\,kg depending on pose and direction,
half of it from the motors' reflected inertia). With a small lag (2\,mm) and a stiffness we consider acceptable near a
plant (400--1000\,N/m), the fork may accelerate only at about 0.04--0.09\,m/s$^2$; we use 0.08\,m/s$^2$ at
1000\,N/m. This section explains the
trade-off, what we measured, and every design choice that moves it.

% ---------------------------------------------------------------------------------------------------------------------
\subsection{Background}
\label{sec:al-background}

\paragraph{What the policy does.} The policy is a neural network trained by reinforcement learning (RL) in
simulation. At 31.25\,Hz (every 32\,ms, one ``policy step'') it reads the state (fork pose, stem shape, target) and
outputs a small displacement and rotation of the fork's \emph{target} pose $x_d$. The next target is the previous
target plus this step, so commanded steps are not lost even if the tool lags behind. Within the 32\,ms the target
moves at constant velocity $\dot x_d$ from the old to the new value. A Cartesian impedance controller at 1\,kHz
(2\,ms in simulation) turns the target into joint torques:
\begin{equation}
  \tau = J^\top\!\left[\,K\,(x_d - x) + D\,(\dot x_d - \dot x)\,\right] + \tau_\text{null} + \tau_\text{Coriolis}.
  \label{eq:al-law}
\end{equation}
$x$ is the measured tool pose, $J$ the tool-tip Jacobian, $K$ the stiffness and $D$ the damping (written here for
the position; the orientation has its own stiffness $K_o$). This is the law of Franka's example controllers, with
two changes: the target velocity $\dot x_d$ is fed forward (Franka's examples use $\dot x_d = 0$), and the damping
$D = 2\sqrt{K}\,\Lambda^{1/2}$ is chosen critical for the arm's actual inertia $\Lambda$ (Section~\ref{sec:al-problem})
instead of Franka's $2\sqrt K$, which is critical only for a 1\,kg tool.

\paragraph{Null-space spring $\tau_\text{null}$.} The FR3 has 7 joints, but the tool pose has only 6 degrees of
freedom. For every tool pose there is therefore a one-parameter family of arm postures (the elbow can swing about the
line from shoulder to wrist without moving the tool): the \emph{null space}. The first term of
Eq.~\eqref{eq:al-law} does not act in this direction, so the elbow would drift freely. A weak joint-space spring
(0.5\,N\,m/rad, as in Franka's examples) pulls the posture towards a rest posture; it is projected so that it does
not push the tool. It plays no role in the problem below (we checked: switching it off does not change the lag).

\paragraph{What RL learns, and why the action space matters for the transfer.} During training the policy tries
millions of commands in simulation and keeps what earns reward. It learns \emph{how the simulated fork reacts to its
commands}, including any lag, overshoot or coupling, and relies on it. If the real robot reacts differently, the
learned behaviour is wrong by that difference: this is the sim-to-real (``reality'') gap. The policy cannot correct
what it never saw in simulation. The usual tools against the gap are:
\begin{center}
{\renewcommand{\arraystretch}{1.15}
\begin{tabular}{>{\raggedright\arraybackslash}p{0.24\textwidth}>{\raggedright\arraybackslash}p{0.42\textwidth}>{\raggedright\arraybackslash}p{0.24\textwidth}}
  \toprule
  Tool & What it does & Status \\
  \midrule
  A. Action-space design & choose commands whose result depends little on what differs between sim and real
    (same controller, small tracking error) & controller matched; this section \\
  \addlinespace
  B. Identification & measure the real robot and correct the simulation & planned, before the first real run \\
  \addlinespace
  C. Domain randomization & vary unknown parameters in training, so the policy works over a range & planned \\
  \addlinespace
  D. Model the robot's safety layer & limits and reflexes of the FR3 become terminations/penalties in sim, so the
    policy does not learn motions the robot refuses & planned \\
  \addlinespace
  E. Adaptation & infer differences online or fine-tune on the robot & not planned \\
  \bottomrule
\end{tabular}}
\end{center}
A large study of action spaces on a Franka arm~\cite{al-aljalbout2024} (13 action spaces, about 250 policies,
reaching and pushing, transferred to the real robot) concludes that \emph{``an action space which yields or
naturally limits the tracking error transfers better''}; torque actions, whose result is set by the robot's
dynamics alone, could not be transferred safely. The same study warns that a limiter which acts on the robot but is
not visible to the policy in simulation (they tried a rate limiter) made learning collapse. A caveat: their controllers were joint
impedance controllers (Cartesian actions were converted to joint targets by inverse kinematics), with gains that
the paper does not report, and the tasks were reaching and pushing a box. The conclusion carries over in principle,
since in every impedance-type controller a larger tracking error means more dependence on the arm's dynamics; the
numbers do not, and the paper does not test a soft Cartesian impedance controller like ours.

\paragraph{Where the gap enters our system.} The chain from a policy output to the stem, step by step:
\begin{center}
{\renewcommand{\arraystretch}{1.15}
\begin{tabular}{>{\raggedright\arraybackslash}p{0.32\textwidth}>{\raggedright\arraybackslash}p{0.20\textwidth}>{\raggedright\arraybackslash}p{0.38\textwidth}}
  \toprule
  Step & Same in sim and real? & Source of differences \\
  \midrule
  policy step $\to$ target pose & yes (same code) & none \\
  target $\to$ torques, law~\eqref{eq:al-law} & yes (ported to C++, checked) & port errors; $J$, $M$ from different
    models \\
  \textbf{torques $\to$ fork motion} & \textbf{no} & arm inertia, joint friction, motor inertia, delays, 1 vs.\ 2\,ms
    control period \\
  fork $\leftrightarrow$ stem contact & no & friction, stem stiffness and damping (randomized) \\
  stem $\to$ observation & no & camera noise, latency \\
  FR3 safety layer & \textbf{missing in sim} & reflexes stop the robot \\
  \bottomrule
\end{tabular}}
\end{center}
The step ``torques $\to$ fork motion'' is where the arm's dynamics act. How much of it reaches the fork's motion
depends on the tracking error: if the tool follows the target within a small error, its motion is essentially the
target's motion, which our code computes identically in sim and real; the larger the error, the more the motion is
shaped by the arm's dynamics, which differ.

% ---------------------------------------------------------------------------------------------------------------------
\subsection{The problem}
\label{sec:al-problem}

\paragraph{What blocks us: agility versus a small tracking error.} We want two things from the fork's motion:
\begin{enumerate}
  \item \textbf{Small lag behind the commanded target}, so that the fork's motion is set by our controller, which
        is the same in simulation and on the FR3, and not by the arm's dynamics, which differ and are not yet
        identified. Action spaces that keep the tracking error small transfer better~\cite{al-aljalbout2024}.
  \item \textbf{Agility}: reach the stem well within the 10\,s episode, brake late before contact, and react quickly
        to how the stem moves.
\end{enumerate}
With an impedance controller the two conflict. The lag grows with how quickly the commanded motion accelerates, so
a small lag allows only gentle accelerations. In numbers: with $K = 1000$\,N/m and at most 2\,mm of lag, the fork may
accelerate at only about 0.09\,m/s$^2$ where the arm is heaviest; at the 0.08\,m/s$^2$ we use, it needs 1.25\,s to
reach 10\,cm/s and 6\,cm to stop from that speed (2.5\,s and 25\,cm for 20\,cm/s). Either
we accept this sluggishness, or we accept more lag (and with it a larger dependence on the real arm), or we change
the controller. This is the decision we need help with. The rest of this section explains where the conflict comes
from, how severe it is for the task, and how severe it will be on the real robot.

\paragraph{Why the tool lags: the impedance law as a mass on a spring.} Consider one direction, free space,
gravity compensated. The controller makes the tool tip behave like a mass $\Lambda$ (the arm's inertia felt at the
tool tip, explained in the next paragraph) attached by a spring $K$ and a damper $D$ to the moving target:
\begin{equation}
  \Lambda\,\ddot x = K\,(x_d - x) + D\,(\dot x_d - \dot x).
\end{equation}
Write the lag as $e = x_d - x$ (target minus tool). Then $\ddot x = \ddot x_d - \ddot e$, and substituting gives
\begin{equation}
  \Lambda\,\ddot e + D\,\dot e + K\,e = \Lambda\,\ddot x_d.
  \label{eq:al-error}
\end{equation}
The lag itself behaves like a damped mass-spring system, driven by the force $\Lambda \ddot x_d$, i.e.\ by the
\emph{acceleration of the target}. Three cases follow:
\begin{enumerate}
  \item \textbf{Target at rest or moving at constant speed} ($\ddot x_d = 0$): the right-hand side is zero, so any lag
        dies out and $e \to 0$. Physically: moving at constant speed in free space needs no force (no friction, no
        gravity), so the spring does not need to be stretched. This holds only because the target velocity is fed
        forward. Without it the damper pulls back with $D\dot x$ all the time, and the spring must balance it:
        $e = D\dot x_d / K$. With $D = 2\sqrt{K\Lambda} = 180$\,N\,s/m ($K = 400$\,N/m, $\Lambda = 20$\,kg) and
        20\,cm/s this gives 90\,mm; we measured 82\,mm (Table~\ref{tab:al-time}).
  \item \textbf{Target accelerating at constant $a$}: the right-hand side is the constant force $\Lambda a$. Once the
        transient has died out ($\ddot e = \dot e = 0$), $K e = \Lambda a$, so
        \begin{equation}
          e = \frac{\Lambda\, a}{K}.
          \label{eq:al-lag}
        \end{equation}
        Physically: Newton's law. To accelerate the mass $\Lambda$ at $a$, the force $\Lambda a$ is needed, and its only
        source is the stretched spring. Example: $\Lambda = 20$\,kg, $a = 0.1$\,m/s$^2$, $K = 1000$\,N/m: $e = 2$\,mm.
        The same holds when braking (the tool then runs ahead of the target, which shows up as overshoot).
  \item \textbf{How fast the lag builds up and dies out}: Eq.~\eqref{eq:al-error} is critically damped, with the time
        constant $\sqrt{\Lambda/K}$: 0.14\,s for 20\,kg and 1000\,N/m, 0.22\,s for 400\,N/m, i.e.\ 4--7 policy steps.
        The lag reaches the value of Eq.~\eqref{eq:al-lag} if the acceleration lasts longer than about three time
        constants, and it disappears again within a similar time once the target moves at constant speed. This
        matters for testing: short test moves end before the lag has settled. Our first sweep (20 policy steps per
        segment, reversals that only braked) therefore showed 1.9\,mm at 0.1\,m/s$^2$; moves that brake \emph{and}
        reverse, or that run up to the speed cap, accelerate long enough and reached 2.15\,mm
        (Table~\ref{tab:al-final}).
\end{enumerate}

\paragraph{The mass in this equation: the apparent mass $\Lambda$ (``8--20\,kg at the tool tip'').} The trade-off
is severe because the mass that the spring has to accelerate is large. If one pushes the tool tip of
the arm (joints free, gravity compensated) with a force $F$, the tip starts to accelerate with $\ddot x = \Lambda^{-1}
F$. $\Lambda$ is the arm's inertia as seen at the tool tip, the \emph{apparent mass} (or operational-space inertia):
\begin{equation}
  \Lambda = \left(J M^{-1} J^\top\right)^{-1},
  \label{eq:al-lambda}
\end{equation}
with $M$ the $7\times7$ joint-space mass matrix and $J$ the $6\times7$ tool-tip Jacobian. $\Lambda$ is a $6\times6$
matrix (force and moment versus linear and angular acceleration) and depends on the pose. Its translational
diagonal entries say how heavy the arm feels when the tip is pushed along $x$, $y$ or $z$. At the earlier start pose
(2026-10-07; tool tip at $(0.30, 0, 0.55)$\,m in front of the robot, fork horizontal), computed from the simulated FR3:
\begin{center}
{\renewcommand{\arraystretch}{1.15}
\begin{tabular}{lrrr}
  \toprule
  Apparent mass, start pose & along $x$ (forward) & along $y$ (sideways) & along $z$ (up) \\
  \midrule
  links only & 11.9\,kg & 4.2\,kg & 5.1\,kg \\
  links + motors' reflected inertia (our simulation) & 20.5\,kg & 7.6\,kg & 10.6\,kg \\
  \bottomrule
\end{tabular}}
\end{center}
Pushing the tip forward with 20\,N accelerates it at about 1\,m/s$^2$, as if it were a 20\,kg mass; the same push
sideways gives 2.6\,m/s$^2$. Forward motion needs the heavy shoulder and elbow joints, sideways motion mostly the
base rotation and the wrist. The second row includes the motors' reflected inertia $n^2 I_\text{motor}$ ($n$: gear
ratio), which we add to $M$ from Franka's description files: 0.61, 0.61, 0.46, 0.46, 0.21, 0.21, 0.21\,kg\,m$^2$ for
joints 1--7 (at the wrist about 100 times the link's own inertia). Without these terms the simulated arm oscillated
and pressed the stem 5 times harder. The lag measurements below give the same value independently: from
Eq.~\eqref{eq:al-lag}, $\Lambda = K e / a \approx 16$--20\,kg along $x$ at the start pose, and up to 21.5\,kg
at the far, high start points of the sweep ($x$ 0.60, $z$ 0.45\,m).

\paragraph{How severe the trade-off is for the task.} Turning Eq.~\eqref{eq:al-lag} around: for a largest acceptable lag $e_\text{max}$, the
target may accelerate at most at
\begin{equation}
  a_\text{max} \approx \frac{K\, e_\text{max}}{\Lambda}.
  \label{eq:al-amax}
\end{equation}
With $\Lambda = 21.5$\,kg (the heaviest pose and direction) and $e_\text{max} = 2$\,mm: 0.04\,m/s$^2$ at
$K = 400$\,N/m and 0.09\,m/s$^2$ at 1000\,N/m. We use 0.08\,m/s$^2$, which leaves about 15\,\% margin (lag
$\le 1.72$\,mm; measured 1.68\,mm, Table~\ref{tab:al-final}). What this means for the task (the
target speeds up at $a_\text{max}$ until the speed cap, then brakes at $a_\text{max}$):
\begin{center}
{\renewcommand{\arraystretch}{1.15}
\begin{tabular}{rrrrr}
  \toprule
  largest acceleration & time from rest & braking distance & 31\,cm, rest to rest & 31\,cm, rest to rest \\
  $a_\text{max}$ & to 20\,cm/s & from 20\,cm/s & (20\,cm/s cap) & (10\,cm/s cap) \\
  \midrule
  0.05\,m/s$^2$ & 4.0\,s & 40\,cm & 5.0\,s & 5.1\,s \\
  \textbf{0.08\,m/s}$^\mathbf{2}$ & 2.5\,s & 25\,cm & 3.9\,s & \textbf{4.4\,s} \\
  0.10\,m/s$^2$ & 2.0\,s & 20\,cm & 3.5\,s & 4.1\,s \\
  0.25\,m/s$^2$ & 0.8\,s & 8\,cm  & 2.3\,s & 3.5\,s \\
  0.50\,m/s$^2$ & 0.4\,s & 4\,cm  & 1.9\,s & 3.3\,s \\
  \bottomrule
\end{tabular}}
\end{center}
31\,cm is the distance from the start pose to the farthest stem tip. The episode lasts 10\,s, and the stem must
still be pushed after the approach. The cost of a small $a_\text{max}$ is therefore not the top speed but agility:
the policy must start braking long before it touches the stem and cannot react quickly to how the stem moves.

\paragraph{Why not simply make the episode longer?} A longer episode gives a slow policy time to finish, but it is
not free. RL learns from simulated experience, so the training time grows roughly in proportion to the simulated
time per episode: a 20\,s episode costs about twice the training time of a 10\,s episode for the same number of
attempts. In addition, the reward for reaching the target arrives later after the actions that caused it, which
makes it harder for the policy to learn which actions were good. We therefore keep 10\,s as long as the task fits,
and lengthen it only if the trained policy demonstrably runs out of time.

\paragraph{Rotations: the same effect, plus a coupling.} Turning the fork works in the same way: the orientation
spring $K_o$ must twist to give the rotational inertia an angular acceleration $\alpha$, so the angle lags by about
$\Lambda_o \alpha / K_o$ ($\Lambda_o$: the rotational part of $\Lambda$). In addition, $\Lambda$ is a full $6\times6$
matrix: its off-diagonal blocks couple translation and rotation. The reason is geometric: the tool tip is 10--20\,cm
away from the wrist's joint axes, and the joints that turn the fork also move the tip, while those that move the
tip also turn the fork. The control law, however, treats position and orientation as two independent springs.
Two consequences, both measured:
\begin{itemize}
  \item \textbf{Tilt while translating.} Accelerating the tip at $a$ needs, besides the force, a moment
        $\Lambda_{op}\, a$ ($\Lambda_{op}$: off-diagonal block). It comes from the orientation spring, which must twist:
        the fork tilts by about $\Lambda_{op} a / K_o$. Table~\ref{tab:al-diag} shows the tilt falling with $1/K_o$
        (2.0$^\circ$ at $K_o = 30$, 0.26$^\circ$ at 300\,N\,m/rad).
  \item \textbf{Tip drift while turning.} Turning the fork at angular acceleration $\alpha$ needs, besides the moment, a
        force at the tip, which the position spring supplies by stretching: the tip moves away from its target by
        about $\Lambda_{po}\,\alpha / K$, although the position command is constant (Tables~\ref{tab:al-rot}
        and~\ref{tab:al-b1}).
\end{itemize}
Both effects scale with the acceleration, like the lag. A controller with inertia feedforward (operational-space
control, option B3 below) would supply these coupling forces and moments from the model and remove both.

\paragraph{Contact.} When the tool is blocked (e.g.\ pressed against the stem), the target keeps moving and the
lag becomes force: $K e$. Small lags therefore also mean small unintended forces. A hard bound is added by
clamping the target: at every control step, $\lVert x_d - x \rVert \le 4$\,mm (spring force $\le 4$\,N at
$K = 1000$\,N/m) and the angle between target and tool orientation $\le 3^\circ$ (moment $\le 5.2$\,N\,m at
$K_o = 100$\,N\,m/rad). In free motion the clamp is inactive, because the lag stays below 2\,mm. While it holds the
target, the target moves with the tool, so the target velocity fed forward is the tool's own velocity: the damper
then adds no force to the bounded spring force, and the controller does not switch back and forth between holding and
releasing. Not bounded by the clamp: the momentum of the arm when the moving fork first hits something (as with any
impedance controller); the policy is meant to learn to approach gently (contact-force and smoothness penalties).

\paragraph{How severe the trade-off will be on the real robot: which $\Lambda$ does the FR3 have?} All numbers above
use the simulated $\Lambda$. The real arm's $\Lambda$ is not known yet, and two observations from Franka's software
suggest that it may differ:
\begin{itemize}
  \item libfranka (current version, September 2026) computes its model (\texttt{franka::Model::mass}) with the
        Pinocchio library from the URDF that the robot provides. The FR3 URDF shipped with libfranka's tests has
        the same link inertias as Franka's description files but \emph{no} motor-inertia entries, and libfranka never
        sets Pinocchio's rotor-inertia (``armature'') term. So libfranka's mass matrix most likely \emph{excludes}
        the motors' reflected inertia. (Read in the source code, not yet checked on our robot.)
  \item The FR3 executes commanded torques with an internal joint torque controller that uses the link-side
        torque sensors. On the DLR lightweight arms, from which Franka's design derives, such torque feedback
        reduces the apparent motor inertia. If the FR3 does the same, the real arm feels lighter than our
        simulation (between the two rows of the table above), and libfranka's model would be the right one.
\end{itemize}
How much this matters depends on the option chosen (Section~\ref{sec:al-choices}):
\begin{center}
{\renewcommand{\arraystretch}{1.15}
\begin{tabular}{>{\raggedright\arraybackslash}p{0.24\textwidth}>{\raggedright\arraybackslash}p{0.66\textwidth}}
  \toprule
  Option & Role of the real $\Lambda$ \\
  \midrule
  B1 (small lag, now) & \textbf{Little}, and this is exactly why small lags transfer well: the unknown $\Lambda$
    acts on the motion only through the lag. The acceleration limit was sized for $\le 2$\,mm with the simulated
    $\le 21.5$\,kg; a real arm of 12\,kg would lag 1\,mm, a heavier one (30\,kg) 2.4\,mm. The policy is then wrong by about a
    millimetre. \textbf{One exception: the damping} $D = 2\sqrt K \Lambda^{1/2}$ is computed from a model of $\Lambda$.
    Before the first real run we must decide which mass matrix the controller on the robot uses (libfranka's, or
    libfranka's plus the motor terms); with the wrong one the tool is under- or over-damped. \\
  \addlinespace
  B2 (larger lag) & \textbf{More}: with 5\,mm of lag set by the arm, a 40\,\% error in $\Lambda$ shifts the motion by
    2\,mm. \\
  \addlinespace
  B3 (inertia feedforward, later) & \textbf{Central}: the controller supplies the accelerating force from the model
    of $\Lambda$, so the remaining lag is set by the model error. Requires the identification first. \\
  \bottomrule
\end{tabular}}
\end{center}

% ---------------------------------------------------------------------------------------------------------------------
\subsection{Measurements}
\label{sec:al-measurements}

\paragraph{Setup.} Simulation (Isaac Lab, PhysX) of the FR3 with the fork, gravity compensated, no joint friction,
control period 2\,ms, policy step 32\,ms, no stem (free space). Every test commands the fork with the same commands
a policy could send, through the full controller.

\paragraph{Commands and limits.} A command is a target step per policy step. Three limits act on it:
\begin{itemize}
  \item \emph{speed cap}: the largest step, e.g.\ 3.2\,mm per 32\,ms = 10\,cm/s (rotation: 1.44$^\circ$ = 45$^\circ$/s);
  \item \emph{acceleration limit}: the largest change of the step from one policy step to the next. A change of
        0.1\,mm per step means $a = 0.1\,\text{mm}/(32\,\text{ms})^2 \approx 0.1$\,m/s$^2$, 0.08\,mm per step
        0.08\,m/s$^2$ (rotation: 0.02$^\circ$ per step $\approx$ 20$^\circ$/s$^2$). Commanding ``full speed'' from rest therefore makes the target speed up at
        this acceleration until the speed cap; commanding ``stop'' makes it brake at the same rate;
  \item the first design (Table~\ref{tab:al-first}) had a speed cap only, so the target jumped from rest to full
        speed in one step.
\end{itemize}

\paragraph{Test moves.} \emph{Straight move}: full command along one axis of the robot base frame ($\pm x$, $\pm y$,
$\pm z$) for a number of policy steps, then zero command (``hold'') for 45 steps (1.4\,s), during which the target
brakes and stops. \emph{Reversal}: full command along $+x$, then immediately along $-x$ (same for $y$, $z$), then hold;
this contains the largest change of velocity. Since 8 October the second segment is twice as long, so that the
target brakes \emph{and} reverses (before, it only braked to rest). \emph{Full-speed move} (since 8 October): from
rest up to the speed cap, 10 steps at the cap, then braking to rest (about 16\,cm; rotations about 120$^\circ$),
one per axis, in a direction with room for it (towards the middle of the workspace; up along $z$).
\emph{Rotation moves}: the same about the base axes, with zero position command. \emph{Start points}: either the
start pose only, or 9 points spanning the workspace in which the stems stand (tool tip at $x$ = 0.35 and 0.60\,m,
$y = \pm0.20$\,m, $z$ = 0.20 and 0.45\,m, plus the start pose). Moves that come within 0.25\,rad of a joint limit
are reported but not counted (decision 8 October): there the arm, not the controller, limits the motion. This
concerns the two start points low and close to the robot ($x$ 0.35, $z$ 0.20\,m), where joint 4 is within
0.21\,rad of its limit.

\paragraph{What is measured (all tables).}
\begin{itemize}
  \item \emph{Lag} (following error): the largest distance between the target position and the actual tool-tip
        position during the move and the hold. The quantity of Eq.~\eqref{eq:al-lag}.
  \item \emph{Overshoot}: how far the tool tip goes past the target's final position after it stops.
  \item \emph{Tilt}: the largest change of the fork's orientation during a straight move, although no rotation is
        commanded.
  \item \emph{Angle lag}: during rotation moves, the largest angle between the target and the actual orientation.
  \item \emph{Tip drift}: during rotation moves, the largest displacement of the tool tip, although no translation is
        commanded.
\end{itemize}
Each table reports the worst value over all moves (and start points). The criterion used so far: lag $\le 2$\,mm,
overshoot $\le 5$\,mm, tilt $\le 1^\circ$; angle lag $\le 2^\circ$, tip drift $\le 2$\,mm.

\begin{table}[h]
  \centering
  \caption{\textbf{First design}: speed cap only, no acceleration limit, $K = 400$\,N/m, $K_o = 30$\,N\,m/rad. Straight
  moves of 30 policy steps at full command from 9 start points, so each move starts with a jump from rest to full speed.
  Lag measured over the last quarter of each move (as defined at the time), overshoot after it. Worst case over all
  moves and start points. This design kept the lag below 2\,mm by limiting the speed to 5.5\,cm/s.}
  \label{tab:al-first}
  {\renewcommand{\arraystretch}{1.15}
  \begin{tabular}{rrrr}
    \toprule
    speed cap (step per 32\,ms) & speed & lag & overshoot \\
    \midrule
    2.0\,mm  & 6.3\,cm/s & 2.25\,mm & 3.9\,mm \\
    1.75\,mm & 5.5\,cm/s & 1.94\,mm & 3.4\,mm \\
    1.5\,mm  & 4.7\,cm/s & 1.63\,mm & 2.9\,mm \\
    \bottomrule
  \end{tabular}}

  \smallskip
  {\small Reading: with a cap of 5.5\,cm/s and 7.8$^\circ$/s the approach to a far stem took more than half of the
  episode, and a 90$^\circ$ turn did not fit in one. We therefore raised the caps to 20\,cm/s and 45$^\circ$/s and
  added the acceleration limit; Tables~\ref{tab:al-time}--\ref{tab:al-b1} show what that does.}
\end{table}

\begin{table}[h]
  \centering
  \caption{\textbf{How the lag evolves over time} (illustrates Eq.~\eqref{eq:al-error}). Start pose, $K = 400$\,N/m,
  $K_o = 30$\,N\,m/rad. One straight move along $+x$ at full command for 40 policy steps. With the acceleration limit
  of 0.49\,m/s$^2$ the target speeds up during the first 13 steps (0.42\,s) and then moves at constant speed.
  Lag of the tool tip behind the target after 10, 20, 30 and 40 policy steps.}
  \label{tab:al-time}
  {\renewcommand{\arraystretch}{1.15}
  \begin{tabular}{lrrrr}
    \toprule
    & \multicolumn{4}{c}{lag after $n$ policy steps} \\
    \cmidrule(lr){2-5}
    & 10 (0.32\,s) & 20 (0.64\,s) & 30 (0.96\,s) & 40 (1.28\,s) \\
    & accelerating & constant speed & constant speed & constant speed \\
    \midrule
    20\,cm/s, velocity feedforward on  & 9.4\,mm & 10.4\,mm & 4.7\,mm & 1.9\,mm \\
    10\,cm/s, velocity feedforward on  & 7.7\,mm & 4.7\,mm  & 2.0\,mm & 0.7\,mm \\
    20\,cm/s, velocity feedforward off & 23\,mm  & 59\,mm   & 75\,mm  & 82\,mm \\
    \bottomrule
  \end{tabular}}

  \smallskip
  {\small With feedforward the lag builds up while the target accelerates and dies out at constant speed (case 1 and 2
  of Section~\ref{sec:al-problem}); without it, the lag grows towards $D\dot x_d/K$. Switching off the null-space
  spring gave the same numbers.}
\end{table}

\begin{table}[h]
  \centering
  \caption{\textbf{Effect of the stiffnesses $K$ and $K_o$.} Start pose, speed cap 20\,cm/s. A straight move along $+x$
  at full command for 25 policy steps (0.8\,s, enough to reach about the speed cap), hold for 45 steps, then the same
  along $-x$ and hold. Two acceleration limits (columns). Largest lag of the tool tip and largest tilt of the fork
  over the whole test.}
  \label{tab:al-diag}
  {\renewcommand{\arraystretch}{1.15}
  \begin{tabular}{rrrrrr}
    \toprule
    & & \multicolumn{2}{c}{acceleration 0.24\,m/s$^2$} & \multicolumn{2}{c}{acceleration 0.49\,m/s$^2$} \\
    \cmidrule(lr){3-4} \cmidrule(lr){5-6}
    $K$ [N/m] & $K_o$ [N\,m/rad] & lag & tilt & lag & tilt \\
    \midrule
    400  & 30  & 9.7\,mm  & 1.80$^\circ$ & 13.6\,mm & 2.40$^\circ$ \\
    400  & 100 & 10.5\,mm & 0.66$^\circ$ & 14.4\,mm & 0.84$^\circ$ \\
    400  & 300 & 10.6\,mm & 0.23$^\circ$ & 14.8\,mm & 0.29$^\circ$ \\
    1000 & 30  & 4.0\,mm  & 2.03$^\circ$ & 6.4\,mm  & 2.69$^\circ$ \\
    1000 & 100 & 4.6\,mm  & 0.73$^\circ$ & 7.1\,mm  & 1.00$^\circ$ \\
    1000 & 300 & 4.8\,mm  & 0.26$^\circ$ & 7.5\,mm  & 0.37$^\circ$ \\
    \bottomrule
  \end{tabular}}

  \smallskip
  {\small Reading: the lag depends on $K$ (2.5 times stiffer, about 2.4 times less lag) and on the acceleration, hardly
  on $K_o$; the tilt depends on $K_o$ (10 times stiffer, about 8 times less tilt). Every entry fails the 2\,mm criterion.}
\end{table}

\begin{table}[h]
  \centering
  \caption{\textbf{Workspace sweep at $K = 1000$\,N/m, $K_o = 30$\,N\,m/rad}, caps 20\,cm/s and 45$^\circ$/s. At each
  start point: straight moves along $\pm x, \pm y, \pm z$ and reversals in $x, y, z$ (30 policy steps per segment, then
  hold), then the same as rotations about the base axes. Each row is one pair of acceleration limits (translation,
  rotation). Worst case over all moves and over the five start points closest to the robot: from the start points at
  $x = 0.60$\,m these moves of up to 20\,cm ran to the edge of the arm's reach, which is a flaw of the test, not of the
  controller.}
  \label{tab:al-rot}
  {\renewcommand{\arraystretch}{1.15}
  \begin{tabular}{rrrrrr}
    \toprule
    \multicolumn{3}{c}{translation moves} & \multicolumn{3}{c}{rotation moves} \\
    \cmidrule(lr){1-3} \cmidrule(lr){4-6}
    acceleration limit & lag & tilt & angular acceleration limit & angle lag & tip drift \\
    \midrule
    0.24\,m/s$^2$ & 5.0\,mm  & 3.0$^\circ$ & 20$^\circ$/s$^2$  & 9.9$^\circ$$^*$ & 14.0\,mm$^*$ \\
    0.49\,m/s$^2$ & 8.5\,mm  & 5.0$^\circ$ & 49$^\circ$/s$^2$  & 4.7$^\circ$ & 4.0\,mm \\
    0.73\,m/s$^2$ & 11.0\,mm & 5.7$^\circ$ & 98$^\circ$/s$^2$  & 6.8$^\circ$ & 7.0\,mm \\
    0.98\,m/s$^2$ & 12.6\,mm & 6.1$^\circ$ & 195$^\circ$/s$^2$ & 9.4$^\circ$ & 14.3\,mm \\
    \bottomrule
  \end{tabular}}

  \smallskip
  {\small Reading: all rows fail the criterion; the lag grows with the acceleration as expected, the tilt is large
  because $K_o$ is soft. $^*$ one start point ($x$ 0.35, $y$ 0.20, $z$ 0.45\,m), reversal about $x$; at the other four
  points at most 4.1$^\circ$ / 3.7\,mm. Not yet understood (likely a wrist joint near its limit).}
\end{table}

\begin{table}[h]
  \centering
  \caption{\textbf{Workspace sweep for the interim choice of 7 October: $K = 1000$\,N/m, $K_o = 100$\,N\,m/rad}, caps
  20\,cm/s and 45$^\circ$/s, smaller acceleration limits. Same moves as Table~\ref{tab:al-rot} but 20 policy steps per
  segment (moves of at most about 6\,cm, so all 9 start points are valid). Worst case over all moves and all 9 start
  points. \emph{Superseded by Table~\ref{tab:al-final}}: these moves reached at most 2\,mm per step (6.25\,cm/s), not
  the caps, and the reversals only braked, so the lag had not settled.}
  \label{tab:al-b1}
  {\renewcommand{\arraystretch}{1.15}
  \begin{tabular}{rrrrrrr}
    \toprule
    \multicolumn{4}{c}{translation moves} & \multicolumn{3}{c}{rotation moves} \\
    \cmidrule(lr){1-4} \cmidrule(lr){5-7}
    acceleration limit & lag & overshoot & tilt & angular accel.\ limit & angle lag & tip drift \\
    \midrule
    0.05\,m/s$^2$ & 1.3\,mm & 0.9\,mm & 0.17$^\circ$ & 10$^\circ$/s$^2$ & 0.30$^\circ$ & 0.55\,mm \\
    \textbf{0.10\,m/s}$^\mathbf{2}$ & \textbf{1.9\,mm} (2.4$^*$) & \textbf{1.8\,mm} & \textbf{0.33}$^\circ$ &
      \textbf{20}$^\circ$\textbf{/s}$^\mathbf{2}$ & \textbf{0.51}$^\circ$ & \textbf{1.35\,mm} \\
    0.15\,m/s$^2$ & 2.9\,mm (3.5$^*$) & 2.6\,mm & 0.47$^\circ$ & 39$^\circ$/s$^2$ & 1.31$^\circ$ & 4.41\,mm \\
    \bottomrule
  \end{tabular}}

  \smallskip
  {\small Reading: the criterion (lag $\le 2$\,mm, tilt $\le 1^\circ$, angle lag $\le 2^\circ$, tip drift $\le 2$\,mm)
  is met at 0.1\,m/s$^2$ and 20$^\circ$/s$^2$ (bold, chosen) and fails at 0.15\,m/s$^2$ / 39$^\circ$/s$^2$. The lag
  of 1.9\,mm matches Eq.~\eqref{eq:al-lag} with $\Lambda \approx 20$\,kg. $^*$ at the two start points low and close
  to the robot ($x$ 0.35, $z$ 0.20\,m), where joint 4 is within 0.21\,rad of its limit; the task will keep the tool
  away from there.}
\end{table}

\begin{table}[h]
  \centering
  \caption{\textbf{Revised sweep (8 October): real reversals and full-speed moves.} $K = 1000$\,N/m,
  $K_o = 100$\,N\,m/rad, rotation limits unchanged (45$^\circ$/s, 20$^\circ$/s$^2$). \emph{Short moves}: as in
  Table~\ref{tab:al-b1}, but reversals that brake and reverse. \emph{Full speed}: one move per axis up to the cap and
  back to rest. Worst case over the counted moves of the 9 start points (moves within 0.25\,rad of a joint limit are
  not counted).}
  \label{tab:al-final}
  {\renewcommand{\arraystretch}{1.15}
  \begin{tabular}{rrrrrrrr}
    \toprule
    & & \multicolumn{3}{c}{short moves, real reversals} & \multicolumn{3}{c}{full speed} \\
    \cmidrule(lr){3-5} \cmidrule(lr){6-8}
    accel.\ limit & speed cap & lag & overshoot & tilt & lag & overshoot & moves counted \\
    \midrule
    0.10\,m/s$^2$ & 20\,cm/s & 2.15\,mm & 1.65\,mm & 0.35$^\circ$ & 3.35\,mm$^*$ & 3.01\,mm & 7 of 27 \\
    0.10\,m/s$^2$ & 10\,cm/s & 2.15\,mm & 1.65\,mm & 0.35$^\circ$ & 2.00\,mm & 1.93\,mm & 21 of 27 \\
    \textbf{0.08\,m/s}$^\mathbf{2}$ & \textbf{10\,cm/s} & \textbf{1.68\,mm} & \textbf{1.32\,mm} & \textbf{0.27}$^\circ$ &
      \textbf{1.65\,mm} & \textbf{1.58\,mm} & 21 of 27 \\
    \bottomrule
  \end{tabular}}

  \smallskip
  {\small Reading: at 0.1\,m/s$^2$ the real reversals at the far, high points ($x$ 0.60, $z$ 0.45\,m) lag 2.15\,mm,
  i.e.\ $\Lambda \approx 21.5$\,kg in Eq.~\eqref{eq:al-lag}, just above the criterion. At 20\,cm/s a full-speed move is
  about 42\,cm long and leaves the comfortable workspace: only 7 of 27 moves stay clear of the joint limits, and the
  counted ones lag up to 3.35\,mm ($^*$ start pose, $+y$, while braking far out; two further moves ran beyond the
  arm's reach). At 10\,cm/s all full-speed moves of the 7 counted start points fit. \textbf{0.08\,m/s$^2$ and
  10\,cm/s} (bold, chosen) meet the criterion everywhere with about 15\,\% margin. Rotations at the same limits: angle
  lag $\le 0.62^\circ$, tip drift $\le 1.23$\,mm (full-speed turns of about 120$^\circ$ about $x$ and $y$ reach the
  wrist's joint limits at most points, so only 1--2 of 3 turns per point are counted). Not counted: the two points
  near joint 4's limit lag up to 2.3\,mm.}
\end{table}

\clearpage
% ---------------------------------------------------------------------------------------------------------------------
\subsection{Design choices and their trade-offs}
\label{sec:al-choices}
These are the building blocks of the design. Every one moves the balance between \emph{transfer robustness}
(small lag, so that the motion is set by the controller), \emph{agility} (how fast the policy can move and react),
\emph{gentleness} (forces on the plant) and \emph{effort} (implementation, identification, training time). The
current value is given in the first column.

\paragraph{Controller (what turns the target into torques).}
\begin{center}
{\small\renewcommand{\arraystretch}{1.2}
\begin{tabular}{>{\raggedright\arraybackslash}p{0.16\textwidth}>{\raggedright\arraybackslash}p{0.23\textwidth}>{\raggedright\arraybackslash}p{0.245\textwidth}>{\raggedright\arraybackslash}p{0.245\textwidth}}
  \toprule
  Building block (current) & What it is & Gains & Costs \\
  \midrule
  Translational stiffness $K$ \newline (\textbf{1000\,N/m}, was 400)
    & spring between target and tool position. Franka's examples use 150--400\,N/m; Franka's internal Cartesian
      impedance accepts up to 3000\,N/m
    & lag $\propto 1/K$, so $a_\text{max} \propto K$ (Eq.~\ref{eq:al-amax}); less position error from joint friction on
      the real arm (error $\approx$ friction$/K$)
    & more force per mm of lag on contact; the damping grows with $\sqrt K$, so encoder noise and loop delay matter
      more; stability margin on the real robot to be checked \\
  \addlinespace\cmidrule(lr){1-4}\addlinespace
  Rotational stiffness $K_o$ \newline (\textbf{100\,N\,m/rad}, was 30)
    & spring between target and tool orientation (10--300\,N\,m/rad)
    & tilt $\propto 1/K_o$ while translating (Table~\ref{tab:al-diag}); smaller angle lag when turning
    & twists a stem caught in the fork more stiffly; same noise/stability caveat as $K$ \\
  \addlinespace\cmidrule(lr){1-4}\addlinespace
  Damping $D$ \newline (\textbf{$2\sqrt K\Lambda^{1/2}$})
    & damper between target and tool velocity
    & critical in every pose and direction; Franka's $2\sqrt K$ (critical only for 1\,kg) overshot by up to 7\,mm
    & uses the model of $\Lambda$ (only for the damping) \\
  \addlinespace\cmidrule(lr){1-4}\addlinespace
  Velocity feedforward $\dot x_d$ \newline (\textbf{on})
    & the damper acts on the velocity \emph{difference}, not on the tool velocity; while the clamp holds the target,
      $\dot x_d$ is the tool's velocity (the held target moves with the tool)
    & no lag at constant speed (without it: 82\,mm at 20\,cm/s, Table~\ref{tab:al-time}); no damping force on top of
      the clamp's bound
    & none found \\
  \addlinespace\cmidrule(lr){1-4}\addlinespace
  Inertia feedforward \newline $J^\top\Lambda\ddot x_d$ \newline (\textbf{off}; option B3)
    & the controller supplies the force needed to accelerate the arm from the model (operational-space
      control~\cite{al-khatib1987}) instead of waiting for the spring to stretch
    & lag drops to $\Delta\Lambda\, a / K$, set only by the \emph{model error} $\Delta\Lambda$ (20\,\% error: 5 times less
      lag); also removes the tilt and tip drift from the coupling
    & the motion now depends on the model of $\Lambda$ (libfranka on the robot); needs the motor-inertia question
      settled and the model checked (identification) \\
  \addlinespace\cmidrule(lr){1-4}\addlinespace
  Control period \newline (\textbf{2\,ms} sim, 1\,ms robot)
    & how often the law is evaluated
    & 1\,ms in sim matches the robot exactly
    & simulation twice as slow, so training twice as long \\
  \bottomrule
\end{tabular}}
\end{center}

\clearpage
\paragraph{Commands (what the policy outputs, and the limits on it).}
\begin{center}
{\small\renewcommand{\arraystretch}{1.2}
\begin{tabular}{>{\raggedright\arraybackslash}p{0.16\textwidth}>{\raggedright\arraybackslash}p{0.23\textwidth}>{\raggedright\arraybackslash}p{0.245\textwidth}>{\raggedright\arraybackslash}p{0.245\textwidth}}
  \toprule
  Building block (current) & What it is & Gains & Costs \\
  \midrule
  Speed cap \newline (\textbf{10\,cm/s, 45$^\circ$/s}; 20\,cm/s until 8 Oct.)
    & largest target step per policy step
    & the policy can approach and turn within one episode; 10\,cm/s is reached and left again within about 13\,cm,
      i.e.\ within the stem area
    & 20\,cm/s needed 42\,cm to reach and leave, more than the stem area, and lagged up to 3.35\,mm on such moves
      (Table~\ref{tab:al-final}) \\
  \addlinespace\cmidrule(lr){1-4}\addlinespace
  Acceleration limit \newline (\textbf{0.08\,m/s$^2$, 20$^\circ$/s$^2$}; 0.1\,m/s$^2$ until 8 Oct.)
    & largest change of the target step from one policy step to the next
    & \emph{the} knob for the lag (Eq.~\ref{eq:al-lag}); part of the policy-side code, so identical in sim and on the
      robot; the policy sees the applied step, so it can learn the limit
    & less agility: long time to speed and long braking distance (table in Section~\ref{sec:al-problem}) \\
  \addlinespace\cmidrule(lr){1-4}\addlinespace
  Lag criterion $e_\text{max}$ \newline (\textbf{2\,mm, 2$^\circ$})
    & how much lag we accept; sets the acceleration limit via Eq.~\eqref{eq:al-amax}
    & smaller: the motion depends less on the arm's dynamics (better transfer~\cite{al-aljalbout2024})
    & smaller $a_\text{max}$; a larger value (e.g.\ 5\,mm, option B2) gives 2.5 times more acceleration \\
  \addlinespace\cmidrule(lr){1-4}\addlinespace
  Low-pass filter on the commands \newline (\textbf{not used})
    & smooths the commands instead of limiting their change
    & smooth motion; common in sim-to-real
    & adds hidden state; Ref.~\cite{al-aljalbout2024} saw learning collapse with a limiter the policy could not see;
      no hard bound on the acceleration \\
  \addlinespace\cmidrule(lr){1-4}\addlinespace
  Target clamp \newline $\lVert x_d - x\rVert \le e_\text{clamp}$ \newline (\textbf{4\,mm, 3$^\circ$}, every control
    step)
    & the target is never further than $e_\text{clamp}$ from the measured tool; checked at every control step, not
      once per policy step (which would also cut steps in free motion)
    & hard bound on the spring force when the tool is blocked: 4\,mm $\times$ 1000\,N/m = 4\,N (about 3 times the
      measured push forces of 0.8--1.4\,N), 3$^\circ \times$ 100\,N\,m/rad = 5.2\,N\,m (43\,\% of the FR3's 12\,N\,m
      wrist torque limit; 10$^\circ$ would exceed it), even if the policy keeps pushing
    & must stay above the free-space lag (2\,mm), otherwise it acts while moving; while it acts, the target does not
      follow the commands (the policy sees this in the applied step); runs in the 1\,kHz loop on the robot \\
  \addlinespace\cmidrule(lr){1-4}\addlinespace
  Policy rate \newline (\textbf{31.25\,Hz})
    & how often the policy sends a new step
    & higher: finer commands, quicker reaction
    & more policy evaluations per second of simulation, so longer training; more sensitive to latency \\
  \addlinespace\cmidrule(lr){1-4}\addlinespace
  Action space \newline (\textbf{Cartesian steps on the previous target})
    & alternatives: joint velocities, Cartesian velocities, absolute poses, joint torques
    & joint velocities transferred best in Ref.~\cite{al-aljalbout2024} (stiff joint-level tracking)
    & joint velocities: no Cartesian compliance at contact, which matters for a fragile stem; torques: motion set by
      the dynamics alone, could not be transferred safely~\cite{al-aljalbout2024} \\
  \bottomrule
\end{tabular}}
\end{center}

\clearpage
\paragraph{Task and training.}
\begin{center}
{\small\renewcommand{\arraystretch}{1.2}
\begin{tabular}{>{\raggedright\arraybackslash}p{0.16\textwidth}>{\raggedright\arraybackslash}p{0.23\textwidth}>{\raggedright\arraybackslash}p{0.245\textwidth}>{\raggedright\arraybackslash}p{0.245\textwidth}}
  \toprule
  Building block (current) & What it is & Gains & Costs \\
  \midrule
  Episode length \newline (\textbf{10\,s})
    & how long the policy has per attempt
    & longer: a slow approach still leaves time to push the stem
    & training time grows about in proportion; the reward comes later after the actions that earned it, which makes
      learning harder \\
  \addlinespace\cmidrule(lr){1-4}\addlinespace
  Workspace bound \newline (\textbf{planned})
    & keep the tool away from poses near joint limits and the edge of the reach
    & there the lag grows and the real FR3 stops with a reflex
    & excludes some stem placements \\
  \addlinespace\cmidrule(lr){1-4}\addlinespace
  Reward terms \newline (\textbf{planned}: contact force, action rate, joint-limit margin)
    & penalties the policy learns to avoid, instead of hard limits
    & the policy learns gentle motion itself, where it matters
    & soft: no guarantee on the robot \\
  \addlinespace\cmidrule(lr){1-4}\addlinespace
  Domain randomization \newline (\textbf{planned}, after a first baseline run)
    & vary $K$, $K_o$, $\Lambda$ (motor inertia, link masses), latency and friction in training
    & policy robust to the remaining differences between sim and real~\cite{al-peng2018}
    & ranges must be realistic; too wide makes the policy cautious and slow to train \\
  \addlinespace\cmidrule(lr){1-4}\addlinespace
  Identification \newline (\textbf{planned}, before the first real run)
    & compare libfranka's model with the simulation; record the real arm's response to the same commands
    & tells how large the dynamics gap is; prerequisite for inertia feedforward and for loosening the criterion
    & robot time \\
  \addlinespace\cmidrule(lr){1-4}\addlinespace
  FR3 safety layer in sim \newline (\textbf{planned})
    & joint velocity, torque and torque-rate limits and the collision reflex as terminations/penalties
    & the policy does not learn motions the robot refuses
    & thresholds must match the robot's configuration \\
  \bottomrule
\end{tabular}}
\end{center}

\paragraph{Main options, in short.}
\begin{center}
{\renewcommand{\arraystretch}{1.2}
\begin{tabular}{>{\raggedright\arraybackslash}p{0.24\textwidth}>{\raggedright\arraybackslash}p{0.18\textwidth}>{\raggedright\arraybackslash}p{0.22\textwidth}>{\raggedright\arraybackslash}p{0.26\textwidth}}
  \toprule
  Option & $a_\text{max}$ at 2\,mm lag & Depends on the arm model & Main cost \\
  \midrule
  B1: $K$ 1000, $K_o$ 100, no inertia feedforward & 0.08\,m/s$^2$ (measured; 0.1 lags 2.15\,mm) & only through the
    damping & sluggish \\
  \addlinespace
  B2: B1 with a 5\,mm criterion & $\approx 0.25$\,m/s$^2$ & yes: 5\,mm of lag are set by the arm & against the transfer
    evidence \\
  \addlinespace
  B3: B1 + inertia feedforward & $\approx 0.5$\,m/s$^2$ or more (not run) & yes: model of $\Lambda$ & identification
    first \\
  \addlinespace
  Stiffer $K$ (3000\,N/m), no inertia feedforward & $\approx 0.3$\,m/s$^2$ & only through the damping & real-robot
    stability, contact forces \\
  \bottomrule
\end{tabular}}
\end{center}

% ---------------------------------------------------------------------------------------------------------------------
\subsection{Plan}
\label{sec:al-plan}
Decision (7 October 2026, revised 8 October): \textbf{B1 for the first training and transfer}, B3 (inertia
feedforward) once a policy trains successfully and the robot's $\Lambda$ has been checked. Values from
Table~\ref{tab:al-final}: $K = 1000$\,N/m, $K_o = 100$\,N\,m/rad, speed caps 10\,cm/s (3.2\,mm per step) and
45$^\circ$/s, acceleration limits 0.08\,m/s$^2$ (0.08\,mm per step) and 20$^\circ$/s$^2$ (0.02$^\circ$ per step),
target clamp 4\,mm and 3$^\circ$ at every control step. On 8 October the speed cap was halved and the acceleration
limit lowered from 0.1\,m/s$^2$, after moves that run up to the cap and reversals that really reverse showed lags of
up to 3.35\,mm and 2.15\,mm (Table~\ref{tab:al-final}).

\paragraph{Is this fast enough for the task?} With these limits (speed up, then brake, from rest to rest): the
farthest stem tip, 31\,cm from the start pose, is reached in 4.4\,s (1.25\,s to reach 10\,cm/s, 1.85\,s at
10\,cm/s, 1.25\,s braking); a 90$^\circ$ turn takes 4.3\,s (peak 42$^\circ$/s), and turning can happen while
approaching. That leaves about 5.5\,s of the 10\,s episode for pushing. Braking from 10\,cm/s takes 6\,cm and 1.25\,s.
We consider this sufficient for a first training. If the trained policy runs out of time, we lengthen the episode
first (cost: training time) and loosen the limits only after the identification.

% ---------------------------------------------------------------------------------------------------------------------
\subsection{Questions}
\label{sec:al-questions}
\begin{enumerate}
  \item Is the analysis right: is the lag of a Cartesian impedance controller on the FR3 set by its apparent mass
        (8--22\,kg at the tool in simulation, half of it from the motors' reflected inertia)?
  \item Does the FR3's internal joint torque controller reduce the apparent motor inertia (as on the DLR
        lightweight arms), and is that why libfranka's model leaves the motor inertia out? This decides how the
        simulation must model the motors, and which mass matrix the damping of our controller should use on the
        robot.
  \item Stiffness: is $K = 1000$\,N/m (and $K_o = 100$\,N\,m/rad) in a libfranka torque loop acceptable on our FR3
        next to a plant? Has the lab run similar gains, and with which velocity filtering?
  \item Inertia feedforward (option B3): how accurate is libfranka's model in your experience? Has the lab used
        operational-space control with acceleration feedforward on the FR3?
  \item Is an acceleration of about 0.08\,m/s$^2$ (1.25\,s to 10\,cm/s) a reasonable limit near a plant, or would
        you accept larger tracking errors for more agility?
  \item Which reflexes and thresholds (collision behaviour, joint velocity limits) does the lab configure, so that we
        model the same ones in simulation?
\end{enumerate}

{\raggedright
\begin{thebibliography}{9}
\bibitem{al-aljalbout2024} E.~Aljalbout, F.~Frank, M.~Karl, P.~van der Smagt. On the role of the action space in
  robot manipulation learning and sim-to-real transfer. \emph{IEEE Robotics and Automation Letters}, 2024.
  arXiv:2312.03673.
\bibitem{al-khatib1987} O.~Khatib. A unified approach for motion and force control of robot manipulators: the
  operational space formulation. \emph{IEEE Journal on Robotics and Automation}, 3(1), 1987.
\bibitem{al-peng2018} X.~B.~Peng, M.~Andrychowicz, W.~Zaremba, P.~Abbeel. Sim-to-real transfer of robotic control
  with dynamics randomization. \emph{IEEE ICRA}, 2018.
\end{thebibliography}
}
% ======================================== END CONTENT ========================================
) or compiled alone via action_limits_report.tex.
% Uses amsmath, booktabs and array (added to main.tex for this section). Numbers: scripts/sweep_action_poses.py
% (default and --full_speed), scripts/check_push_env.py and docs/notes/2026-10-07.md, 2026-10-08.md in the project repo.

% ======================================= BEGIN CONTENT =======================================
\section{Open problem: how fast may the policy move the fork?}
\label{sec:al}

\paragraph{Status (7--8 October 2026).} Problem identified in simulation; values for the first training runs set
on 7 October (option B1) and revised on 8 October after longer test moves (Section~\ref{sec:al-plan}); the
long-term solution is open. We ask for feedback on the analysis and on the
questions at the end (Section~\ref{sec:al-questions}).

\paragraph{In one paragraph.} The policy commands the fork through a Cartesian impedance controller, the same in
simulation and on the FR3. What blocks us is a trade-off between \emph{agility} and a \emph{small tracking error}.
A small error transfers better~\cite{al-aljalbout2024}, because then the fork's motion is set by our controller
(identical on both sides) rather than by the arm's dynamics (not identical, not yet identified). But an impedance
controller is a spring of stiffness $K$ between the commanded and the actual tool position, and it can only
accelerate the arm by being stretched: while the commanded motion accelerates at $a$, the tool lags behind by about
$\Lambda a/K$, where $\Lambda$ is the arm's apparent mass at the tool tip (8--22\,kg depending on pose and direction,
half of it from the motors' reflected inertia). With a small lag (2\,mm) and a stiffness we consider acceptable near a
plant (400--1000\,N/m), the fork may accelerate only at about 0.04--0.09\,m/s$^2$; we use 0.08\,m/s$^2$ at
1000\,N/m. This section explains the
trade-off, what we measured, and every design choice that moves it.

% ---------------------------------------------------------------------------------------------------------------------
\subsection{Background}
\label{sec:al-background}

\paragraph{What the policy does.} The policy is a neural network trained by reinforcement learning (RL) in
simulation. At 31.25\,Hz (every 32\,ms, one ``policy step'') it reads the state (fork pose, stem shape, target) and
outputs a small displacement and rotation of the fork's \emph{target} pose $x_d$. The next target is the previous
target plus this step, so commanded steps are not lost even if the tool lags behind. Within the 32\,ms the target
moves at constant velocity $\dot x_d$ from the old to the new value. A Cartesian impedance controller at 1\,kHz
(2\,ms in simulation) turns the target into joint torques:
\begin{equation}
  \tau = J^\top\!\left[\,K\,(x_d - x) + D\,(\dot x_d - \dot x)\,\right] + \tau_\text{null} + \tau_\text{Coriolis}.
  \label{eq:al-law}
\end{equation}
$x$ is the measured tool pose, $J$ the tool-tip Jacobian, $K$ the stiffness and $D$ the damping (written here for
the position; the orientation has its own stiffness $K_o$). This is the law of Franka's example controllers, with
two changes: the target velocity $\dot x_d$ is fed forward (Franka's examples use $\dot x_d = 0$), and the damping
$D = 2\sqrt{K}\,\Lambda^{1/2}$ is chosen critical for the arm's actual inertia $\Lambda$ (Section~\ref{sec:al-problem})
instead of Franka's $2\sqrt K$, which is critical only for a 1\,kg tool.

\paragraph{Null-space spring $\tau_\text{null}$.} The FR3 has 7 joints, but the tool pose has only 6 degrees of
freedom. For every tool pose there is therefore a one-parameter family of arm postures (the elbow can swing about the
line from shoulder to wrist without moving the tool): the \emph{null space}. The first term of
Eq.~\eqref{eq:al-law} does not act in this direction, so the elbow would drift freely. A weak joint-space spring
(0.5\,N\,m/rad, as in Franka's examples) pulls the posture towards a rest posture; it is projected so that it does
not push the tool. It plays no role in the problem below (we checked: switching it off does not change the lag).

\paragraph{What RL learns, and why the action space matters for the transfer.} During training the policy tries
millions of commands in simulation and keeps what earns reward. It learns \emph{how the simulated fork reacts to its
commands}, including any lag, overshoot or coupling, and relies on it. If the real robot reacts differently, the
learned behaviour is wrong by that difference: this is the sim-to-real (``reality'') gap. The policy cannot correct
what it never saw in simulation. The usual tools against the gap are:
\begin{center}
{\renewcommand{\arraystretch}{1.15}
\begin{tabular}{>{\raggedright\arraybackslash}p{0.24\textwidth}>{\raggedright\arraybackslash}p{0.42\textwidth}>{\raggedright\arraybackslash}p{0.24\textwidth}}
  \toprule
  Tool & What it does & Status \\
  \midrule
  A. Action-space design & choose commands whose result depends little on what differs between sim and real
    (same controller, small tracking error) & controller matched; this section \\
  \addlinespace
  B. Identification & measure the real robot and correct the simulation & planned, before the first real run \\
  \addlinespace
  C. Domain randomization & vary unknown parameters in training, so the policy works over a range & planned \\
  \addlinespace
  D. Model the robot's safety layer & limits and reflexes of the FR3 become terminations/penalties in sim, so the
    policy does not learn motions the robot refuses & planned \\
  \addlinespace
  E. Adaptation & infer differences online or fine-tune on the robot & not planned \\
  \bottomrule
\end{tabular}}
\end{center}
A large study of action spaces on a Franka arm~\cite{al-aljalbout2024} (13 action spaces, about 250 policies,
reaching and pushing, transferred to the real robot) concludes that \emph{``an action space which yields or
naturally limits the tracking error transfers better''}; torque actions, whose result is set by the robot's
dynamics alone, could not be transferred safely. The same study warns that a limiter which acts on the robot but is
not visible to the policy in simulation (they tried a rate limiter) made learning collapse. A caveat: their controllers were joint
impedance controllers (Cartesian actions were converted to joint targets by inverse kinematics), with gains that
the paper does not report, and the tasks were reaching and pushing a box. The conclusion carries over in principle,
since in every impedance-type controller a larger tracking error means more dependence on the arm's dynamics; the
numbers do not, and the paper does not test a soft Cartesian impedance controller like ours.

\paragraph{Where the gap enters our system.} The chain from a policy output to the stem, step by step:
\begin{center}
{\renewcommand{\arraystretch}{1.15}
\begin{tabular}{>{\raggedright\arraybackslash}p{0.32\textwidth}>{\raggedright\arraybackslash}p{0.20\textwidth}>{\raggedright\arraybackslash}p{0.38\textwidth}}
  \toprule
  Step & Same in sim and real? & Source of differences \\
  \midrule
  policy step $\to$ target pose & yes (same code) & none \\
  target $\to$ torques, law~\eqref{eq:al-law} & yes (ported to C++, checked) & port errors; $J$, $M$ from different
    models \\
  \textbf{torques $\to$ fork motion} & \textbf{no} & arm inertia, joint friction, motor inertia, delays, 1 vs.\ 2\,ms
    control period \\
  fork $\leftrightarrow$ stem contact & no & friction, stem stiffness and damping (randomized) \\
  stem $\to$ observation & no & camera noise, latency \\
  FR3 safety layer & \textbf{missing in sim} & reflexes stop the robot \\
  \bottomrule
\end{tabular}}
\end{center}
The step ``torques $\to$ fork motion'' is where the arm's dynamics act. How much of it reaches the fork's motion
depends on the tracking error: if the tool follows the target within a small error, its motion is essentially the
target's motion, which our code computes identically in sim and real; the larger the error, the more the motion is
shaped by the arm's dynamics, which differ.

% ---------------------------------------------------------------------------------------------------------------------
\subsection{The problem}
\label{sec:al-problem}

\paragraph{What blocks us: agility versus a small tracking error.} We want two things from the fork's motion:
\begin{enumerate}
  \item \textbf{Small lag behind the commanded target}, so that the fork's motion is set by our controller, which
        is the same in simulation and on the FR3, and not by the arm's dynamics, which differ and are not yet
        identified. Action spaces that keep the tracking error small transfer better~\cite{al-aljalbout2024}.
  \item \textbf{Agility}: reach the stem well within the 10\,s episode, brake late before contact, and react quickly
        to how the stem moves.
\end{enumerate}
With an impedance controller the two conflict. The lag grows with how quickly the commanded motion accelerates, so
a small lag allows only gentle accelerations. In numbers: with $K = 1000$\,N/m and at most 2\,mm of lag, the fork may
accelerate at only about 0.09\,m/s$^2$ where the arm is heaviest; at the 0.08\,m/s$^2$ we use, it needs 1.25\,s to
reach 10\,cm/s and 6\,cm to stop from that speed (2.5\,s and 25\,cm for 20\,cm/s). Either
we accept this sluggishness, or we accept more lag (and with it a larger dependence on the real arm), or we change
the controller. This is the decision we need help with. The rest of this section explains where the conflict comes
from, how severe it is for the task, and how severe it will be on the real robot.

\paragraph{Why the tool lags: the impedance law as a mass on a spring.} Consider one direction, free space,
gravity compensated. The controller makes the tool tip behave like a mass $\Lambda$ (the arm's inertia felt at the
tool tip, explained in the next paragraph) attached by a spring $K$ and a damper $D$ to the moving target:
\begin{equation}
  \Lambda\,\ddot x = K\,(x_d - x) + D\,(\dot x_d - \dot x).
\end{equation}
Write the lag as $e = x_d - x$ (target minus tool). Then $\ddot x = \ddot x_d - \ddot e$, and substituting gives
\begin{equation}
  \Lambda\,\ddot e + D\,\dot e + K\,e = \Lambda\,\ddot x_d.
  \label{eq:al-error}
\end{equation}
The lag itself behaves like a damped mass-spring system, driven by the force $\Lambda \ddot x_d$, i.e.\ by the
\emph{acceleration of the target}. Three cases follow:
\begin{enumerate}
  \item \textbf{Target at rest or moving at constant speed} ($\ddot x_d = 0$): the right-hand side is zero, so any lag
        dies out and $e \to 0$. Physically: moving at constant speed in free space needs no force (no friction, no
        gravity), so the spring does not need to be stretched. This holds only because the target velocity is fed
        forward. Without it the damper pulls back with $D\dot x$ all the time, and the spring must balance it:
        $e = D\dot x_d / K$. With $D = 2\sqrt{K\Lambda} = 180$\,N\,s/m ($K = 400$\,N/m, $\Lambda = 20$\,kg) and
        20\,cm/s this gives 90\,mm; we measured 82\,mm (Table~\ref{tab:al-time}).
  \item \textbf{Target accelerating at constant $a$}: the right-hand side is the constant force $\Lambda a$. Once the
        transient has died out ($\ddot e = \dot e = 0$), $K e = \Lambda a$, so
        \begin{equation}
          e = \frac{\Lambda\, a}{K}.
          \label{eq:al-lag}
        \end{equation}
        Physically: Newton's law. To accelerate the mass $\Lambda$ at $a$, the force $\Lambda a$ is needed, and its only
        source is the stretched spring. Example: $\Lambda = 20$\,kg, $a = 0.1$\,m/s$^2$, $K = 1000$\,N/m: $e = 2$\,mm.
        The same holds when braking (the tool then runs ahead of the target, which shows up as overshoot).
  \item \textbf{How fast the lag builds up and dies out}: Eq.~\eqref{eq:al-error} is critically damped, with the time
        constant $\sqrt{\Lambda/K}$: 0.14\,s for 20\,kg and 1000\,N/m, 0.22\,s for 400\,N/m, i.e.\ 4--7 policy steps.
        The lag reaches the value of Eq.~\eqref{eq:al-lag} if the acceleration lasts longer than about three time
        constants, and it disappears again within a similar time once the target moves at constant speed. This
        matters for testing: short test moves end before the lag has settled. Our first sweep (20 policy steps per
        segment, reversals that only braked) therefore showed 1.9\,mm at 0.1\,m/s$^2$; moves that brake \emph{and}
        reverse, or that run up to the speed cap, accelerate long enough and reached 2.15\,mm
        (Table~\ref{tab:al-final}).
\end{enumerate}

\paragraph{The mass in this equation: the apparent mass $\Lambda$ (``8--20\,kg at the tool tip'').} The trade-off
is severe because the mass that the spring has to accelerate is large. If one pushes the tool tip of
the arm (joints free, gravity compensated) with a force $F$, the tip starts to accelerate with $\ddot x = \Lambda^{-1}
F$. $\Lambda$ is the arm's inertia as seen at the tool tip, the \emph{apparent mass} (or operational-space inertia):
\begin{equation}
  \Lambda = \left(J M^{-1} J^\top\right)^{-1},
  \label{eq:al-lambda}
\end{equation}
with $M$ the $7\times7$ joint-space mass matrix and $J$ the $6\times7$ tool-tip Jacobian. $\Lambda$ is a $6\times6$
matrix (force and moment versus linear and angular acceleration) and depends on the pose. Its translational
diagonal entries say how heavy the arm feels when the tip is pushed along $x$, $y$ or $z$. At the earlier start pose
(2026-10-07; tool tip at $(0.30, 0, 0.55)$\,m in front of the robot, fork horizontal), computed from the simulated FR3:
\begin{center}
{\renewcommand{\arraystretch}{1.15}
\begin{tabular}{lrrr}
  \toprule
  Apparent mass, start pose & along $x$ (forward) & along $y$ (sideways) & along $z$ (up) \\
  \midrule
  links only & 11.9\,kg & 4.2\,kg & 5.1\,kg \\
  links + motors' reflected inertia (our simulation) & 20.5\,kg & 7.6\,kg & 10.6\,kg \\
  \bottomrule
\end{tabular}}
\end{center}
Pushing the tip forward with 20\,N accelerates it at about 1\,m/s$^2$, as if it were a 20\,kg mass; the same push
sideways gives 2.6\,m/s$^2$. Forward motion needs the heavy shoulder and elbow joints, sideways motion mostly the
base rotation and the wrist. The second row includes the motors' reflected inertia $n^2 I_\text{motor}$ ($n$: gear
ratio), which we add to $M$ from Franka's description files: 0.61, 0.61, 0.46, 0.46, 0.21, 0.21, 0.21\,kg\,m$^2$ for
joints 1--7 (at the wrist about 100 times the link's own inertia). Without these terms the simulated arm oscillated
and pressed the stem 5 times harder. The lag measurements below give the same value independently: from
Eq.~\eqref{eq:al-lag}, $\Lambda = K e / a \approx 16$--20\,kg along $x$ at the start pose, and up to 21.5\,kg
at the far, high start points of the sweep ($x$ 0.60, $z$ 0.45\,m).

\paragraph{How severe the trade-off is for the task.} Turning Eq.~\eqref{eq:al-lag} around: for a largest acceptable lag $e_\text{max}$, the
target may accelerate at most at
\begin{equation}
  a_\text{max} \approx \frac{K\, e_\text{max}}{\Lambda}.
  \label{eq:al-amax}
\end{equation}
With $\Lambda = 21.5$\,kg (the heaviest pose and direction) and $e_\text{max} = 2$\,mm: 0.04\,m/s$^2$ at
$K = 400$\,N/m and 0.09\,m/s$^2$ at 1000\,N/m. We use 0.08\,m/s$^2$, which leaves about 15\,\% margin (lag
$\le 1.72$\,mm; measured 1.68\,mm, Table~\ref{tab:al-final}). What this means for the task (the
target speeds up at $a_\text{max}$ until the speed cap, then brakes at $a_\text{max}$):
\begin{center}
{\renewcommand{\arraystretch}{1.15}
\begin{tabular}{rrrrr}
  \toprule
  largest acceleration & time from rest & braking distance & 31\,cm, rest to rest & 31\,cm, rest to rest \\
  $a_\text{max}$ & to 20\,cm/s & from 20\,cm/s & (20\,cm/s cap) & (10\,cm/s cap) \\
  \midrule
  0.05\,m/s$^2$ & 4.0\,s & 40\,cm & 5.0\,s & 5.1\,s \\
  \textbf{0.08\,m/s}$^\mathbf{2}$ & 2.5\,s & 25\,cm & 3.9\,s & \textbf{4.4\,s} \\
  0.10\,m/s$^2$ & 2.0\,s & 20\,cm & 3.5\,s & 4.1\,s \\
  0.25\,m/s$^2$ & 0.8\,s & 8\,cm  & 2.3\,s & 3.5\,s \\
  0.50\,m/s$^2$ & 0.4\,s & 4\,cm  & 1.9\,s & 3.3\,s \\
  \bottomrule
\end{tabular}}
\end{center}
31\,cm is the distance from the start pose to the farthest stem tip. The episode lasts 10\,s, and the stem must
still be pushed after the approach. The cost of a small $a_\text{max}$ is therefore not the top speed but agility:
the policy must start braking long before it touches the stem and cannot react quickly to how the stem moves.

\paragraph{Why not simply make the episode longer?} A longer episode gives a slow policy time to finish, but it is
not free. RL learns from simulated experience, so the training time grows roughly in proportion to the simulated
time per episode: a 20\,s episode costs about twice the training time of a 10\,s episode for the same number of
attempts. In addition, the reward for reaching the target arrives later after the actions that caused it, which
makes it harder for the policy to learn which actions were good. We therefore keep 10\,s as long as the task fits,
and lengthen it only if the trained policy demonstrably runs out of time.

\paragraph{Rotations: the same effect, plus a coupling.} Turning the fork works in the same way: the orientation
spring $K_o$ must twist to give the rotational inertia an angular acceleration $\alpha$, so the angle lags by about
$\Lambda_o \alpha / K_o$ ($\Lambda_o$: the rotational part of $\Lambda$). In addition, $\Lambda$ is a full $6\times6$
matrix: its off-diagonal blocks couple translation and rotation. The reason is geometric: the tool tip is 10--20\,cm
away from the wrist's joint axes, and the joints that turn the fork also move the tip, while those that move the
tip also turn the fork. The control law, however, treats position and orientation as two independent springs.
Two consequences, both measured:
\begin{itemize}
  \item \textbf{Tilt while translating.} Accelerating the tip at $a$ needs, besides the force, a moment
        $\Lambda_{op}\, a$ ($\Lambda_{op}$: off-diagonal block). It comes from the orientation spring, which must twist:
        the fork tilts by about $\Lambda_{op} a / K_o$. Table~\ref{tab:al-diag} shows the tilt falling with $1/K_o$
        (2.0$^\circ$ at $K_o = 30$, 0.26$^\circ$ at 300\,N\,m/rad).
  \item \textbf{Tip drift while turning.} Turning the fork at angular acceleration $\alpha$ needs, besides the moment, a
        force at the tip, which the position spring supplies by stretching: the tip moves away from its target by
        about $\Lambda_{po}\,\alpha / K$, although the position command is constant (Tables~\ref{tab:al-rot}
        and~\ref{tab:al-b1}).
\end{itemize}
Both effects scale with the acceleration, like the lag. A controller with inertia feedforward (operational-space
control, option B3 below) would supply these coupling forces and moments from the model and remove both.

\paragraph{Contact.} When the tool is blocked (e.g.\ pressed against the stem), the target keeps moving and the
lag becomes force: $K e$. Small lags therefore also mean small unintended forces. A hard bound is added by
clamping the target: at every control step, $\lVert x_d - x \rVert \le 4$\,mm (spring force $\le 4$\,N at
$K = 1000$\,N/m) and the angle between target and tool orientation $\le 3^\circ$ (moment $\le 5.2$\,N\,m at
$K_o = 100$\,N\,m/rad). In free motion the clamp is inactive, because the lag stays below 2\,mm. While it holds the
target, the target moves with the tool, so the target velocity fed forward is the tool's own velocity: the damper
then adds no force to the bounded spring force, and the controller does not switch back and forth between holding and
releasing. Not bounded by the clamp: the momentum of the arm when the moving fork first hits something (as with any
impedance controller); the policy is meant to learn to approach gently (contact-force and smoothness penalties).

\paragraph{How severe the trade-off will be on the real robot: which $\Lambda$ does the FR3 have?} All numbers above
use the simulated $\Lambda$. The real arm's $\Lambda$ is not known yet, and two observations from Franka's software
suggest that it may differ:
\begin{itemize}
  \item libfranka (current version, September 2026) computes its model (\texttt{franka::Model::mass}) with the
        Pinocchio library from the URDF that the robot provides. The FR3 URDF shipped with libfranka's tests has
        the same link inertias as Franka's description files but \emph{no} motor-inertia entries, and libfranka never
        sets Pinocchio's rotor-inertia (``armature'') term. So libfranka's mass matrix most likely \emph{excludes}
        the motors' reflected inertia. (Read in the source code, not yet checked on our robot.)
  \item The FR3 executes commanded torques with an internal joint torque controller that uses the link-side
        torque sensors. On the DLR lightweight arms, from which Franka's design derives, such torque feedback
        reduces the apparent motor inertia. If the FR3 does the same, the real arm feels lighter than our
        simulation (between the two rows of the table above), and libfranka's model would be the right one.
\end{itemize}
How much this matters depends on the option chosen (Section~\ref{sec:al-choices}):
\begin{center}
{\renewcommand{\arraystretch}{1.15}
\begin{tabular}{>{\raggedright\arraybackslash}p{0.24\textwidth}>{\raggedright\arraybackslash}p{0.66\textwidth}}
  \toprule
  Option & Role of the real $\Lambda$ \\
  \midrule
  B1 (small lag, now) & \textbf{Little}, and this is exactly why small lags transfer well: the unknown $\Lambda$
    acts on the motion only through the lag. The acceleration limit was sized for $\le 2$\,mm with the simulated
    $\le 21.5$\,kg; a real arm of 12\,kg would lag 1\,mm, a heavier one (30\,kg) 2.4\,mm. The policy is then wrong by about a
    millimetre. \textbf{One exception: the damping} $D = 2\sqrt K \Lambda^{1/2}$ is computed from a model of $\Lambda$.
    Before the first real run we must decide which mass matrix the controller on the robot uses (libfranka's, or
    libfranka's plus the motor terms); with the wrong one the tool is under- or over-damped. \\
  \addlinespace
  B2 (larger lag) & \textbf{More}: with 5\,mm of lag set by the arm, a 40\,\% error in $\Lambda$ shifts the motion by
    2\,mm. \\
  \addlinespace
  B3 (inertia feedforward, later) & \textbf{Central}: the controller supplies the accelerating force from the model
    of $\Lambda$, so the remaining lag is set by the model error. Requires the identification first. \\
  \bottomrule
\end{tabular}}
\end{center}

% ---------------------------------------------------------------------------------------------------------------------
\subsection{Measurements}
\label{sec:al-measurements}

\paragraph{Setup.} Simulation (Isaac Lab, PhysX) of the FR3 with the fork, gravity compensated, no joint friction,
control period 2\,ms, policy step 32\,ms, no stem (free space). Every test commands the fork with the same commands
a policy could send, through the full controller.

\paragraph{Commands and limits.} A command is a target step per policy step. Three limits act on it:
\begin{itemize}
  \item \emph{speed cap}: the largest step, e.g.\ 3.2\,mm per 32\,ms = 10\,cm/s (rotation: 1.44$^\circ$ = 45$^\circ$/s);
  \item \emph{acceleration limit}: the largest change of the step from one policy step to the next. A change of
        0.1\,mm per step means $a = 0.1\,\text{mm}/(32\,\text{ms})^2 \approx 0.1$\,m/s$^2$, 0.08\,mm per step
        0.08\,m/s$^2$ (rotation: 0.02$^\circ$ per step $\approx$ 20$^\circ$/s$^2$). Commanding ``full speed'' from rest therefore makes the target speed up at
        this acceleration until the speed cap; commanding ``stop'' makes it brake at the same rate;
  \item the first design (Table~\ref{tab:al-first}) had a speed cap only, so the target jumped from rest to full
        speed in one step.
\end{itemize}

\paragraph{Test moves.} \emph{Straight move}: full command along one axis of the robot base frame ($\pm x$, $\pm y$,
$\pm z$) for a number of policy steps, then zero command (``hold'') for 45 steps (1.4\,s), during which the target
brakes and stops. \emph{Reversal}: full command along $+x$, then immediately along $-x$ (same for $y$, $z$), then hold;
this contains the largest change of velocity. Since 8 October the second segment is twice as long, so that the
target brakes \emph{and} reverses (before, it only braked to rest). \emph{Full-speed move} (since 8 October): from
rest up to the speed cap, 10 steps at the cap, then braking to rest (about 16\,cm; rotations about 120$^\circ$),
one per axis, in a direction with room for it (towards the middle of the workspace; up along $z$).
\emph{Rotation moves}: the same about the base axes, with zero position command. \emph{Start points}: either the
start pose only, or 9 points spanning the workspace in which the stems stand (tool tip at $x$ = 0.35 and 0.60\,m,
$y = \pm0.20$\,m, $z$ = 0.20 and 0.45\,m, plus the start pose). Moves that come within 0.25\,rad of a joint limit
are reported but not counted (decision 8 October): there the arm, not the controller, limits the motion. This
concerns the two start points low and close to the robot ($x$ 0.35, $z$ 0.20\,m), where joint 4 is within
0.21\,rad of its limit.

\paragraph{What is measured (all tables).}
\begin{itemize}
  \item \emph{Lag} (following error): the largest distance between the target position and the actual tool-tip
        position during the move and the hold. The quantity of Eq.~\eqref{eq:al-lag}.
  \item \emph{Overshoot}: how far the tool tip goes past the target's final position after it stops.
  \item \emph{Tilt}: the largest change of the fork's orientation during a straight move, although no rotation is
        commanded.
  \item \emph{Angle lag}: during rotation moves, the largest angle between the target and the actual orientation.
  \item \emph{Tip drift}: during rotation moves, the largest displacement of the tool tip, although no translation is
        commanded.
\end{itemize}
Each table reports the worst value over all moves (and start points). The criterion used so far: lag $\le 2$\,mm,
overshoot $\le 5$\,mm, tilt $\le 1^\circ$; angle lag $\le 2^\circ$, tip drift $\le 2$\,mm.

\begin{table}[h]
  \centering
  \caption{\textbf{First design}: speed cap only, no acceleration limit, $K = 400$\,N/m, $K_o = 30$\,N\,m/rad. Straight
  moves of 30 policy steps at full command from 9 start points, so each move starts with a jump from rest to full speed.
  Lag measured over the last quarter of each move (as defined at the time), overshoot after it. Worst case over all
  moves and start points. This design kept the lag below 2\,mm by limiting the speed to 5.5\,cm/s.}
  \label{tab:al-first}
  {\renewcommand{\arraystretch}{1.15}
  \begin{tabular}{rrrr}
    \toprule
    speed cap (step per 32\,ms) & speed & lag & overshoot \\
    \midrule
    2.0\,mm  & 6.3\,cm/s & 2.25\,mm & 3.9\,mm \\
    1.75\,mm & 5.5\,cm/s & 1.94\,mm & 3.4\,mm \\
    1.5\,mm  & 4.7\,cm/s & 1.63\,mm & 2.9\,mm \\
    \bottomrule
  \end{tabular}}

  \smallskip
  {\small Reading: with a cap of 5.5\,cm/s and 7.8$^\circ$/s the approach to a far stem took more than half of the
  episode, and a 90$^\circ$ turn did not fit in one. We therefore raised the caps to 20\,cm/s and 45$^\circ$/s and
  added the acceleration limit; Tables~\ref{tab:al-time}--\ref{tab:al-b1} show what that does.}
\end{table}

\begin{table}[h]
  \centering
  \caption{\textbf{How the lag evolves over time} (illustrates Eq.~\eqref{eq:al-error}). Start pose, $K = 400$\,N/m,
  $K_o = 30$\,N\,m/rad. One straight move along $+x$ at full command for 40 policy steps. With the acceleration limit
  of 0.49\,m/s$^2$ the target speeds up during the first 13 steps (0.42\,s) and then moves at constant speed.
  Lag of the tool tip behind the target after 10, 20, 30 and 40 policy steps.}
  \label{tab:al-time}
  {\renewcommand{\arraystretch}{1.15}
  \begin{tabular}{lrrrr}
    \toprule
    & \multicolumn{4}{c}{lag after $n$ policy steps} \\
    \cmidrule(lr){2-5}
    & 10 (0.32\,s) & 20 (0.64\,s) & 30 (0.96\,s) & 40 (1.28\,s) \\
    & accelerating & constant speed & constant speed & constant speed \\
    \midrule
    20\,cm/s, velocity feedforward on  & 9.4\,mm & 10.4\,mm & 4.7\,mm & 1.9\,mm \\
    10\,cm/s, velocity feedforward on  & 7.7\,mm & 4.7\,mm  & 2.0\,mm & 0.7\,mm \\
    20\,cm/s, velocity feedforward off & 23\,mm  & 59\,mm   & 75\,mm  & 82\,mm \\
    \bottomrule
  \end{tabular}}

  \smallskip
  {\small With feedforward the lag builds up while the target accelerates and dies out at constant speed (case 1 and 2
  of Section~\ref{sec:al-problem}); without it, the lag grows towards $D\dot x_d/K$. Switching off the null-space
  spring gave the same numbers.}
\end{table}

\begin{table}[h]
  \centering
  \caption{\textbf{Effect of the stiffnesses $K$ and $K_o$.} Start pose, speed cap 20\,cm/s. A straight move along $+x$
  at full command for 25 policy steps (0.8\,s, enough to reach about the speed cap), hold for 45 steps, then the same
  along $-x$ and hold. Two acceleration limits (columns). Largest lag of the tool tip and largest tilt of the fork
  over the whole test.}
  \label{tab:al-diag}
  {\renewcommand{\arraystretch}{1.15}
  \begin{tabular}{rrrrrr}
    \toprule
    & & \multicolumn{2}{c}{acceleration 0.24\,m/s$^2$} & \multicolumn{2}{c}{acceleration 0.49\,m/s$^2$} \\
    \cmidrule(lr){3-4} \cmidrule(lr){5-6}
    $K$ [N/m] & $K_o$ [N\,m/rad] & lag & tilt & lag & tilt \\
    \midrule
    400  & 30  & 9.7\,mm  & 1.80$^\circ$ & 13.6\,mm & 2.40$^\circ$ \\
    400  & 100 & 10.5\,mm & 0.66$^\circ$ & 14.4\,mm & 0.84$^\circ$ \\
    400  & 300 & 10.6\,mm & 0.23$^\circ$ & 14.8\,mm & 0.29$^\circ$ \\
    1000 & 30  & 4.0\,mm  & 2.03$^\circ$ & 6.4\,mm  & 2.69$^\circ$ \\
    1000 & 100 & 4.6\,mm  & 0.73$^\circ$ & 7.1\,mm  & 1.00$^\circ$ \\
    1000 & 300 & 4.8\,mm  & 0.26$^\circ$ & 7.5\,mm  & 0.37$^\circ$ \\
    \bottomrule
  \end{tabular}}

  \smallskip
  {\small Reading: the lag depends on $K$ (2.5 times stiffer, about 2.4 times less lag) and on the acceleration, hardly
  on $K_o$; the tilt depends on $K_o$ (10 times stiffer, about 8 times less tilt). Every entry fails the 2\,mm criterion.}
\end{table}

\begin{table}[h]
  \centering
  \caption{\textbf{Workspace sweep at $K = 1000$\,N/m, $K_o = 30$\,N\,m/rad}, caps 20\,cm/s and 45$^\circ$/s. At each
  start point: straight moves along $\pm x, \pm y, \pm z$ and reversals in $x, y, z$ (30 policy steps per segment, then
  hold), then the same as rotations about the base axes. Each row is one pair of acceleration limits (translation,
  rotation). Worst case over all moves and over the five start points closest to the robot: from the start points at
  $x = 0.60$\,m these moves of up to 20\,cm ran to the edge of the arm's reach, which is a flaw of the test, not of the
  controller.}
  \label{tab:al-rot}
  {\renewcommand{\arraystretch}{1.15}
  \begin{tabular}{rrrrrr}
    \toprule
    \multicolumn{3}{c}{translation moves} & \multicolumn{3}{c}{rotation moves} \\
    \cmidrule(lr){1-3} \cmidrule(lr){4-6}
    acceleration limit & lag & tilt & angular acceleration limit & angle lag & tip drift \\
    \midrule
    0.24\,m/s$^2$ & 5.0\,mm  & 3.0$^\circ$ & 20$^\circ$/s$^2$  & 9.9$^\circ$$^*$ & 14.0\,mm$^*$ \\
    0.49\,m/s$^2$ & 8.5\,mm  & 5.0$^\circ$ & 49$^\circ$/s$^2$  & 4.7$^\circ$ & 4.0\,mm \\
    0.73\,m/s$^2$ & 11.0\,mm & 5.7$^\circ$ & 98$^\circ$/s$^2$  & 6.8$^\circ$ & 7.0\,mm \\
    0.98\,m/s$^2$ & 12.6\,mm & 6.1$^\circ$ & 195$^\circ$/s$^2$ & 9.4$^\circ$ & 14.3\,mm \\
    \bottomrule
  \end{tabular}}

  \smallskip
  {\small Reading: all rows fail the criterion; the lag grows with the acceleration as expected, the tilt is large
  because $K_o$ is soft. $^*$ one start point ($x$ 0.35, $y$ 0.20, $z$ 0.45\,m), reversal about $x$; at the other four
  points at most 4.1$^\circ$ / 3.7\,mm. Not yet understood (likely a wrist joint near its limit).}
\end{table}

\begin{table}[h]
  \centering
  \caption{\textbf{Workspace sweep for the interim choice of 7 October: $K = 1000$\,N/m, $K_o = 100$\,N\,m/rad}, caps
  20\,cm/s and 45$^\circ$/s, smaller acceleration limits. Same moves as Table~\ref{tab:al-rot} but 20 policy steps per
  segment (moves of at most about 6\,cm, so all 9 start points are valid). Worst case over all moves and all 9 start
  points. \emph{Superseded by Table~\ref{tab:al-final}}: these moves reached at most 2\,mm per step (6.25\,cm/s), not
  the caps, and the reversals only braked, so the lag had not settled.}
  \label{tab:al-b1}
  {\renewcommand{\arraystretch}{1.15}
  \begin{tabular}{rrrrrrr}
    \toprule
    \multicolumn{4}{c}{translation moves} & \multicolumn{3}{c}{rotation moves} \\
    \cmidrule(lr){1-4} \cmidrule(lr){5-7}
    acceleration limit & lag & overshoot & tilt & angular accel.\ limit & angle lag & tip drift \\
    \midrule
    0.05\,m/s$^2$ & 1.3\,mm & 0.9\,mm & 0.17$^\circ$ & 10$^\circ$/s$^2$ & 0.30$^\circ$ & 0.55\,mm \\
    \textbf{0.10\,m/s}$^\mathbf{2}$ & \textbf{1.9\,mm} (2.4$^*$) & \textbf{1.8\,mm} & \textbf{0.33}$^\circ$ &
      \textbf{20}$^\circ$\textbf{/s}$^\mathbf{2}$ & \textbf{0.51}$^\circ$ & \textbf{1.35\,mm} \\
    0.15\,m/s$^2$ & 2.9\,mm (3.5$^*$) & 2.6\,mm & 0.47$^\circ$ & 39$^\circ$/s$^2$ & 1.31$^\circ$ & 4.41\,mm \\
    \bottomrule
  \end{tabular}}

  \smallskip
  {\small Reading: the criterion (lag $\le 2$\,mm, tilt $\le 1^\circ$, angle lag $\le 2^\circ$, tip drift $\le 2$\,mm)
  is met at 0.1\,m/s$^2$ and 20$^\circ$/s$^2$ (bold, chosen) and fails at 0.15\,m/s$^2$ / 39$^\circ$/s$^2$. The lag
  of 1.9\,mm matches Eq.~\eqref{eq:al-lag} with $\Lambda \approx 20$\,kg. $^*$ at the two start points low and close
  to the robot ($x$ 0.35, $z$ 0.20\,m), where joint 4 is within 0.21\,rad of its limit; the task will keep the tool
  away from there.}
\end{table}

\begin{table}[h]
  \centering
  \caption{\textbf{Revised sweep (8 October): real reversals and full-speed moves.} $K = 1000$\,N/m,
  $K_o = 100$\,N\,m/rad, rotation limits unchanged (45$^\circ$/s, 20$^\circ$/s$^2$). \emph{Short moves}: as in
  Table~\ref{tab:al-b1}, but reversals that brake and reverse. \emph{Full speed}: one move per axis up to the cap and
  back to rest. Worst case over the counted moves of the 9 start points (moves within 0.25\,rad of a joint limit are
  not counted).}
  \label{tab:al-final}
  {\renewcommand{\arraystretch}{1.15}
  \begin{tabular}{rrrrrrrr}
    \toprule
    & & \multicolumn{3}{c}{short moves, real reversals} & \multicolumn{3}{c}{full speed} \\
    \cmidrule(lr){3-5} \cmidrule(lr){6-8}
    accel.\ limit & speed cap & lag & overshoot & tilt & lag & overshoot & moves counted \\
    \midrule
    0.10\,m/s$^2$ & 20\,cm/s & 2.15\,mm & 1.65\,mm & 0.35$^\circ$ & 3.35\,mm$^*$ & 3.01\,mm & 7 of 27 \\
    0.10\,m/s$^2$ & 10\,cm/s & 2.15\,mm & 1.65\,mm & 0.35$^\circ$ & 2.00\,mm & 1.93\,mm & 21 of 27 \\
    \textbf{0.08\,m/s}$^\mathbf{2}$ & \textbf{10\,cm/s} & \textbf{1.68\,mm} & \textbf{1.32\,mm} & \textbf{0.27}$^\circ$ &
      \textbf{1.65\,mm} & \textbf{1.58\,mm} & 21 of 27 \\
    \bottomrule
  \end{tabular}}

  \smallskip
  {\small Reading: at 0.1\,m/s$^2$ the real reversals at the far, high points ($x$ 0.60, $z$ 0.45\,m) lag 2.15\,mm,
  i.e.\ $\Lambda \approx 21.5$\,kg in Eq.~\eqref{eq:al-lag}, just above the criterion. At 20\,cm/s a full-speed move is
  about 42\,cm long and leaves the comfortable workspace: only 7 of 27 moves stay clear of the joint limits, and the
  counted ones lag up to 3.35\,mm ($^*$ start pose, $+y$, while braking far out; two further moves ran beyond the
  arm's reach). At 10\,cm/s all full-speed moves of the 7 counted start points fit. \textbf{0.08\,m/s$^2$ and
  10\,cm/s} (bold, chosen) meet the criterion everywhere with about 15\,\% margin. Rotations at the same limits: angle
  lag $\le 0.62^\circ$, tip drift $\le 1.23$\,mm (full-speed turns of about 120$^\circ$ about $x$ and $y$ reach the
  wrist's joint limits at most points, so only 1--2 of 3 turns per point are counted). Not counted: the two points
  near joint 4's limit lag up to 2.3\,mm.}
\end{table}

\clearpage
% ---------------------------------------------------------------------------------------------------------------------
\subsection{Design choices and their trade-offs}
\label{sec:al-choices}
These are the building blocks of the design. Every one moves the balance between \emph{transfer robustness}
(small lag, so that the motion is set by the controller), \emph{agility} (how fast the policy can move and react),
\emph{gentleness} (forces on the plant) and \emph{effort} (implementation, identification, training time). The
current value is given in the first column.

\paragraph{Controller (what turns the target into torques).}
\begin{center}
{\small\renewcommand{\arraystretch}{1.2}
\begin{tabular}{>{\raggedright\arraybackslash}p{0.16\textwidth}>{\raggedright\arraybackslash}p{0.23\textwidth}>{\raggedright\arraybackslash}p{0.245\textwidth}>{\raggedright\arraybackslash}p{0.245\textwidth}}
  \toprule
  Building block (current) & What it is & Gains & Costs \\
  \midrule
  Translational stiffness $K$ \newline (\textbf{1000\,N/m}, was 400)
    & spring between target and tool position. Franka's examples use 150--400\,N/m; Franka's internal Cartesian
      impedance accepts up to 3000\,N/m
    & lag $\propto 1/K$, so $a_\text{max} \propto K$ (Eq.~\ref{eq:al-amax}); less position error from joint friction on
      the real arm (error $\approx$ friction$/K$)
    & more force per mm of lag on contact; the damping grows with $\sqrt K$, so encoder noise and loop delay matter
      more; stability margin on the real robot to be checked \\
  \addlinespace\cmidrule(lr){1-4}\addlinespace
  Rotational stiffness $K_o$ \newline (\textbf{100\,N\,m/rad}, was 30)
    & spring between target and tool orientation (10--300\,N\,m/rad)
    & tilt $\propto 1/K_o$ while translating (Table~\ref{tab:al-diag}); smaller angle lag when turning
    & twists a stem caught in the fork more stiffly; same noise/stability caveat as $K$ \\
  \addlinespace\cmidrule(lr){1-4}\addlinespace
  Damping $D$ \newline (\textbf{$2\sqrt K\Lambda^{1/2}$})
    & damper between target and tool velocity
    & critical in every pose and direction; Franka's $2\sqrt K$ (critical only for 1\,kg) overshot by up to 7\,mm
    & uses the model of $\Lambda$ (only for the damping) \\
  \addlinespace\cmidrule(lr){1-4}\addlinespace
  Velocity feedforward $\dot x_d$ \newline (\textbf{on})
    & the damper acts on the velocity \emph{difference}, not on the tool velocity; while the clamp holds the target,
      $\dot x_d$ is the tool's velocity (the held target moves with the tool)
    & no lag at constant speed (without it: 82\,mm at 20\,cm/s, Table~\ref{tab:al-time}); no damping force on top of
      the clamp's bound
    & none found \\
  \addlinespace\cmidrule(lr){1-4}\addlinespace
  Inertia feedforward \newline $J^\top\Lambda\ddot x_d$ \newline (\textbf{off}; option B3)
    & the controller supplies the force needed to accelerate the arm from the model (operational-space
      control~\cite{al-khatib1987}) instead of waiting for the spring to stretch
    & lag drops to $\Delta\Lambda\, a / K$, set only by the \emph{model error} $\Delta\Lambda$ (20\,\% error: 5 times less
      lag); also removes the tilt and tip drift from the coupling
    & the motion now depends on the model of $\Lambda$ (libfranka on the robot); needs the motor-inertia question
      settled and the model checked (identification) \\
  \addlinespace\cmidrule(lr){1-4}\addlinespace
  Control period \newline (\textbf{2\,ms} sim, 1\,ms robot)
    & how often the law is evaluated
    & 1\,ms in sim matches the robot exactly
    & simulation twice as slow, so training twice as long \\
  \bottomrule
\end{tabular}}
\end{center}

\clearpage
\paragraph{Commands (what the policy outputs, and the limits on it).}
\begin{center}
{\small\renewcommand{\arraystretch}{1.2}
\begin{tabular}{>{\raggedright\arraybackslash}p{0.16\textwidth}>{\raggedright\arraybackslash}p{0.23\textwidth}>{\raggedright\arraybackslash}p{0.245\textwidth}>{\raggedright\arraybackslash}p{0.245\textwidth}}
  \toprule
  Building block (current) & What it is & Gains & Costs \\
  \midrule
  Speed cap \newline (\textbf{10\,cm/s, 45$^\circ$/s}; 20\,cm/s until 8 Oct.)
    & largest target step per policy step
    & the policy can approach and turn within one episode; 10\,cm/s is reached and left again within about 13\,cm,
      i.e.\ within the stem area
    & 20\,cm/s needed 42\,cm to reach and leave, more than the stem area, and lagged up to 3.35\,mm on such moves
      (Table~\ref{tab:al-final}) \\
  \addlinespace\cmidrule(lr){1-4}\addlinespace
  Acceleration limit \newline (\textbf{0.08\,m/s$^2$, 20$^\circ$/s$^2$}; 0.1\,m/s$^2$ until 8 Oct.)
    & largest change of the target step from one policy step to the next
    & \emph{the} knob for the lag (Eq.~\ref{eq:al-lag}); part of the policy-side code, so identical in sim and on the
      robot; the policy sees the applied step, so it can learn the limit
    & less agility: long time to speed and long braking distance (table in Section~\ref{sec:al-problem}) \\
  \addlinespace\cmidrule(lr){1-4}\addlinespace
  Lag criterion $e_\text{max}$ \newline (\textbf{2\,mm, 2$^\circ$})
    & how much lag we accept; sets the acceleration limit via Eq.~\eqref{eq:al-amax}
    & smaller: the motion depends less on the arm's dynamics (better transfer~\cite{al-aljalbout2024})
    & smaller $a_\text{max}$; a larger value (e.g.\ 5\,mm, option B2) gives 2.5 times more acceleration \\
  \addlinespace\cmidrule(lr){1-4}\addlinespace
  Low-pass filter on the commands \newline (\textbf{not used})
    & smooths the commands instead of limiting their change
    & smooth motion; common in sim-to-real
    & adds hidden state; Ref.~\cite{al-aljalbout2024} saw learning collapse with a limiter the policy could not see;
      no hard bound on the acceleration \\
  \addlinespace\cmidrule(lr){1-4}\addlinespace
  Target clamp \newline $\lVert x_d - x\rVert \le e_\text{clamp}$ \newline (\textbf{4\,mm, 3$^\circ$}, every control
    step)
    & the target is never further than $e_\text{clamp}$ from the measured tool; checked at every control step, not
      once per policy step (which would also cut steps in free motion)
    & hard bound on the spring force when the tool is blocked: 4\,mm $\times$ 1000\,N/m = 4\,N (about 3 times the
      measured push forces of 0.8--1.4\,N), 3$^\circ \times$ 100\,N\,m/rad = 5.2\,N\,m (43\,\% of the FR3's 12\,N\,m
      wrist torque limit; 10$^\circ$ would exceed it), even if the policy keeps pushing
    & must stay above the free-space lag (2\,mm), otherwise it acts while moving; while it acts, the target does not
      follow the commands (the policy sees this in the applied step); runs in the 1\,kHz loop on the robot \\
  \addlinespace\cmidrule(lr){1-4}\addlinespace
  Policy rate \newline (\textbf{31.25\,Hz})
    & how often the policy sends a new step
    & higher: finer commands, quicker reaction
    & more policy evaluations per second of simulation, so longer training; more sensitive to latency \\
  \addlinespace\cmidrule(lr){1-4}\addlinespace
  Action space \newline (\textbf{Cartesian steps on the previous target})
    & alternatives: joint velocities, Cartesian velocities, absolute poses, joint torques
    & joint velocities transferred best in Ref.~\cite{al-aljalbout2024} (stiff joint-level tracking)
    & joint velocities: no Cartesian compliance at contact, which matters for a fragile stem; torques: motion set by
      the dynamics alone, could not be transferred safely~\cite{al-aljalbout2024} \\
  \bottomrule
\end{tabular}}
\end{center}

\clearpage
\paragraph{Task and training.}
\begin{center}
{\small\renewcommand{\arraystretch}{1.2}
\begin{tabular}{>{\raggedright\arraybackslash}p{0.16\textwidth}>{\raggedright\arraybackslash}p{0.23\textwidth}>{\raggedright\arraybackslash}p{0.245\textwidth}>{\raggedright\arraybackslash}p{0.245\textwidth}}
  \toprule
  Building block (current) & What it is & Gains & Costs \\
  \midrule
  Episode length \newline (\textbf{10\,s})
    & how long the policy has per attempt
    & longer: a slow approach still leaves time to push the stem
    & training time grows about in proportion; the reward comes later after the actions that earned it, which makes
      learning harder \\
  \addlinespace\cmidrule(lr){1-4}\addlinespace
  Workspace bound \newline (\textbf{planned})
    & keep the tool away from poses near joint limits and the edge of the reach
    & there the lag grows and the real FR3 stops with a reflex
    & excludes some stem placements \\
  \addlinespace\cmidrule(lr){1-4}\addlinespace
  Reward terms \newline (\textbf{planned}: contact force, action rate, joint-limit margin)
    & penalties the policy learns to avoid, instead of hard limits
    & the policy learns gentle motion itself, where it matters
    & soft: no guarantee on the robot \\
  \addlinespace\cmidrule(lr){1-4}\addlinespace
  Domain randomization \newline (\textbf{planned}, after a first baseline run)
    & vary $K$, $K_o$, $\Lambda$ (motor inertia, link masses), latency and friction in training
    & policy robust to the remaining differences between sim and real~\cite{al-peng2018}
    & ranges must be realistic; too wide makes the policy cautious and slow to train \\
  \addlinespace\cmidrule(lr){1-4}\addlinespace
  Identification \newline (\textbf{planned}, before the first real run)
    & compare libfranka's model with the simulation; record the real arm's response to the same commands
    & tells how large the dynamics gap is; prerequisite for inertia feedforward and for loosening the criterion
    & robot time \\
  \addlinespace\cmidrule(lr){1-4}\addlinespace
  FR3 safety layer in sim \newline (\textbf{planned})
    & joint velocity, torque and torque-rate limits and the collision reflex as terminations/penalties
    & the policy does not learn motions the robot refuses
    & thresholds must match the robot's configuration \\
  \bottomrule
\end{tabular}}
\end{center}

\paragraph{Main options, in short.}
\begin{center}
{\renewcommand{\arraystretch}{1.2}
\begin{tabular}{>{\raggedright\arraybackslash}p{0.24\textwidth}>{\raggedright\arraybackslash}p{0.18\textwidth}>{\raggedright\arraybackslash}p{0.22\textwidth}>{\raggedright\arraybackslash}p{0.26\textwidth}}
  \toprule
  Option & $a_\text{max}$ at 2\,mm lag & Depends on the arm model & Main cost \\
  \midrule
  B1: $K$ 1000, $K_o$ 100, no inertia feedforward & 0.08\,m/s$^2$ (measured; 0.1 lags 2.15\,mm) & only through the
    damping & sluggish \\
  \addlinespace
  B2: B1 with a 5\,mm criterion & $\approx 0.25$\,m/s$^2$ & yes: 5\,mm of lag are set by the arm & against the transfer
    evidence \\
  \addlinespace
  B3: B1 + inertia feedforward & $\approx 0.5$\,m/s$^2$ or more (not run) & yes: model of $\Lambda$ & identification
    first \\
  \addlinespace
  Stiffer $K$ (3000\,N/m), no inertia feedforward & $\approx 0.3$\,m/s$^2$ & only through the damping & real-robot
    stability, contact forces \\
  \bottomrule
\end{tabular}}
\end{center}

% ---------------------------------------------------------------------------------------------------------------------
\subsection{Plan}
\label{sec:al-plan}
Decision (7 October 2026, revised 8 October): \textbf{B1 for the first training and transfer}, B3 (inertia
feedforward) once a policy trains successfully and the robot's $\Lambda$ has been checked. Values from
Table~\ref{tab:al-final}: $K = 1000$\,N/m, $K_o = 100$\,N\,m/rad, speed caps 10\,cm/s (3.2\,mm per step) and
45$^\circ$/s, acceleration limits 0.08\,m/s$^2$ (0.08\,mm per step) and 20$^\circ$/s$^2$ (0.02$^\circ$ per step),
target clamp 4\,mm and 3$^\circ$ at every control step. On 8 October the speed cap was halved and the acceleration
limit lowered from 0.1\,m/s$^2$, after moves that run up to the cap and reversals that really reverse showed lags of
up to 3.35\,mm and 2.15\,mm (Table~\ref{tab:al-final}).

\paragraph{Is this fast enough for the task?} With these limits (speed up, then brake, from rest to rest): the
farthest stem tip, 31\,cm from the start pose, is reached in 4.4\,s (1.25\,s to reach 10\,cm/s, 1.85\,s at
10\,cm/s, 1.25\,s braking); a 90$^\circ$ turn takes 4.3\,s (peak 42$^\circ$/s), and turning can happen while
approaching. That leaves about 5.5\,s of the 10\,s episode for pushing. Braking from 10\,cm/s takes 6\,cm and 1.25\,s.
We consider this sufficient for a first training. If the trained policy runs out of time, we lengthen the episode
first (cost: training time) and loosen the limits only after the identification.

% ---------------------------------------------------------------------------------------------------------------------
\subsection{Questions}
\label{sec:al-questions}
\begin{enumerate}
  \item Is the analysis right: is the lag of a Cartesian impedance controller on the FR3 set by its apparent mass
        (8--22\,kg at the tool in simulation, half of it from the motors' reflected inertia)?
  \item Does the FR3's internal joint torque controller reduce the apparent motor inertia (as on the DLR
        lightweight arms), and is that why libfranka's model leaves the motor inertia out? This decides how the
        simulation must model the motors, and which mass matrix the damping of our controller should use on the
        robot.
  \item Stiffness: is $K = 1000$\,N/m (and $K_o = 100$\,N\,m/rad) in a libfranka torque loop acceptable on our FR3
        next to a plant? Has the lab run similar gains, and with which velocity filtering?
  \item Inertia feedforward (option B3): how accurate is libfranka's model in your experience? Has the lab used
        operational-space control with acceleration feedforward on the FR3?
  \item Is an acceleration of about 0.08\,m/s$^2$ (1.25\,s to 10\,cm/s) a reasonable limit near a plant, or would
        you accept larger tracking errors for more agility?
  \item Which reflexes and thresholds (collision behaviour, joint velocity limits) does the lab configure, so that we
        model the same ones in simulation?
\end{enumerate}

{\raggedright
\begin{thebibliography}{9}
\bibitem{al-aljalbout2024} E.~Aljalbout, F.~Frank, M.~Karl, P.~van der Smagt. On the role of the action space in
  robot manipulation learning and sim-to-real transfer. \emph{IEEE Robotics and Automation Letters}, 2024.
  arXiv:2312.03673.
\bibitem{al-khatib1987} O.~Khatib. A unified approach for motion and force control of robot manipulators: the
  operational space formulation. \emph{IEEE Journal on Robotics and Automation}, 3(1), 1987.
\bibitem{al-peng2018} X.~B.~Peng, M.~Andrychowicz, W.~Zaremba, P.~Abbeel. Sim-to-real transfer of robotic control
  with dynamics randomization. \emph{IEEE ICRA}, 2018.
\end{thebibliography}
}
% ======================================== END CONTENT ========================================
) or compiled alone via action_limits_report.tex.
% Uses amsmath, booktabs and array (added to main.tex for this section). Numbers: scripts/sweep_action_poses.py
% (default and --full_speed), scripts/check_push_env.py and docs/notes/2026-10-07.md, 2026-10-08.md in the project repo.

% ======================================= BEGIN CONTENT =======================================
\section{Open problem: how fast may the policy move the fork?}
\label{sec:al}

\paragraph{Status (7--8 October 2026).} Problem identified in simulation; values for the first training runs set
on 7 October (option B1) and revised on 8 October after longer test moves (Section~\ref{sec:al-plan}); the
long-term solution is open. We ask for feedback on the analysis and on the
questions at the end (Section~\ref{sec:al-questions}).

\paragraph{In one paragraph.} The policy commands the fork through a Cartesian impedance controller, the same in
simulation and on the FR3. What blocks us is a trade-off between \emph{agility} and a \emph{small tracking error}.
A small error transfers better~\cite{al-aljalbout2024}, because then the fork's motion is set by our controller
(identical on both sides) rather than by the arm's dynamics (not identical, not yet identified). But an impedance
controller is a spring of stiffness $K$ between the commanded and the actual tool position, and it can only
accelerate the arm by being stretched: while the commanded motion accelerates at $a$, the tool lags behind by about
$\Lambda a/K$, where $\Lambda$ is the arm's apparent mass at the tool tip (8--22\,kg depending on pose and direction,
half of it from the motors' reflected inertia). With a small lag (2\,mm) and a stiffness we consider acceptable near a
plant (400--1000\,N/m), the fork may accelerate only at about 0.04--0.09\,m/s$^2$; we use 0.08\,m/s$^2$ at
1000\,N/m. This section explains the
trade-off, what we measured, and every design choice that moves it.

% ---------------------------------------------------------------------------------------------------------------------
\subsection{Background}
\label{sec:al-background}

\paragraph{What the policy does.} The policy is a neural network trained by reinforcement learning (RL) in
simulation. At 31.25\,Hz (every 32\,ms, one ``policy step'') it reads the state (fork pose, stem shape, target) and
outputs a small displacement and rotation of the fork's \emph{target} pose $x_d$. The next target is the previous
target plus this step, so commanded steps are not lost even if the tool lags behind. Within the 32\,ms the target
moves at constant velocity $\dot x_d$ from the old to the new value. A Cartesian impedance controller at 1\,kHz
(2\,ms in simulation) turns the target into joint torques:
\begin{equation}
  \tau = J^\top\!\left[\,K\,(x_d - x) + D\,(\dot x_d - \dot x)\,\right] + \tau_\text{null} + \tau_\text{Coriolis}.
  \label{eq:al-law}
\end{equation}
$x$ is the measured tool pose, $J$ the tool-tip Jacobian, $K$ the stiffness and $D$ the damping (written here for
the position; the orientation has its own stiffness $K_o$). This is the law of Franka's example controllers, with
two changes: the target velocity $\dot x_d$ is fed forward (Franka's examples use $\dot x_d = 0$), and the damping
$D = 2\sqrt{K}\,\Lambda^{1/2}$ is chosen critical for the arm's actual inertia $\Lambda$ (Section~\ref{sec:al-problem})
instead of Franka's $2\sqrt K$, which is critical only for a 1\,kg tool.

\paragraph{Null-space spring $\tau_\text{null}$.} The FR3 has 7 joints, but the tool pose has only 6 degrees of
freedom. For every tool pose there is therefore a one-parameter family of arm postures (the elbow can swing about the
line from shoulder to wrist without moving the tool): the \emph{null space}. The first term of
Eq.~\eqref{eq:al-law} does not act in this direction, so the elbow would drift freely. A weak joint-space spring
(0.5\,N\,m/rad, as in Franka's examples) pulls the posture towards a rest posture; it is projected so that it does
not push the tool. It plays no role in the problem below (we checked: switching it off does not change the lag).

\paragraph{What RL learns, and why the action space matters for the transfer.} During training the policy tries
millions of commands in simulation and keeps what earns reward. It learns \emph{how the simulated fork reacts to its
commands}, including any lag, overshoot or coupling, and relies on it. If the real robot reacts differently, the
learned behaviour is wrong by that difference: this is the sim-to-real (``reality'') gap. The policy cannot correct
what it never saw in simulation. The usual tools against the gap are:
\begin{center}
{\renewcommand{\arraystretch}{1.15}
\begin{tabular}{>{\raggedright\arraybackslash}p{0.24\textwidth}>{\raggedright\arraybackslash}p{0.42\textwidth}>{\raggedright\arraybackslash}p{0.24\textwidth}}
  \toprule
  Tool & What it does & Status \\
  \midrule
  A. Action-space design & choose commands whose result depends little on what differs between sim and real
    (same controller, small tracking error) & controller matched; this section \\
  \addlinespace
  B. Identification & measure the real robot and correct the simulation & planned, before the first real run \\
  \addlinespace
  C. Domain randomization & vary unknown parameters in training, so the policy works over a range & planned \\
  \addlinespace
  D. Model the robot's safety layer & limits and reflexes of the FR3 become terminations/penalties in sim, so the
    policy does not learn motions the robot refuses & planned \\
  \addlinespace
  E. Adaptation & infer differences online or fine-tune on the robot & not planned \\
  \bottomrule
\end{tabular}}
\end{center}
A large study of action spaces on a Franka arm~\cite{al-aljalbout2024} (13 action spaces, about 250 policies,
reaching and pushing, transferred to the real robot) concludes that \emph{``an action space which yields or
naturally limits the tracking error transfers better''}; torque actions, whose result is set by the robot's
dynamics alone, could not be transferred safely. The same study warns that a limiter which acts on the robot but is
not visible to the policy in simulation (they tried a rate limiter) made learning collapse. A caveat: their controllers were joint
impedance controllers (Cartesian actions were converted to joint targets by inverse kinematics), with gains that
the paper does not report, and the tasks were reaching and pushing a box. The conclusion carries over in principle,
since in every impedance-type controller a larger tracking error means more dependence on the arm's dynamics; the
numbers do not, and the paper does not test a soft Cartesian impedance controller like ours.

\paragraph{Where the gap enters our system.} The chain from a policy output to the stem, step by step:
\begin{center}
{\renewcommand{\arraystretch}{1.15}
\begin{tabular}{>{\raggedright\arraybackslash}p{0.32\textwidth}>{\raggedright\arraybackslash}p{0.20\textwidth}>{\raggedright\arraybackslash}p{0.38\textwidth}}
  \toprule
  Step & Same in sim and real? & Source of differences \\
  \midrule
  policy step $\to$ target pose & yes (same code) & none \\
  target $\to$ torques, law~\eqref{eq:al-law} & yes (ported to C++, checked) & port errors; $J$, $M$ from different
    models \\
  \textbf{torques $\to$ fork motion} & \textbf{no} & arm inertia, joint friction, motor inertia, delays, 1 vs.\ 2\,ms
    control period \\
  fork $\leftrightarrow$ stem contact & no & friction, stem stiffness and damping (randomized) \\
  stem $\to$ observation & no & camera noise, latency \\
  FR3 safety layer & \textbf{missing in sim} & reflexes stop the robot \\
  \bottomrule
\end{tabular}}
\end{center}
The step ``torques $\to$ fork motion'' is where the arm's dynamics act. How much of it reaches the fork's motion
depends on the tracking error: if the tool follows the target within a small error, its motion is essentially the
target's motion, which our code computes identically in sim and real; the larger the error, the more the motion is
shaped by the arm's dynamics, which differ.

% ---------------------------------------------------------------------------------------------------------------------
\subsection{The problem}
\label{sec:al-problem}

\paragraph{What blocks us: agility versus a small tracking error.} We want two things from the fork's motion:
\begin{enumerate}
  \item \textbf{Small lag behind the commanded target}, so that the fork's motion is set by our controller, which
        is the same in simulation and on the FR3, and not by the arm's dynamics, which differ and are not yet
        identified. Action spaces that keep the tracking error small transfer better~\cite{al-aljalbout2024}.
  \item \textbf{Agility}: reach the stem well within the 10\,s episode, brake late before contact, and react quickly
        to how the stem moves.
\end{enumerate}
With an impedance controller the two conflict. The lag grows with how quickly the commanded motion accelerates, so
a small lag allows only gentle accelerations. In numbers: with $K = 1000$\,N/m and at most 2\,mm of lag, the fork may
accelerate at only about 0.09\,m/s$^2$ where the arm is heaviest; at the 0.08\,m/s$^2$ we use, it needs 1.25\,s to
reach 10\,cm/s and 6\,cm to stop from that speed (2.5\,s and 25\,cm for 20\,cm/s). Either
we accept this sluggishness, or we accept more lag (and with it a larger dependence on the real arm), or we change
the controller. This is the decision we need help with. The rest of this section explains where the conflict comes
from, how severe it is for the task, and how severe it will be on the real robot.

\paragraph{Why the tool lags: the impedance law as a mass on a spring.} Consider one direction, free space,
gravity compensated. The controller makes the tool tip behave like a mass $\Lambda$ (the arm's inertia felt at the
tool tip, explained in the next paragraph) attached by a spring $K$ and a damper $D$ to the moving target:
\begin{equation}
  \Lambda\,\ddot x = K\,(x_d - x) + D\,(\dot x_d - \dot x).
\end{equation}
Write the lag as $e = x_d - x$ (target minus tool). Then $\ddot x = \ddot x_d - \ddot e$, and substituting gives
\begin{equation}
  \Lambda\,\ddot e + D\,\dot e + K\,e = \Lambda\,\ddot x_d.
  \label{eq:al-error}
\end{equation}
The lag itself behaves like a damped mass-spring system, driven by the force $\Lambda \ddot x_d$, i.e.\ by the
\emph{acceleration of the target}. Three cases follow:
\begin{enumerate}
  \item \textbf{Target at rest or moving at constant speed} ($\ddot x_d = 0$): the right-hand side is zero, so any lag
        dies out and $e \to 0$. Physically: moving at constant speed in free space needs no force (no friction, no
        gravity), so the spring does not need to be stretched. This holds only because the target velocity is fed
        forward. Without it the damper pulls back with $D\dot x$ all the time, and the spring must balance it:
        $e = D\dot x_d / K$. With $D = 2\sqrt{K\Lambda} = 180$\,N\,s/m ($K = 400$\,N/m, $\Lambda = 20$\,kg) and
        20\,cm/s this gives 90\,mm; we measured 82\,mm (Table~\ref{tab:al-time}).
  \item \textbf{Target accelerating at constant $a$}: the right-hand side is the constant force $\Lambda a$. Once the
        transient has died out ($\ddot e = \dot e = 0$), $K e = \Lambda a$, so
        \begin{equation}
          e = \frac{\Lambda\, a}{K}.
          \label{eq:al-lag}
        \end{equation}
        Physically: Newton's law. To accelerate the mass $\Lambda$ at $a$, the force $\Lambda a$ is needed, and its only
        source is the stretched spring. Example: $\Lambda = 20$\,kg, $a = 0.1$\,m/s$^2$, $K = 1000$\,N/m: $e = 2$\,mm.
        The same holds when braking (the tool then runs ahead of the target, which shows up as overshoot).
  \item \textbf{How fast the lag builds up and dies out}: Eq.~\eqref{eq:al-error} is critically damped, with the time
        constant $\sqrt{\Lambda/K}$: 0.14\,s for 20\,kg and 1000\,N/m, 0.22\,s for 400\,N/m, i.e.\ 4--7 policy steps.
        The lag reaches the value of Eq.~\eqref{eq:al-lag} if the acceleration lasts longer than about three time
        constants, and it disappears again within a similar time once the target moves at constant speed. This
        matters for testing: short test moves end before the lag has settled. Our first sweep (20 policy steps per
        segment, reversals that only braked) therefore showed 1.9\,mm at 0.1\,m/s$^2$; moves that brake \emph{and}
        reverse, or that run up to the speed cap, accelerate long enough and reached 2.15\,mm
        (Table~\ref{tab:al-final}).
\end{enumerate}

\paragraph{The mass in this equation: the apparent mass $\Lambda$ (``8--20\,kg at the tool tip'').} The trade-off
is severe because the mass that the spring has to accelerate is large. If one pushes the tool tip of
the arm (joints free, gravity compensated) with a force $F$, the tip starts to accelerate with $\ddot x = \Lambda^{-1}
F$. $\Lambda$ is the arm's inertia as seen at the tool tip, the \emph{apparent mass} (or operational-space inertia):
\begin{equation}
  \Lambda = \left(J M^{-1} J^\top\right)^{-1},
  \label{eq:al-lambda}
\end{equation}
with $M$ the $7\times7$ joint-space mass matrix and $J$ the $6\times7$ tool-tip Jacobian. $\Lambda$ is a $6\times6$
matrix (force and moment versus linear and angular acceleration) and depends on the pose. Its translational
diagonal entries say how heavy the arm feels when the tip is pushed along $x$, $y$ or $z$. At the earlier start pose
(2026-10-07; tool tip at $(0.30, 0, 0.55)$\,m in front of the robot, fork horizontal), computed from the simulated FR3:
\begin{center}
{\renewcommand{\arraystretch}{1.15}
\begin{tabular}{lrrr}
  \toprule
  Apparent mass, start pose & along $x$ (forward) & along $y$ (sideways) & along $z$ (up) \\
  \midrule
  links only & 11.9\,kg & 4.2\,kg & 5.1\,kg \\
  links + motors' reflected inertia (our simulation) & 20.5\,kg & 7.6\,kg & 10.6\,kg \\
  \bottomrule
\end{tabular}}
\end{center}
Pushing the tip forward with 20\,N accelerates it at about 1\,m/s$^2$, as if it were a 20\,kg mass; the same push
sideways gives 2.6\,m/s$^2$. Forward motion needs the heavy shoulder and elbow joints, sideways motion mostly the
base rotation and the wrist. The second row includes the motors' reflected inertia $n^2 I_\text{motor}$ ($n$: gear
ratio), which we add to $M$ from Franka's description files: 0.61, 0.61, 0.46, 0.46, 0.21, 0.21, 0.21\,kg\,m$^2$ for
joints 1--7 (at the wrist about 100 times the link's own inertia). Without these terms the simulated arm oscillated
and pressed the stem 5 times harder. The lag measurements below give the same value independently: from
Eq.~\eqref{eq:al-lag}, $\Lambda = K e / a \approx 16$--20\,kg along $x$ at the start pose, and up to 21.5\,kg
at the far, high start points of the sweep ($x$ 0.60, $z$ 0.45\,m).

\paragraph{How severe the trade-off is for the task.} Turning Eq.~\eqref{eq:al-lag} around: for a largest acceptable lag $e_\text{max}$, the
target may accelerate at most at
\begin{equation}
  a_\text{max} \approx \frac{K\, e_\text{max}}{\Lambda}.
  \label{eq:al-amax}
\end{equation}
With $\Lambda = 21.5$\,kg (the heaviest pose and direction) and $e_\text{max} = 2$\,mm: 0.04\,m/s$^2$ at
$K = 400$\,N/m and 0.09\,m/s$^2$ at 1000\,N/m. We use 0.08\,m/s$^2$, which leaves about 15\,\% margin (lag
$\le 1.72$\,mm; measured 1.68\,mm, Table~\ref{tab:al-final}). What this means for the task (the
target speeds up at $a_\text{max}$ until the speed cap, then brakes at $a_\text{max}$):
\begin{center}
{\renewcommand{\arraystretch}{1.15}
\begin{tabular}{rrrrr}
  \toprule
  largest acceleration & time from rest & braking distance & 31\,cm, rest to rest & 31\,cm, rest to rest \\
  $a_\text{max}$ & to 20\,cm/s & from 20\,cm/s & (20\,cm/s cap) & (10\,cm/s cap) \\
  \midrule
  0.05\,m/s$^2$ & 4.0\,s & 40\,cm & 5.0\,s & 5.1\,s \\
  \textbf{0.08\,m/s}$^\mathbf{2}$ & 2.5\,s & 25\,cm & 3.9\,s & \textbf{4.4\,s} \\
  0.10\,m/s$^2$ & 2.0\,s & 20\,cm & 3.5\,s & 4.1\,s \\
  0.25\,m/s$^2$ & 0.8\,s & 8\,cm  & 2.3\,s & 3.5\,s \\
  0.50\,m/s$^2$ & 0.4\,s & 4\,cm  & 1.9\,s & 3.3\,s \\
  \bottomrule
\end{tabular}}
\end{center}
31\,cm is the distance from the start pose to the farthest stem tip. The episode lasts 10\,s, and the stem must
still be pushed after the approach. The cost of a small $a_\text{max}$ is therefore not the top speed but agility:
the policy must start braking long before it touches the stem and cannot react quickly to how the stem moves.

\paragraph{Why not simply make the episode longer?} A longer episode gives a slow policy time to finish, but it is
not free. RL learns from simulated experience, so the training time grows roughly in proportion to the simulated
time per episode: a 20\,s episode costs about twice the training time of a 10\,s episode for the same number of
attempts. In addition, the reward for reaching the target arrives later after the actions that caused it, which
makes it harder for the policy to learn which actions were good. We therefore keep 10\,s as long as the task fits,
and lengthen it only if the trained policy demonstrably runs out of time.

\paragraph{Rotations: the same effect, plus a coupling.} Turning the fork works in the same way: the orientation
spring $K_o$ must twist to give the rotational inertia an angular acceleration $\alpha$, so the angle lags by about
$\Lambda_o \alpha / K_o$ ($\Lambda_o$: the rotational part of $\Lambda$). In addition, $\Lambda$ is a full $6\times6$
matrix: its off-diagonal blocks couple translation and rotation. The reason is geometric: the tool tip is 10--20\,cm
away from the wrist's joint axes, and the joints that turn the fork also move the tip, while those that move the
tip also turn the fork. The control law, however, treats position and orientation as two independent springs.
Two consequences, both measured:
\begin{itemize}
  \item \textbf{Tilt while translating.} Accelerating the tip at $a$ needs, besides the force, a moment
        $\Lambda_{op}\, a$ ($\Lambda_{op}$: off-diagonal block). It comes from the orientation spring, which must twist:
        the fork tilts by about $\Lambda_{op} a / K_o$. Table~\ref{tab:al-diag} shows the tilt falling with $1/K_o$
        (2.0$^\circ$ at $K_o = 30$, 0.26$^\circ$ at 300\,N\,m/rad).
  \item \textbf{Tip drift while turning.} Turning the fork at angular acceleration $\alpha$ needs, besides the moment, a
        force at the tip, which the position spring supplies by stretching: the tip moves away from its target by
        about $\Lambda_{po}\,\alpha / K$, although the position command is constant (Tables~\ref{tab:al-rot}
        and~\ref{tab:al-b1}).
\end{itemize}
Both effects scale with the acceleration, like the lag. A controller with inertia feedforward (operational-space
control, option B3 below) would supply these coupling forces and moments from the model and remove both.

\paragraph{Contact.} When the tool is blocked (e.g.\ pressed against the stem), the target keeps moving and the
lag becomes force: $K e$. Small lags therefore also mean small unintended forces. A hard bound is added by
clamping the target: at every control step, $\lVert x_d - x \rVert \le 4$\,mm (spring force $\le 4$\,N at
$K = 1000$\,N/m) and the angle between target and tool orientation $\le 3^\circ$ (moment $\le 5.2$\,N\,m at
$K_o = 100$\,N\,m/rad). In free motion the clamp is inactive, because the lag stays below 2\,mm. While it holds the
target, the target moves with the tool, so the target velocity fed forward is the tool's own velocity: the damper
then adds no force to the bounded spring force, and the controller does not switch back and forth between holding and
releasing. Not bounded by the clamp: the momentum of the arm when the moving fork first hits something (as with any
impedance controller); the policy is meant to learn to approach gently (contact-force and smoothness penalties).

\paragraph{How severe the trade-off will be on the real robot: which $\Lambda$ does the FR3 have?} All numbers above
use the simulated $\Lambda$. The real arm's $\Lambda$ is not known yet, and two observations from Franka's software
suggest that it may differ:
\begin{itemize}
  \item libfranka (current version, September 2026) computes its model (\texttt{franka::Model::mass}) with the
        Pinocchio library from the URDF that the robot provides. The FR3 URDF shipped with libfranka's tests has
        the same link inertias as Franka's description files but \emph{no} motor-inertia entries, and libfranka never
        sets Pinocchio's rotor-inertia (``armature'') term. So libfranka's mass matrix most likely \emph{excludes}
        the motors' reflected inertia. (Read in the source code, not yet checked on our robot.)
  \item The FR3 executes commanded torques with an internal joint torque controller that uses the link-side
        torque sensors. On the DLR lightweight arms, from which Franka's design derives, such torque feedback
        reduces the apparent motor inertia. If the FR3 does the same, the real arm feels lighter than our
        simulation (between the two rows of the table above), and libfranka's model would be the right one.
\end{itemize}
How much this matters depends on the option chosen (Section~\ref{sec:al-choices}):
\begin{center}
{\renewcommand{\arraystretch}{1.15}
\begin{tabular}{>{\raggedright\arraybackslash}p{0.24\textwidth}>{\raggedright\arraybackslash}p{0.66\textwidth}}
  \toprule
  Option & Role of the real $\Lambda$ \\
  \midrule
  B1 (small lag, now) & \textbf{Little}, and this is exactly why small lags transfer well: the unknown $\Lambda$
    acts on the motion only through the lag. The acceleration limit was sized for $\le 2$\,mm with the simulated
    $\le 21.5$\,kg; a real arm of 12\,kg would lag 1\,mm, a heavier one (30\,kg) 2.4\,mm. The policy is then wrong by about a
    millimetre. \textbf{One exception: the damping} $D = 2\sqrt K \Lambda^{1/2}$ is computed from a model of $\Lambda$.
    Before the first real run we must decide which mass matrix the controller on the robot uses (libfranka's, or
    libfranka's plus the motor terms); with the wrong one the tool is under- or over-damped. \\
  \addlinespace
  B2 (larger lag) & \textbf{More}: with 5\,mm of lag set by the arm, a 40\,\% error in $\Lambda$ shifts the motion by
    2\,mm. \\
  \addlinespace
  B3 (inertia feedforward, later) & \textbf{Central}: the controller supplies the accelerating force from the model
    of $\Lambda$, so the remaining lag is set by the model error. Requires the identification first. \\
  \bottomrule
\end{tabular}}
\end{center}

% ---------------------------------------------------------------------------------------------------------------------
\subsection{Measurements}
\label{sec:al-measurements}

\paragraph{Setup.} Simulation (Isaac Lab, PhysX) of the FR3 with the fork, gravity compensated, no joint friction,
control period 2\,ms, policy step 32\,ms, no stem (free space). Every test commands the fork with the same commands
a policy could send, through the full controller.

\paragraph{Commands and limits.} A command is a target step per policy step. Three limits act on it:
\begin{itemize}
  \item \emph{speed cap}: the largest step, e.g.\ 3.2\,mm per 32\,ms = 10\,cm/s (rotation: 1.44$^\circ$ = 45$^\circ$/s);
  \item \emph{acceleration limit}: the largest change of the step from one policy step to the next. A change of
        0.1\,mm per step means $a = 0.1\,\text{mm}/(32\,\text{ms})^2 \approx 0.1$\,m/s$^2$, 0.08\,mm per step
        0.08\,m/s$^2$ (rotation: 0.02$^\circ$ per step $\approx$ 20$^\circ$/s$^2$). Commanding ``full speed'' from rest therefore makes the target speed up at
        this acceleration until the speed cap; commanding ``stop'' makes it brake at the same rate;
  \item the first design (Table~\ref{tab:al-first}) had a speed cap only, so the target jumped from rest to full
        speed in one step.
\end{itemize}

\paragraph{Test moves.} \emph{Straight move}: full command along one axis of the robot base frame ($\pm x$, $\pm y$,
$\pm z$) for a number of policy steps, then zero command (``hold'') for 45 steps (1.4\,s), during which the target
brakes and stops. \emph{Reversal}: full command along $+x$, then immediately along $-x$ (same for $y$, $z$), then hold;
this contains the largest change of velocity. Since 8 October the second segment is twice as long, so that the
target brakes \emph{and} reverses (before, it only braked to rest). \emph{Full-speed move} (since 8 October): from
rest up to the speed cap, 10 steps at the cap, then braking to rest (about 16\,cm; rotations about 120$^\circ$),
one per axis, in a direction with room for it (towards the middle of the workspace; up along $z$).
\emph{Rotation moves}: the same about the base axes, with zero position command. \emph{Start points}: either the
start pose only, or 9 points spanning the workspace in which the stems stand (tool tip at $x$ = 0.35 and 0.60\,m,
$y = \pm0.20$\,m, $z$ = 0.20 and 0.45\,m, plus the start pose). Moves that come within 0.25\,rad of a joint limit
are reported but not counted (decision 8 October): there the arm, not the controller, limits the motion. This
concerns the two start points low and close to the robot ($x$ 0.35, $z$ 0.20\,m), where joint 4 is within
0.21\,rad of its limit.

\paragraph{What is measured (all tables).}
\begin{itemize}
  \item \emph{Lag} (following error): the largest distance between the target position and the actual tool-tip
        position during the move and the hold. The quantity of Eq.~\eqref{eq:al-lag}.
  \item \emph{Overshoot}: how far the tool tip goes past the target's final position after it stops.
  \item \emph{Tilt}: the largest change of the fork's orientation during a straight move, although no rotation is
        commanded.
  \item \emph{Angle lag}: during rotation moves, the largest angle between the target and the actual orientation.
  \item \emph{Tip drift}: during rotation moves, the largest displacement of the tool tip, although no translation is
        commanded.
\end{itemize}
Each table reports the worst value over all moves (and start points). The criterion used so far: lag $\le 2$\,mm,
overshoot $\le 5$\,mm, tilt $\le 1^\circ$; angle lag $\le 2^\circ$, tip drift $\le 2$\,mm.

\begin{table}[h]
  \centering
  \caption{\textbf{First design}: speed cap only, no acceleration limit, $K = 400$\,N/m, $K_o = 30$\,N\,m/rad. Straight
  moves of 30 policy steps at full command from 9 start points, so each move starts with a jump from rest to full speed.
  Lag measured over the last quarter of each move (as defined at the time), overshoot after it. Worst case over all
  moves and start points. This design kept the lag below 2\,mm by limiting the speed to 5.5\,cm/s.}
  \label{tab:al-first}
  {\renewcommand{\arraystretch}{1.15}
  \begin{tabular}{rrrr}
    \toprule
    speed cap (step per 32\,ms) & speed & lag & overshoot \\
    \midrule
    2.0\,mm  & 6.3\,cm/s & 2.25\,mm & 3.9\,mm \\
    1.75\,mm & 5.5\,cm/s & 1.94\,mm & 3.4\,mm \\
    1.5\,mm  & 4.7\,cm/s & 1.63\,mm & 2.9\,mm \\
    \bottomrule
  \end{tabular}}

  \smallskip
  {\small Reading: with a cap of 5.5\,cm/s and 7.8$^\circ$/s the approach to a far stem took more than half of the
  episode, and a 90$^\circ$ turn did not fit in one. We therefore raised the caps to 20\,cm/s and 45$^\circ$/s and
  added the acceleration limit; Tables~\ref{tab:al-time}--\ref{tab:al-b1} show what that does.}
\end{table}

\begin{table}[h]
  \centering
  \caption{\textbf{How the lag evolves over time} (illustrates Eq.~\eqref{eq:al-error}). Start pose, $K = 400$\,N/m,
  $K_o = 30$\,N\,m/rad. One straight move along $+x$ at full command for 40 policy steps. With the acceleration limit
  of 0.49\,m/s$^2$ the target speeds up during the first 13 steps (0.42\,s) and then moves at constant speed.
  Lag of the tool tip behind the target after 10, 20, 30 and 40 policy steps.}
  \label{tab:al-time}
  {\renewcommand{\arraystretch}{1.15}
  \begin{tabular}{lrrrr}
    \toprule
    & \multicolumn{4}{c}{lag after $n$ policy steps} \\
    \cmidrule(lr){2-5}
    & 10 (0.32\,s) & 20 (0.64\,s) & 30 (0.96\,s) & 40 (1.28\,s) \\
    & accelerating & constant speed & constant speed & constant speed \\
    \midrule
    20\,cm/s, velocity feedforward on  & 9.4\,mm & 10.4\,mm & 4.7\,mm & 1.9\,mm \\
    10\,cm/s, velocity feedforward on  & 7.7\,mm & 4.7\,mm  & 2.0\,mm & 0.7\,mm \\
    20\,cm/s, velocity feedforward off & 23\,mm  & 59\,mm   & 75\,mm  & 82\,mm \\
    \bottomrule
  \end{tabular}}

  \smallskip
  {\small With feedforward the lag builds up while the target accelerates and dies out at constant speed (case 1 and 2
  of Section~\ref{sec:al-problem}); without it, the lag grows towards $D\dot x_d/K$. Switching off the null-space
  spring gave the same numbers.}
\end{table}

\begin{table}[h]
  \centering
  \caption{\textbf{Effect of the stiffnesses $K$ and $K_o$.} Start pose, speed cap 20\,cm/s. A straight move along $+x$
  at full command for 25 policy steps (0.8\,s, enough to reach about the speed cap), hold for 45 steps, then the same
  along $-x$ and hold. Two acceleration limits (columns). Largest lag of the tool tip and largest tilt of the fork
  over the whole test.}
  \label{tab:al-diag}
  {\renewcommand{\arraystretch}{1.15}
  \begin{tabular}{rrrrrr}
    \toprule
    & & \multicolumn{2}{c}{acceleration 0.24\,m/s$^2$} & \multicolumn{2}{c}{acceleration 0.49\,m/s$^2$} \\
    \cmidrule(lr){3-4} \cmidrule(lr){5-6}
    $K$ [N/m] & $K_o$ [N\,m/rad] & lag & tilt & lag & tilt \\
    \midrule
    400  & 30  & 9.7\,mm  & 1.80$^\circ$ & 13.6\,mm & 2.40$^\circ$ \\
    400  & 100 & 10.5\,mm & 0.66$^\circ$ & 14.4\,mm & 0.84$^\circ$ \\
    400  & 300 & 10.6\,mm & 0.23$^\circ$ & 14.8\,mm & 0.29$^\circ$ \\
    1000 & 30  & 4.0\,mm  & 2.03$^\circ$ & 6.4\,mm  & 2.69$^\circ$ \\
    1000 & 100 & 4.6\,mm  & 0.73$^\circ$ & 7.1\,mm  & 1.00$^\circ$ \\
    1000 & 300 & 4.8\,mm  & 0.26$^\circ$ & 7.5\,mm  & 0.37$^\circ$ \\
    \bottomrule
  \end{tabular}}

  \smallskip
  {\small Reading: the lag depends on $K$ (2.5 times stiffer, about 2.4 times less lag) and on the acceleration, hardly
  on $K_o$; the tilt depends on $K_o$ (10 times stiffer, about 8 times less tilt). Every entry fails the 2\,mm criterion.}
\end{table}

\begin{table}[h]
  \centering
  \caption{\textbf{Workspace sweep at $K = 1000$\,N/m, $K_o = 30$\,N\,m/rad}, caps 20\,cm/s and 45$^\circ$/s. At each
  start point: straight moves along $\pm x, \pm y, \pm z$ and reversals in $x, y, z$ (30 policy steps per segment, then
  hold), then the same as rotations about the base axes. Each row is one pair of acceleration limits (translation,
  rotation). Worst case over all moves and over the five start points closest to the robot: from the start points at
  $x = 0.60$\,m these moves of up to 20\,cm ran to the edge of the arm's reach, which is a flaw of the test, not of the
  controller.}
  \label{tab:al-rot}
  {\renewcommand{\arraystretch}{1.15}
  \begin{tabular}{rrrrrr}
    \toprule
    \multicolumn{3}{c}{translation moves} & \multicolumn{3}{c}{rotation moves} \\
    \cmidrule(lr){1-3} \cmidrule(lr){4-6}
    acceleration limit & lag & tilt & angular acceleration limit & angle lag & tip drift \\
    \midrule
    0.24\,m/s$^2$ & 5.0\,mm  & 3.0$^\circ$ & 20$^\circ$/s$^2$  & 9.9$^\circ$$^*$ & 14.0\,mm$^*$ \\
    0.49\,m/s$^2$ & 8.5\,mm  & 5.0$^\circ$ & 49$^\circ$/s$^2$  & 4.7$^\circ$ & 4.0\,mm \\
    0.73\,m/s$^2$ & 11.0\,mm & 5.7$^\circ$ & 98$^\circ$/s$^2$  & 6.8$^\circ$ & 7.0\,mm \\
    0.98\,m/s$^2$ & 12.6\,mm & 6.1$^\circ$ & 195$^\circ$/s$^2$ & 9.4$^\circ$ & 14.3\,mm \\
    \bottomrule
  \end{tabular}}

  \smallskip
  {\small Reading: all rows fail the criterion; the lag grows with the acceleration as expected, the tilt is large
  because $K_o$ is soft. $^*$ one start point ($x$ 0.35, $y$ 0.20, $z$ 0.45\,m), reversal about $x$; at the other four
  points at most 4.1$^\circ$ / 3.7\,mm. Not yet understood (likely a wrist joint near its limit).}
\end{table}

\begin{table}[h]
  \centering
  \caption{\textbf{Workspace sweep for the interim choice of 7 October: $K = 1000$\,N/m, $K_o = 100$\,N\,m/rad}, caps
  20\,cm/s and 45$^\circ$/s, smaller acceleration limits. Same moves as Table~\ref{tab:al-rot} but 20 policy steps per
  segment (moves of at most about 6\,cm, so all 9 start points are valid). Worst case over all moves and all 9 start
  points. \emph{Superseded by Table~\ref{tab:al-final}}: these moves reached at most 2\,mm per step (6.25\,cm/s), not
  the caps, and the reversals only braked, so the lag had not settled.}
  \label{tab:al-b1}
  {\renewcommand{\arraystretch}{1.15}
  \begin{tabular}{rrrrrrr}
    \toprule
    \multicolumn{4}{c}{translation moves} & \multicolumn{3}{c}{rotation moves} \\
    \cmidrule(lr){1-4} \cmidrule(lr){5-7}
    acceleration limit & lag & overshoot & tilt & angular accel.\ limit & angle lag & tip drift \\
    \midrule
    0.05\,m/s$^2$ & 1.3\,mm & 0.9\,mm & 0.17$^\circ$ & 10$^\circ$/s$^2$ & 0.30$^\circ$ & 0.55\,mm \\
    \textbf{0.10\,m/s}$^\mathbf{2}$ & \textbf{1.9\,mm} (2.4$^*$) & \textbf{1.8\,mm} & \textbf{0.33}$^\circ$ &
      \textbf{20}$^\circ$\textbf{/s}$^\mathbf{2}$ & \textbf{0.51}$^\circ$ & \textbf{1.35\,mm} \\
    0.15\,m/s$^2$ & 2.9\,mm (3.5$^*$) & 2.6\,mm & 0.47$^\circ$ & 39$^\circ$/s$^2$ & 1.31$^\circ$ & 4.41\,mm \\
    \bottomrule
  \end{tabular}}

  \smallskip
  {\small Reading: the criterion (lag $\le 2$\,mm, tilt $\le 1^\circ$, angle lag $\le 2^\circ$, tip drift $\le 2$\,mm)
  is met at 0.1\,m/s$^2$ and 20$^\circ$/s$^2$ (bold, chosen) and fails at 0.15\,m/s$^2$ / 39$^\circ$/s$^2$. The lag
  of 1.9\,mm matches Eq.~\eqref{eq:al-lag} with $\Lambda \approx 20$\,kg. $^*$ at the two start points low and close
  to the robot ($x$ 0.35, $z$ 0.20\,m), where joint 4 is within 0.21\,rad of its limit; the task will keep the tool
  away from there.}
\end{table}

\begin{table}[h]
  \centering
  \caption{\textbf{Revised sweep (8 October): real reversals and full-speed moves.} $K = 1000$\,N/m,
  $K_o = 100$\,N\,m/rad, rotation limits unchanged (45$^\circ$/s, 20$^\circ$/s$^2$). \emph{Short moves}: as in
  Table~\ref{tab:al-b1}, but reversals that brake and reverse. \emph{Full speed}: one move per axis up to the cap and
  back to rest. Worst case over the counted moves of the 9 start points (moves within 0.25\,rad of a joint limit are
  not counted).}
  \label{tab:al-final}
  {\renewcommand{\arraystretch}{1.15}
  \begin{tabular}{rrrrrrrr}
    \toprule
    & & \multicolumn{3}{c}{short moves, real reversals} & \multicolumn{3}{c}{full speed} \\
    \cmidrule(lr){3-5} \cmidrule(lr){6-8}
    accel.\ limit & speed cap & lag & overshoot & tilt & lag & overshoot & moves counted \\
    \midrule
    0.10\,m/s$^2$ & 20\,cm/s & 2.15\,mm & 1.65\,mm & 0.35$^\circ$ & 3.35\,mm$^*$ & 3.01\,mm & 7 of 27 \\
    0.10\,m/s$^2$ & 10\,cm/s & 2.15\,mm & 1.65\,mm & 0.35$^\circ$ & 2.00\,mm & 1.93\,mm & 21 of 27 \\
    \textbf{0.08\,m/s}$^\mathbf{2}$ & \textbf{10\,cm/s} & \textbf{1.68\,mm} & \textbf{1.32\,mm} & \textbf{0.27}$^\circ$ &
      \textbf{1.65\,mm} & \textbf{1.58\,mm} & 21 of 27 \\
    \bottomrule
  \end{tabular}}

  \smallskip
  {\small Reading: at 0.1\,m/s$^2$ the real reversals at the far, high points ($x$ 0.60, $z$ 0.45\,m) lag 2.15\,mm,
  i.e.\ $\Lambda \approx 21.5$\,kg in Eq.~\eqref{eq:al-lag}, just above the criterion. At 20\,cm/s a full-speed move is
  about 42\,cm long and leaves the comfortable workspace: only 7 of 27 moves stay clear of the joint limits, and the
  counted ones lag up to 3.35\,mm ($^*$ start pose, $+y$, while braking far out; two further moves ran beyond the
  arm's reach). At 10\,cm/s all full-speed moves of the 7 counted start points fit. \textbf{0.08\,m/s$^2$ and
  10\,cm/s} (bold, chosen) meet the criterion everywhere with about 15\,\% margin. Rotations at the same limits: angle
  lag $\le 0.62^\circ$, tip drift $\le 1.23$\,mm (full-speed turns of about 120$^\circ$ about $x$ and $y$ reach the
  wrist's joint limits at most points, so only 1--2 of 3 turns per point are counted). Not counted: the two points
  near joint 4's limit lag up to 2.3\,mm.}
\end{table}

\clearpage
% ---------------------------------------------------------------------------------------------------------------------
\subsection{Design choices and their trade-offs}
\label{sec:al-choices}
These are the building blocks of the design. Every one moves the balance between \emph{transfer robustness}
(small lag, so that the motion is set by the controller), \emph{agility} (how fast the policy can move and react),
\emph{gentleness} (forces on the plant) and \emph{effort} (implementation, identification, training time). The
current value is given in the first column.

\paragraph{Controller (what turns the target into torques).}
\begin{center}
{\small\renewcommand{\arraystretch}{1.2}
\begin{tabular}{>{\raggedright\arraybackslash}p{0.16\textwidth}>{\raggedright\arraybackslash}p{0.23\textwidth}>{\raggedright\arraybackslash}p{0.245\textwidth}>{\raggedright\arraybackslash}p{0.245\textwidth}}
  \toprule
  Building block (current) & What it is & Gains & Costs \\
  \midrule
  Translational stiffness $K$ \newline (\textbf{1000\,N/m}, was 400)
    & spring between target and tool position. Franka's examples use 150--400\,N/m; Franka's internal Cartesian
      impedance accepts up to 3000\,N/m
    & lag $\propto 1/K$, so $a_\text{max} \propto K$ (Eq.~\ref{eq:al-amax}); less position error from joint friction on
      the real arm (error $\approx$ friction$/K$)
    & more force per mm of lag on contact; the damping grows with $\sqrt K$, so encoder noise and loop delay matter
      more; stability margin on the real robot to be checked \\
  \addlinespace\cmidrule(lr){1-4}\addlinespace
  Rotational stiffness $K_o$ \newline (\textbf{100\,N\,m/rad}, was 30)
    & spring between target and tool orientation (10--300\,N\,m/rad)
    & tilt $\propto 1/K_o$ while translating (Table~\ref{tab:al-diag}); smaller angle lag when turning
    & twists a stem caught in the fork more stiffly; same noise/stability caveat as $K$ \\
  \addlinespace\cmidrule(lr){1-4}\addlinespace
  Damping $D$ \newline (\textbf{$2\sqrt K\Lambda^{1/2}$})
    & damper between target and tool velocity
    & critical in every pose and direction; Franka's $2\sqrt K$ (critical only for 1\,kg) overshot by up to 7\,mm
    & uses the model of $\Lambda$ (only for the damping) \\
  \addlinespace\cmidrule(lr){1-4}\addlinespace
  Velocity feedforward $\dot x_d$ \newline (\textbf{on})
    & the damper acts on the velocity \emph{difference}, not on the tool velocity; while the clamp holds the target,
      $\dot x_d$ is the tool's velocity (the held target moves with the tool)
    & no lag at constant speed (without it: 82\,mm at 20\,cm/s, Table~\ref{tab:al-time}); no damping force on top of
      the clamp's bound
    & none found \\
  \addlinespace\cmidrule(lr){1-4}\addlinespace
  Inertia feedforward \newline $J^\top\Lambda\ddot x_d$ \newline (\textbf{off}; option B3)
    & the controller supplies the force needed to accelerate the arm from the model (operational-space
      control~\cite{al-khatib1987}) instead of waiting for the spring to stretch
    & lag drops to $\Delta\Lambda\, a / K$, set only by the \emph{model error} $\Delta\Lambda$ (20\,\% error: 5 times less
      lag); also removes the tilt and tip drift from the coupling
    & the motion now depends on the model of $\Lambda$ (libfranka on the robot); needs the motor-inertia question
      settled and the model checked (identification) \\
  \addlinespace\cmidrule(lr){1-4}\addlinespace
  Control period \newline (\textbf{2\,ms} sim, 1\,ms robot)
    & how often the law is evaluated
    & 1\,ms in sim matches the robot exactly
    & simulation twice as slow, so training twice as long \\
  \bottomrule
\end{tabular}}
\end{center}

\clearpage
\paragraph{Commands (what the policy outputs, and the limits on it).}
\begin{center}
{\small\renewcommand{\arraystretch}{1.2}
\begin{tabular}{>{\raggedright\arraybackslash}p{0.16\textwidth}>{\raggedright\arraybackslash}p{0.23\textwidth}>{\raggedright\arraybackslash}p{0.245\textwidth}>{\raggedright\arraybackslash}p{0.245\textwidth}}
  \toprule
  Building block (current) & What it is & Gains & Costs \\
  \midrule
  Speed cap \newline (\textbf{10\,cm/s, 45$^\circ$/s}; 20\,cm/s until 8 Oct.)
    & largest target step per policy step
    & the policy can approach and turn within one episode; 10\,cm/s is reached and left again within about 13\,cm,
      i.e.\ within the stem area
    & 20\,cm/s needed 42\,cm to reach and leave, more than the stem area, and lagged up to 3.35\,mm on such moves
      (Table~\ref{tab:al-final}) \\
  \addlinespace\cmidrule(lr){1-4}\addlinespace
  Acceleration limit \newline (\textbf{0.08\,m/s$^2$, 20$^\circ$/s$^2$}; 0.1\,m/s$^2$ until 8 Oct.)
    & largest change of the target step from one policy step to the next
    & \emph{the} knob for the lag (Eq.~\ref{eq:al-lag}); part of the policy-side code, so identical in sim and on the
      robot; the policy sees the applied step, so it can learn the limit
    & less agility: long time to speed and long braking distance (table in Section~\ref{sec:al-problem}) \\
  \addlinespace\cmidrule(lr){1-4}\addlinespace
  Lag criterion $e_\text{max}$ \newline (\textbf{2\,mm, 2$^\circ$})
    & how much lag we accept; sets the acceleration limit via Eq.~\eqref{eq:al-amax}
    & smaller: the motion depends less on the arm's dynamics (better transfer~\cite{al-aljalbout2024})
    & smaller $a_\text{max}$; a larger value (e.g.\ 5\,mm, option B2) gives 2.5 times more acceleration \\
  \addlinespace\cmidrule(lr){1-4}\addlinespace
  Low-pass filter on the commands \newline (\textbf{not used})
    & smooths the commands instead of limiting their change
    & smooth motion; common in sim-to-real
    & adds hidden state; Ref.~\cite{al-aljalbout2024} saw learning collapse with a limiter the policy could not see;
      no hard bound on the acceleration \\
  \addlinespace\cmidrule(lr){1-4}\addlinespace
  Target clamp \newline $\lVert x_d - x\rVert \le e_\text{clamp}$ \newline (\textbf{4\,mm, 3$^\circ$}, every control
    step)
    & the target is never further than $e_\text{clamp}$ from the measured tool; checked at every control step, not
      once per policy step (which would also cut steps in free motion)
    & hard bound on the spring force when the tool is blocked: 4\,mm $\times$ 1000\,N/m = 4\,N (about 3 times the
      measured push forces of 0.8--1.4\,N), 3$^\circ \times$ 100\,N\,m/rad = 5.2\,N\,m (43\,\% of the FR3's 12\,N\,m
      wrist torque limit; 10$^\circ$ would exceed it), even if the policy keeps pushing
    & must stay above the free-space lag (2\,mm), otherwise it acts while moving; while it acts, the target does not
      follow the commands (the policy sees this in the applied step); runs in the 1\,kHz loop on the robot \\
  \addlinespace\cmidrule(lr){1-4}\addlinespace
  Policy rate \newline (\textbf{31.25\,Hz})
    & how often the policy sends a new step
    & higher: finer commands, quicker reaction
    & more policy evaluations per second of simulation, so longer training; more sensitive to latency \\
  \addlinespace\cmidrule(lr){1-4}\addlinespace
  Action space \newline (\textbf{Cartesian steps on the previous target})
    & alternatives: joint velocities, Cartesian velocities, absolute poses, joint torques
    & joint velocities transferred best in Ref.~\cite{al-aljalbout2024} (stiff joint-level tracking)
    & joint velocities: no Cartesian compliance at contact, which matters for a fragile stem; torques: motion set by
      the dynamics alone, could not be transferred safely~\cite{al-aljalbout2024} \\
  \bottomrule
\end{tabular}}
\end{center}

\clearpage
\paragraph{Task and training.}
\begin{center}
{\small\renewcommand{\arraystretch}{1.2}
\begin{tabular}{>{\raggedright\arraybackslash}p{0.16\textwidth}>{\raggedright\arraybackslash}p{0.23\textwidth}>{\raggedright\arraybackslash}p{0.245\textwidth}>{\raggedright\arraybackslash}p{0.245\textwidth}}
  \toprule
  Building block (current) & What it is & Gains & Costs \\
  \midrule
  Episode length \newline (\textbf{10\,s})
    & how long the policy has per attempt
    & longer: a slow approach still leaves time to push the stem
    & training time grows about in proportion; the reward comes later after the actions that earned it, which makes
      learning harder \\
  \addlinespace\cmidrule(lr){1-4}\addlinespace
  Workspace bound \newline (\textbf{planned})
    & keep the tool away from poses near joint limits and the edge of the reach
    & there the lag grows and the real FR3 stops with a reflex
    & excludes some stem placements \\
  \addlinespace\cmidrule(lr){1-4}\addlinespace
  Reward terms \newline (\textbf{planned}: contact force, action rate, joint-limit margin)
    & penalties the policy learns to avoid, instead of hard limits
    & the policy learns gentle motion itself, where it matters
    & soft: no guarantee on the robot \\
  \addlinespace\cmidrule(lr){1-4}\addlinespace
  Domain randomization \newline (\textbf{planned}, after a first baseline run)
    & vary $K$, $K_o$, $\Lambda$ (motor inertia, link masses), latency and friction in training
    & policy robust to the remaining differences between sim and real~\cite{al-peng2018}
    & ranges must be realistic; too wide makes the policy cautious and slow to train \\
  \addlinespace\cmidrule(lr){1-4}\addlinespace
  Identification \newline (\textbf{planned}, before the first real run)
    & compare libfranka's model with the simulation; record the real arm's response to the same commands
    & tells how large the dynamics gap is; prerequisite for inertia feedforward and for loosening the criterion
    & robot time \\
  \addlinespace\cmidrule(lr){1-4}\addlinespace
  FR3 safety layer in sim \newline (\textbf{planned})
    & joint velocity, torque and torque-rate limits and the collision reflex as terminations/penalties
    & the policy does not learn motions the robot refuses
    & thresholds must match the robot's configuration \\
  \bottomrule
\end{tabular}}
\end{center}

\paragraph{Main options, in short.}
\begin{center}
{\renewcommand{\arraystretch}{1.2}
\begin{tabular}{>{\raggedright\arraybackslash}p{0.24\textwidth}>{\raggedright\arraybackslash}p{0.18\textwidth}>{\raggedright\arraybackslash}p{0.22\textwidth}>{\raggedright\arraybackslash}p{0.26\textwidth}}
  \toprule
  Option & $a_\text{max}$ at 2\,mm lag & Depends on the arm model & Main cost \\
  \midrule
  B1: $K$ 1000, $K_o$ 100, no inertia feedforward & 0.08\,m/s$^2$ (measured; 0.1 lags 2.15\,mm) & only through the
    damping & sluggish \\
  \addlinespace
  B2: B1 with a 5\,mm criterion & $\approx 0.25$\,m/s$^2$ & yes: 5\,mm of lag are set by the arm & against the transfer
    evidence \\
  \addlinespace
  B3: B1 + inertia feedforward & $\approx 0.5$\,m/s$^2$ or more (not run) & yes: model of $\Lambda$ & identification
    first \\
  \addlinespace
  Stiffer $K$ (3000\,N/m), no inertia feedforward & $\approx 0.3$\,m/s$^2$ & only through the damping & real-robot
    stability, contact forces \\
  \bottomrule
\end{tabular}}
\end{center}

% ---------------------------------------------------------------------------------------------------------------------
\subsection{Plan}
\label{sec:al-plan}
Decision (7 October 2026, revised 8 October): \textbf{B1 for the first training and transfer}, B3 (inertia
feedforward) once a policy trains successfully and the robot's $\Lambda$ has been checked. Values from
Table~\ref{tab:al-final}: $K = 1000$\,N/m, $K_o = 100$\,N\,m/rad, speed caps 10\,cm/s (3.2\,mm per step) and
45$^\circ$/s, acceleration limits 0.08\,m/s$^2$ (0.08\,mm per step) and 20$^\circ$/s$^2$ (0.02$^\circ$ per step),
target clamp 4\,mm and 3$^\circ$ at every control step. On 8 October the speed cap was halved and the acceleration
limit lowered from 0.1\,m/s$^2$, after moves that run up to the cap and reversals that really reverse showed lags of
up to 3.35\,mm and 2.15\,mm (Table~\ref{tab:al-final}).

\paragraph{Is this fast enough for the task?} With these limits (speed up, then brake, from rest to rest): the
farthest stem tip, 31\,cm from the start pose, is reached in 4.4\,s (1.25\,s to reach 10\,cm/s, 1.85\,s at
10\,cm/s, 1.25\,s braking); a 90$^\circ$ turn takes 4.3\,s (peak 42$^\circ$/s), and turning can happen while
approaching. That leaves about 5.5\,s of the 10\,s episode for pushing. Braking from 10\,cm/s takes 6\,cm and 1.25\,s.
We consider this sufficient for a first training. If the trained policy runs out of time, we lengthen the episode
first (cost: training time) and loosen the limits only after the identification.

% ---------------------------------------------------------------------------------------------------------------------
\subsection{Questions}
\label{sec:al-questions}
\begin{enumerate}
  \item Is the analysis right: is the lag of a Cartesian impedance controller on the FR3 set by its apparent mass
        (8--22\,kg at the tool in simulation, half of it from the motors' reflected inertia)?
  \item Does the FR3's internal joint torque controller reduce the apparent motor inertia (as on the DLR
        lightweight arms), and is that why libfranka's model leaves the motor inertia out? This decides how the
        simulation must model the motors, and which mass matrix the damping of our controller should use on the
        robot.
  \item Stiffness: is $K = 1000$\,N/m (and $K_o = 100$\,N\,m/rad) in a libfranka torque loop acceptable on our FR3
        next to a plant? Has the lab run similar gains, and with which velocity filtering?
  \item Inertia feedforward (option B3): how accurate is libfranka's model in your experience? Has the lab used
        operational-space control with acceleration feedforward on the FR3?
  \item Is an acceleration of about 0.08\,m/s$^2$ (1.25\,s to 10\,cm/s) a reasonable limit near a plant, or would
        you accept larger tracking errors for more agility?
  \item Which reflexes and thresholds (collision behaviour, joint velocity limits) does the lab configure, so that we
        model the same ones in simulation?
\end{enumerate}

{\raggedright
\begin{thebibliography}{9}
\bibitem{al-aljalbout2024} E.~Aljalbout, F.~Frank, M.~Karl, P.~van der Smagt. On the role of the action space in
  robot manipulation learning and sim-to-real transfer. \emph{IEEE Robotics and Automation Letters}, 2024.
  arXiv:2312.03673.
\bibitem{al-khatib1987} O.~Khatib. A unified approach for motion and force control of robot manipulators: the
  operational space formulation. \emph{IEEE Journal on Robotics and Automation}, 3(1), 1987.
\bibitem{al-peng2018} X.~B.~Peng, M.~Andrychowicz, W.~Zaremba, P.~Abbeel. Sim-to-real transfer of robotic control
  with dynamics randomization. \emph{IEEE ICRA}, 2018.
\end{thebibliography}
}
% ======================================== END CONTENT ========================================
